% Options for packages loaded elsewhere
\PassOptionsToPackage{unicode}{hyperref}
\PassOptionsToPackage{hyphens}{url}
\documentclass[
  11pt,
]{article}
\usepackage{xcolor}
\usepackage[margin=25mm]{geometry}
\usepackage{amsmath,amssymb}
\setcounter{secnumdepth}{-\maxdimen} % remove section numbering
\usepackage{iftex}
\ifPDFTeX
  \usepackage[T1]{fontenc}
  \usepackage[utf8]{inputenc}
  \usepackage{textcomp} % provide euro and other symbols
\else % if luatex or xetex
  \usepackage{unicode-math} % this also loads fontspec
  \defaultfontfeatures{Scale=MatchLowercase}
  \defaultfontfeatures[\rmfamily]{Ligatures=TeX,Scale=1}
\fi
\usepackage{lmodern}
\ifPDFTeX\else
  % xetex/luatex font selection
\fi
% Use upquote if available, for straight quotes in verbatim environments
\IfFileExists{upquote.sty}{\usepackage{upquote}}{}
\IfFileExists{microtype.sty}{% use microtype if available
  \usepackage[]{microtype}
  \UseMicrotypeSet[protrusion]{basicmath} % disable protrusion for tt fonts
}{}
\makeatletter
\@ifundefined{KOMAClassName}{% if non-KOMA class
  \IfFileExists{parskip.sty}{%
    \usepackage{parskip}
  }{% else
    \setlength{\parindent}{0pt}
    \setlength{\parskip}{6pt plus 2pt minus 1pt}}
}{% if KOMA class
  \KOMAoptions{parskip=half}}
\makeatother
\setlength{\emergencystretch}{3em} % prevent overfull lines
\providecommand{\tightlist}{%
  \setlength{\itemsep}{0pt}\setlength{\parskip}{0pt}}
\usepackage{bookmark}
\IfFileExists{xurl.sty}{\usepackage{xurl}}{} % add URL line breaks if available
\urlstyle{same}
\hypersetup{
  hidelinks,
  pdfcreator={LaTeX via pandoc}}

\author{}
\date{}

\begin{document}

\section{Proper actions, universal proper spaces and equivariant
K-homology}\label{proper-actions-universal-proper-spaces-and-equivariant-k-homology}

\emph{Written by GPT-6.1 Sol (OpenAI). Public domain (CC0).}

An equivariant Fredholm cycle needs a space on which the group acts
properly. There is no preferred such space for a general group. A
universal proper space collects them, and equivariant K-homology with
compact orbit support removes the choice of model. We construct the
universal comparison space, prove the comparison of its locally compact
stages, and give Euclidean, tree and enlarged local-field building
models.

Throughout, \(G\) is a second countable locally compact Hausdorff group
with a fixed left Haar measure \(dg\). Spaces on which \(C_0\) and
\(KK\) are used are locally compact, Hausdorff and second countable.
Equivariant cycles and their homotopies are those of
\href{KT-KK-17.html}{Equivariant KK-theory and the Green--Julg theorem},
Sections 1--2; the topology and Haar integration tools are proved in
\href{supporting/noncompact-foundations/noncompact-foundations.html}{Noncompact
foundations for equivariant induction}, Sections NCF.1 and NCF.3--NCF.5.
The root-group prerequisites for the building construction are stated
precisely at the beginning of its appendices. Basic freely accessible
references are {[}Lück{]}, {[}Gómez Aparicio--Julg--Valette{]} and
{[}Baum--Higson--Schick{]}. The universal proper-space formulation in
this setting is due to Baum, Connes and Higson.

\subsection{1. Properness and normalized
cutoffs}\label{properness-and-normalized-cutoffs}

\textbf{Definition 1.1.} A continuous action on \(X\) is \textbf{proper}
if \[
 G\times X\longrightarrow X\times X,\qquad(g,x)\longmapsto(gx,x)
 \tag{1.1}
\] is proper, meaning that the inverse image of every compact set is
compact. A proper space is \textbf{\(G\)-compact} if \(X/G\) is compact.
These two properties have different roles: properness controls
transporters and stabilizers; \(G\)-compactness controls orbit support.

For compact \(K,L\subset X\), write \[
 T(K,L)=\{g\in G:gK\cap L\ne\varnothing\}.
 \tag{1.2}
\] For locally compact Hausdorff \(X\), (1.1) is proper exactly when
every \(T(K,L)\) is compact. In one direction project the compact
inverse image of \(L\times K\) onto \(G\). Conversely the inverse image
of a compact subset of \(X^2\) is closed in \(T(K,L)\times K\), for its
two compact coordinate projections \(L,K\). The transporter itself is
closed: if group elements converge, their witnessing points in \(K\)
have a convergent subnet. In particular every stabilizer
\(G_x=T(\{x\},\{x\})\) is compact.

\textbf{Lemma 1.2 (orbit topology).} The quotient \(Y=X/G\) of a proper
space in the standing category is locally compact, Hausdorff and second
countable. Its quotient map \(\pi\) is open. If \(Y\) is compact, there
is a compact \(K\subset X\) with \(GK=X\).

\textbf{Proof.} Saturations of open sets are open. The orbit relation is
closed: for a convergent net \((g_i x_i,x_i)\to(y,x)\), put both
coordinates eventually in compact neighborhoods of their limits.
Properness puts the lifted pairs \((g_i,x_i)\) in a compact set; a
subnet gives \(g_i\to g\) and \(y=gx\). The map \(\pi\times\pi\) is an
open quotient map, so the closed relation descends to a closed diagonal
in \(Y^2\). Thus \(Y\) is Hausdorff. If \(V\subset X\) is precompact and
open, \(\pi(V)\) is open and its closure is contained in the compact,
hence closed, set \(\pi(\overline V)\). Such sets prove local
compactness. Images of a countable base form a countable base. Finally
finitely many \(\pi(V_j)\), with each \(V_j\) precompact, cover a
compact quotient. Set \(K=\bigcup_j\overline V_j\). \(\square\)

\textbf{Theorem 1.3 (squared normalization).} Every proper \(G\)-compact
space \(X\) in the standing category has \(c\in C_c(X)\), \(c\ge0\),
with \[
                 \int_G c(g^{-1}x)^2\,dg=1\qquad(x\in X).
 \tag{1.3}
\] More generally, every proper \(X\) has a continuous
\(b:X\to[0,\infty)\) such that \[
 \int_G b(g^{-1}x)\,dg=1,
 \qquad \operatorname{supp}b\cap\pi^{-1}(L)
       \text{ is compact for every compact }L\subset X/G.
 \tag{1.4}
\] One may take \(c=\sqrt b\); in the cocompact case \(c\) has compact
support.

\textbf{Proof.} First suppose \(X/G\) is compact. Take compact \(K\) as
in Lemma 1.2 and choose \(f\in C_c(X)\), \(f\ge0\), strictly positive on
\(K\). The compact cutoff construction of NCF.1 supplies such an \(f\).
Put \[
 a(x)=\int_G f(g^{-1}x)^2\,dg.
 \tag{1.5}
\] It is finite and continuous. Indeed, for \(x\) in a compact
neighborhood \(Q\), every contributing \(g\) is in
\(T(\operatorname{supp}f,Q)\). The integrand is jointly continuous
there, and the compact integration domain works throughout that
neighborhood. Every orbit meets the positive set of \(f\), so the
integrand is positive on a nonempty open subset of \(G\); Haar measure
gives \(a(x)>0\). Substitution \(g=sh\) proves \(a(sx)=a(x)\). Thus
\(c(x)=f(x)/\sqrt{a(x)}\) is continuous, supported in
\(\operatorname{supp}f\), and satisfies (1.3).

For arbitrary \(X\), choose nonnegative \(f_i\in C_c(X)\) whose positive
sets \(W_i\) meet every orbit. Let \(V_i=\pi(W_i)\) and \[
 h_i(x)=\int_G f_i(g^{-1}x)\,dg.
\] The same transporter argument makes \(h_i\) invariant and continuous;
its descended function \(\bar h_i\) is positive on \(V_i\). By Lemma 1.2
and the complete partition construction in NCF.1, choose a countable
locally finite partition \(\rho_i\) on \(Y\), with compact supports
contained in \(V_i\). Repeating a cover label after refinement is
allowed. Set \[
       b(x)=\sum_i\frac{\rho_i(\pi x)}{\bar h_i(\pi x)}f_i(x).
 \tag{1.6}
\] Extend each summand by zero outside \(\pi^{-1}(V_i)\). This is
continuous because \(\operatorname{supp}\rho_i\) is a compact subset of
the positive set of \(\bar h_i\). The sum is locally finite on the
quotient. Its orbit integral is \(\sum_i\rho_i(\pi x)=1\). Above a
compact quotient set, only finitely many partition supports occur. The
corresponding support of \(b\) is a closed subset of the union of
finitely many \(\operatorname{supp}f_i\), and is therefore compact. This
proves (1.4). \(\square\)

The support condition has a useful local consequence. For any compact
\(Q\subset X\), all \(g\) with \(b(g^{-1}x)\ne0\) for some \(x\in Q\)
lie in one compact set: use the compact support of \(b\) over
\(\pi(Q)\), then (1.2). For a discrete group, there are only finitely
many such \(g\). This local finiteness will make the Euclidean and tree
constructions continuous.

A slice at a point with stabilizer \(H\) means an \(H\)-invariant
subspace \(S\) for which the induced map \(G\times_H S\to X\) identifies
an invariant neighborhood with that associated space. Slices give useful
local descriptions when a slice theorem applies. The proofs of Theorem
1.3 and the models below use transporters and cutoffs directly; no
slice-existence assertion is needed in them.

\subsection{2. An explicit universal comparison
space}\label{an-explicit-universal-comparison-space}

An ambient model may fail to be locally compact. We will only form
\(C_0(Z)\) on its closed, locally compact, second countable,
\(G\)-compact stages \(Z\). Say that a proper Hausdorff ambient space
\(E\) has the \textbf{universal property for locally compact sources}
when every proper space in the standing category has an equivariant map
into \(E\), and any two such maps are equivariantly homotopic. We use
compactly generated ambient spaces: a subset is closed if its
intersection with each compact subset is closed in that compact subset.
Metric spaces have this property. The usual full universal proper model
has the same mapping and homotopy property on the full chosen category
of proper spaces, including ambient models; Section 5 explains the
resulting global uniqueness. The stage comparison proved here requires
only the explicitly stated locally compact-source property.

\textbf{Theorem 2.1 (the positive unit sphere).} With left translation,
the metric space \[
 \mathcal U=\{u\in L^2(G,dg):u\text{ is real and nonnegative a.e.},\ \|u\|_2=1\},
 \quad (\lambda_su)(g)=u(s^{-1}g),
 \tag{2.1}
\] is proper and has the universal property for locally compact sources.
For every compact subgroup \(H\), \(\mathcal U^H\) is nonempty and
contractible; no noncompact closed subgroup has a fixed point. The two
projections \(\mathcal U\times\mathcal U\to\mathcal U\) are
equivariantly homotopic.

\textbf{Proof.} The positive cone is norm closed, so \(\mathcal U\) is a
closed subset of the unit sphere. NCF.1 and NCF.3--NCF.5 imply that
\(C_c(G)\) is dense in \(L^2(G)\): truncate over a compact exhaustion,
approximate by simple functions, use Radon inner and outer
approximation, and separate the compact inner set from the complement of
its open outer set by a continuous cutoff. A countable relatively
compact base and rational simple approximations make \(L^2(G)\)
separable. Translation is norm continuous on \(C_c(G)\), by joint
continuity on a compact support neighborhood; density and translation
isometry give strong continuity on \(L^2(G)\). Thus the action is
continuous.

For \(u,v\in C_c(G)\), the coefficient \(\langle\lambda_g u,v\rangle\)
vanishes unless
\(g\in(\operatorname{supp}v)(\operatorname{supp}u)^{-1}\).
Cauchy--Schwarz and approximation give \[
                   \langle\lambda_gu,v\rangle\longrightarrow0
                   \quad(g\longrightarrow\infty)
 \tag{2.2}
\] for every \(u,v\in L^2(G)\). Finite norm nets and Cauchy--Schwarz
show that this convergence is uniform for \(u,v\) in fixed compact
subsets. For compact \(K,L\subset\mathcal U\), the transporter is closed
and lies in a compact subset of \(G\), since a matching pair has
coefficient one. The inverse image of a compact subset of
\(\mathcal U^2\) under the action map is therefore closed in a compact
transporter times a compact second projection. This proves properness,
including for this ambient space which need not be locally compact.

For compact \(H\), average a nonzero nonnegative \(f\in C_c(G)\) over
normalized Haar measure on \(H\): \[
 v=\int_H\lambda_hf\,dh.
\] This is nonnegative, \(H\)-fixed and compactly supported, and
\(\int_Gv=\int_Gf>0\). Normalize it to \(u_H\in\mathcal U^H\).
Normalized convex interpolation \[
              P(u,v,t)=\frac{(1-t)u+tv}{\|(1-t)u+tv\|_2}
 \tag{2.3}
\] is continuous and equivariant. Its denominator is at least
\(1/\sqrt2\), because \(\langle u,v\rangle\ge0\). Taking \(v=u_H\)
contracts \(\mathcal U^H\), and taking arbitrary \(u,v\) gives the
homotopy of projections. If a noncompact closed subgroup fixed \(u\),
its coefficient with itself would be one throughout that subgroup,
contrary to (2.2).

For a proper \(X\), use Theorem 1.3 and define \[
                    F_b(x)(g)=\sqrt{b(g^{-1}x)}.
 \tag{2.4}
\] Equation (1.4) makes its norm one. On a compact neighborhood in
\(X\), its integrands are supported in one compact subset of \(G\), as
observed after Theorem 1.3. Joint continuity and integration give norm
continuity. Substitution gives \(F_b(sx)=\lambda_sF_b(x)\). Any two
equivariant maps \(X\to\mathcal U\), whether or not constructed from
cutoffs, are homotopic by (2.3). \(\square\)

\textbf{Lemma 2.2 (locally compact stages).} For every compact
\(K\subset\mathcal U\), the saturation \(Z_K=GK\) is closed, locally
compact, Hausdorff, second countable, proper and \(G\)-compact.

\textbf{Proof.} Fix \(v\in\mathcal U\). By uniform (2.2), there is
compact \(C\subset G\) such that \(|\langle\lambda_gk,v\rangle|<1/2\)
for \(k\in K\), \(g\notin C\). If \(\|\lambda_gk-v\|<1/2\), then
\(\langle\lambda_gk,v\rangle=1-\|\lambda_gk-v\|^2/2>7/8\), so
\(g\in C\). Consequently the part of \(GK\) near \(v\) is contained in
the compact image \(CK\). Every closure point of \(GK\) lies in such a
compact image, hence lies in \(GK\); this proves closedness. At
\(v\in GK\), the closure of a sufficiently small relative ball is a
closed subset of \(CK\), proving local compactness. Separation and
second countability are inherited from the separable metric space
\(\mathcal U\). Properness restricts to a closed invariant subspace, and
the quotient is the image of \(K\). \(\square\)

\textbf{Lemma 2.3 (proper maps and their images).} If \(X\) is proper,
locally compact and \(G\)-compact, and \(T\) is a proper Hausdorff
ambient space with compact transporters, every continuous equivariant
map \(f:X\to T\) is proper. If \(T\) is compactly generated, its image
is closed, locally compact and second countable, with compact orbit
space.

\textbf{Proof.} Choose compact \(K\) with \(GK=X\). For compact
\(L\subset T\), any \(gk\) mapping into \(L\) has \(g\in T(f(K),L)\).
Thus \(f^{-1}(L)\) is closed in the compact image of
\(T(f(K),L)\times K\). This is properness.

For closed \(D\subset X\), the intersection
\(f(D)\cap L=f(D\cap f^{-1}(L))\) is compact and closed in \(L\).
Compact generation of \(T\) makes \(f(D)\) closed. Thus \(f\) is a
closed map and \(S=f(X)\) is closed in \(T\). To see that \(S\) is
locally compact, cover the compact fiber over \(s\in S\) by finitely
many relatively compact open sets of \(X\), with union \(V\). The set
\(S\setminus f(X\setminus V)\) is an open neighborhood of \(s\) in
\(S\), and its closure is contained in the compact set
\(f(\overline V)\). If \(\mathcal B\) is a countable base of \(X\), the
sets \(S\setminus f(X\setminus V)\), for finite unions \(V\) of members
of \(\mathcal B\), form a countable base: cover each compact fiber
inside the inverse image of a prescribed open set by finitely many such
members. Finally \(S/G\) is the image of \(f(K)\). \(\square\)

\subsection{3. Universality by squared-distance
minimization}\label{universality-by-squared-distance-minimization}

We will use a complete geodesic metric space \(Y\) with the CAT(0)
squared-distance inequality \[
 d(z,\gamma_{p,q}(t))^2
 \le(1-t)d(z,p)^2+t d(z,q)^2-t(1-t)d(p,q)^2.
 \tag{3.1}
\] The geodesic \(\gamma_{p,q}\) is parametrized on \([0,1]\). In a
CAT(0) space it is unique. Here the targets are also proper metric
spaces, meaning that closed bounded balls are compact. Geodesic
interpolation is then continuous directly: for converging endpoints and
parameters, the intervening points lie in a common compact ball; every
convergent subsequence has the two prescribed endpoint distances and
hence is the unique corresponding point on the limiting geodesic.

\textbf{Lemma 3.1 (barycenters).} A compactly supported probability
measure \(\mu\) on a complete space satisfying (3.1) has a unique
minimizer of \[
                           E_\mu(y)=\int_Y d(y,z)^2\,d\mu(z).
 \tag{3.2}
\] For a family with a common bounded support and densities continuous
in \(L^1\), its minimizer is continuous.

\textbf{Proof.} Choose \(o\in Y\) and \(R\) with the support inside
\(\overline B(o,R)\). For \(d(y,o)\ge R\), the triangle inequality gives
\(E_\mu(y)\ge(d(y,o)-R)^2\), while \(E_\mu(o)\le R^2\). Put
\(m=\inf E_\mu\), and choose \(y_n\) with \(E_\mu(y_n)\le m+1/n\).
Integrated midpoint convexity gives \[
 m\le E_\mu(\gamma_{y_n,y_j}(1/2))
 \le m+\frac1{2n}+\frac1{2j}-\frac14d(y_n,y_j)^2.
\] Thus the sequence is Cauchy. Its limit exists by completeness, and
continuity of the energy on bounded balls shows it attains the infimum.
The same inequality makes two minimizers equal. Notice that compactness
of a minimizing ball is not needed for this existence argument.

For continuity, suppose the densities \(w_x\) on a common finite measure
space converge in \(L^1\), and their support points all lie in
\(\overline B(o,R)\). Every minimizer is in \(\overline B(o,2R)\). On
this ball, \[
       |E_x(y)-E_{x_0}(y)|\le9R^2\|w_x-w_{x_0}\|_1=:\varepsilon_x.
 \tag{3.3}
\] If \(R=0\), the minimizer is always \(o\). Otherwise minimality gives
\(E_{x_0}(m_x)-E_{x_0}(m_{x_0})\le2\varepsilon_x\). Midpoint convexity
and minimality for \(E_{x_0}\) give
\(d(m_x,m_{x_0})^2\le2(E_{x_0}(m_x)-E_{x_0}(m_{x_0}))\le4\varepsilon_x\).
This tends to zero. \(\square\)

\textbf{Theorem 3.2 (a CAT(0) model).} Suppose \(G\) acts continuously,
properly and isometrically on a proper complete CAT(0) metric space
\(Y\). Then \(Y\) has the universal property for locally compact proper
sources. Every compact subgroup has a nonempty contractible fixed set,
and the two projections \(Y\times Y\to Y\) are equivariantly homotopic.

\textbf{Proof.} For a proper source \(X\), choose \(b\) from (1.4) and
fix \(y_0\in Y\). Put \[
 \mu_x=(g\mapsto gy_0)_*\bigl(b(g^{-1}x)\,dg\bigr),
 \qquad f(x)=\operatorname*{argmin}_{y\in Y}
                 \int_Gd(y,gy_0)^2b(g^{-1}x)\,dg.
 \tag{3.4}
\] These are probability measures. On a compact neighborhood of any
\(x_0\), the contributing \(g\)'s lie in one compact subset
\(C\subset G\). Its image \(Cy_0\) is compact and bounded; joint
continuity makes the densities continuous in \(L^1(C)\). Lemma 3.1
proves existence, uniqueness and continuity of \(f\). For a discrete
group the sum has finitely many terms throughout that neighborhood.
Substitution \(g=sh\), using left Haar measure, gives
\(\mu_{sx}=s_*\mu_x\). Isometry and uniqueness imply \(f(sx)=sf(x)\).

For two equivariant maps \(f_0,f_1\), the homotopy
\(\gamma_{f_0(x),f_1(x)}(t)\) is continuous by geodesic interpolation
and equivariant by uniqueness. Taking the two coordinate projections
gives the projection homotopy. If \(H\) is compact, the measure obtained
by averaging \(hy_0\) over normalized Haar measure on \(H\) has compact
support and is \(H\)-invariant. Its unique minimizer is \(H\)-fixed. The
fixed set is closed and convex; contraction along the geodesics from
that fixed minimizer is a contraction through fixed points. \(\square\)

\textbf{Example 3.3 (lattices).} Let \(\mathbb Z^n\) act on
\(\mathbb R^n\) by translations. Transporters of compact sets lie in a
bounded subset of the discrete lattice, hence are finite. The quotient
is compact because \([0,1]^n\) meets every orbit. Euclidean space is
proper and complete, and (3.1) is the squared-norm identity. Theorem 3.2
proves it is a universal proper model in the stated category. Its map
(3.4) is the explicit affine average \[
                   f(x)=\sum_{m\in\mathbb Z^n}b(m^{-1}x)m.
 \tag{3.5}
\] Here \(y_0=0\). The weights sum to one and are locally finite, so
\(f(sx)=s+f(x)\). In this example the geodesic homotopy is straight-line
interpolation.

\textbf{Example 3.4 (free groups).} Let \(F_n\), \(n\ge1\), be the group
of reduced words in \(a_1^{\pm1},\ldots,a_n^{\pm1}\). Its Cayley graph
joins \(w\) to \(wa_i\); give each edge length one. Cancellation gives a
unique reduced word between two vertices, so there is a unique simple
path and no circuit. The geometric realization \(T_n\) is therefore a
tree. Its degree is \(2n\), hence every bounded ball lies in a finite
subtree. This proves properness of the metric and, by compactness of a
ball containing a Cauchy sequence, completeness.

Here is a direct verification of (3.1). For a segment \([p,q]\) of
length \(L\), let \(w\) be the branch point where the path from \(z\)
meets that segment. Put \(a=d(p,w)\), \(h=d(z,w)\). Then \[
 d(z,\gamma_{p,q}(t))=h+|a-tL|,
 \quad d(z,p)=h+a,\quad d(z,q)=h+L-a.
\] The right side of (3.1) equals \[
 h^2+(a-tL)^2+2h\bigl((1-t)a+t(L-a)\bigr).
\] The last parenthesis is at least \(|a-tL|\): subtracting this
absolute value gives \(2t(L-a)\) if \(a\ge tL\), and \(2(1-t)a\)
otherwise. This proves the inequality for every point of the metric
tree.

Left multiplication acts by isometries, freely on vertices and without
edge inversions. An inversion of \([w,wa_i]\) would imply both
\(gw=wa_i\) and \(gwa_i=w\), hence \(a_i^2=1\), which is impossible for
reduced words. Thus every point stabilizer is trivial. Compact sets meet
finitely many edges; only finitely many translates of either of their
finite vertex sets can meet the other. If an edge-interior intersection
occurs, their endpoint sets also meet after the same translation. Hence
transporters are finite, proving the action proper. The finitely many
edges \([1,a_i]\), together with their endpoints, meet every orbit, so
it is \(F_n\)-compact. Theorem 3.2 proves that \(T_n\) is a universal
proper model. The case \(F_2\) is the four-valent tree.

For a torsion-free discrete group, compact subgroups are finite and
consequently trivial. Proper stabilizers are therefore trivial, so the
family of proper actions coincides with the family of free proper
actions. This is the reason \(\underline E\Gamma\) and the free
universal model \(E\Gamma\) coincide in that case. In particular \(F_n\)
is torsion free: conjugate a nontrivial reduced word to a cyclically
reduced word \(v\); the reduced length of \(v^m\) is \(m|v|\) for every
positive \(m\).

\textbf{Example 3.5 (compact groups).} If \(G\) is compact, its action
on a point is proper. Every space has exactly one map to that point, so
both existence and homotopy uniqueness are immediate. The point is
already \(G\)-compact.

\textbf{Theorem 3.6 (enlarged local-field buildings).} Let \(k\) be a
nonarchimedean local field and \(\mathbf G\) a connected reductive
\(k\)-group. There is an enlarged building \(\mathcal B(\mathbf G,k)\)
with a continuous, proper, cocompact action of \(\mathbf G(k)\). It is
locally compact and second countable, and its metric is proper, complete
and CAT(0). Its apartments retain the entire Euclidean factor belonging
to the \(k\)-split centre. Consequently it is a universal proper model
for locally compact sources.

\textbf{Proof.} The split integral coordinates and the nonreduced
unitary calculations are proved in Sections BF.0--BF.7 below. Sections
AB.0--AB.10 prove finite-extension valuations, the quotient topology,
the apartment and metric comparisons, and integral unramified descent.
Sections GB.0--GB.23 complete the relative-root construction and descent
for an arbitrary connected reductive group. More precisely, GB.16 proves
torus-cocycle vanishing, GB.17 the rational regular-class
representative, GB.18 quasi-splitting over a finite unramified
extension, GB.19 the anisotropic centralizer and its full central
valuation lattice, GB.20 the exact rational root bounds, GB.21 the
residue/coface correspondence, and GB.22 common maximal rational
apartments and wall exchange. GB.23 proves the full stabilizers, proper
quotient topology, equality with the inherited CAT(0) metric, and
compactness of metric balls and of the orbit space. Apply Theorem 3.2 to
this proper complete metric. \(\square\)

The enlargement is necessary even for \(\mathbf G=\mathbf G_m\). Its
reduced building is a point, on which \(k^\times\) would have a
noncompact stabilizer. Its enlarged building is \(\mathbb R\), with
translation \(x\mapsto x-\omega(t)\); the kernel \(O_k^\times\) is
compact and the valuation quotient is a lattice. Sections AB.10 and
GB.19 prove the corresponding assertion for every torus and for the
centralizer used in descent.

\subsection{4. Equivariant K-homology with compact orbit
support}\label{equivariant-k-homology-with-compact-orbit-support}

\textbf{Definition 4.1.} For a proper locally compact \(G\)-compact
space \(X\), put \[
                   K_i^G(X)=KK_i^G(C_0(X),\mathbb C),\qquad i\in\mathbb Z/2.
 \tag{4.1}
\] A proper equivariant map \(f:X\to Y\) induces \(f_*\) by pulling
functions back: \(f^*:C_0(Y)\to C_0(X)\). On a cycle it replaces
\(\phi\) by \(\phi\circ f^*\), retaining its Hilbert space, action and
operator. The cycle defects, direct sums and interval homotopies are
preserved because every tested function is pulled back to \(C_0(X)\).
Thus \(f_*\) is a group homomorphism, \((h\circ f)_*=h_*f_*\), and the
identity map induces the identity. This describes covariance in spaces
explicitly.

Proper equivariant homotopies induce equal maps. If
\(H:X\times[0,1]\to Y\) is proper, \(a(H(x,t))\) is continuous and
vanishes uniformly at infinity in \(x\): for a compact support of \(a\),
properness puts all its preimages in one compact subset of the cylinder.
Approximation by compactly supported functions gives the assertion for
every \(a\in C_0(Y)\). The representations on the interval Hilbert
module therefore give a Kasparov homotopy between the two pullback
cycles. Lemma 2.3 makes every equivariant homotopy out of a
\(G\)-compact proper cylinder proper into its image stage.

\textbf{Definition 4.2.} For a proper compactly generated ambient model
\(E\), let \(\mathscr S(E)\) be its closed invariant, locally compact,
second countable, \(G\)-compact subspaces. Set \[
                 RK_i^G(E)=\varinjlim_{Z\in\mathscr S(E)}K_i^G(Z).
 \tag{4.2}
\] The transition map is induced by the closed inclusion. The stages are
directed: the union of two stages is closed and \(G\)-compact; it is
locally compact by taking neighborhoods with compact closure in the two
closed summands, and second countable by the closed-map countable-base
argument of Lemma 2.3 applied to their finite disjoint union. An element
of (4.2) is represented on one stage, and two representatives agree
exactly when their images agree on a larger stage. The notation refers
to equivariant K-homology with \(G\)-compact supports; it is distinct
from Kasparov's bivariant \(RKK\) notation for families over a base.

For \(\mathcal U\), the subspaces \(Z_K\) of Lemma 2.2 are cofinal: if
\(Z\) is any stage, Lemma 1.2 gives compact \(K\subset Z\) with
\(GK=Z\). Hence \[
        RK_i^G(\mathcal U)=\varinjlim_{K\subset\mathcal U\ \mathrm{compact}}
                            KK_i^G(C_0(GK),\mathbb C).
 \tag{4.3}
\] Only locally compact spaces occur inside \(C_0\) in this formula.

\textbf{Theorem 4.3 (canonical model independence).} If \(E\) and \(F\)
are proper Hausdorff compactly generated ambient spaces with the
universal property for locally compact sources, their groups (4.2) are
canonically isomorphic. These isomorphisms compose canonically for three
models and respect equivariant cycles and proper maps of stages.

\textbf{Proof.} For each stage \(Z\subset E\), choose an equivariant map
\(u_Z:Z\to F\). By Lemma 2.3 it is proper, and its image \(W\) is a
stage of \(F\). Thus it induces \(K_i^G(Z)\to K_i^G(W)\to RK_i^G(F)\).
For two choices \(u_Z,v_Z\), universal homotopy uniqueness gives a
homotopy \(Z\times[0,1]\to F\); Lemma 2.3 makes its image a stage
containing the two endpoint images, and the preceding proper-homotopy
argument identifies the induced maps. The same argument for a stage
inclusion proves compatibility with transition maps. Therefore these
choices give a well-defined map \(\Phi_{E,F}\) of direct limits.

Choose the reverse maps stage by stage. On an original \(Z\subset E\),
their composite is an equivariant map back into \(E\). It is
equivariantly homotopic to the original inclusion
\(Z\hookrightarrow E\); the homotopy image is again a stage by Lemma
2.3. Thus \(\Phi_{F,E}\Phi_{E,F}\) acts as the identity on every
representative, hence on the limit. The other composite is identical by
the same reasoning. With three models, the chosen composite stage maps
are homotopic to the maps chosen directly; this proves
\(\Phi_{F,T}\Phi_{E,F}=\Phi_{E,T}\). All equalities use source pullbacks
and homotopies of actual locally compact stages, so they preserve the
full Kasparov cycle equivalence relations. \(\square\)

If a locally compact model \(E\) is itself \(G\)-compact, its whole
space is a final stage. The examples therefore give \[
 RK_i^{\mathbb Z^n}(E\mathbb Z^n)=KK_i^{\mathbb Z^n}(C_0(\mathbb R^n),\mathbb C),
 \quad RK_i^{F_n}(EF_n)=KK_i^{F_n}(C_0(T_n),\mathbb C),
 \tag{4.4}
\] and, for compact \(G\), \(RK_i^G(EG)=KK_i^G(\mathbb C,\mathbb C)\).
These are identifications of the source groups; the assembly
homomorphism and its further properties belong to the next lesson.

\textbf{Proposition 4.4 (finite groups).} If \(G\) is finite, then \[
                         RK_0^G(EG)=R(G),\qquad RK_1^G(EG)=0,
 \tag{4.5}
\] where \(R(G)\) is the Grothendieck ring of finite-dimensional complex
representations, with direct sum and tensor product.

\textbf{Proof.} Use the point model. The compact scalar-cycle comparison
and its homotopy relations are proved in
\href{KT-KK-17.html}{Equivariant KK-theory and the Green--Julg theorem},
Theorem 4.3 and Proposition 6.1. Here is the concrete index description.
Replace a scalar cycle by its unital summand, average its operator over
\(G\), and take its self-adjoint part. These are equivariant compact
perturbations of the represented operator and preserve its class, by
Proposition 1.2 of that lesson. An even cycle now has an invariant
Fredholm block \(T:H^+\to H^-\), with finite-dimensional invariant
kernel and cokernel. On their orthogonal complements its polar unitary
differs from \(T\) by a compact operator: the image of \(|T|\) in the
Calkin algebra is one. This complementary cycle is degenerate. The class
is therefore \[
                              [\ker T]-[\ker T^*].
 \tag{4.6}
\] Every finite representation occurs as a cycle on a finite-dimensional
even Hilbert space with zero operator.

For completeness, independence and injectivity under module homotopy
follow from the invariant compact algebra comparison in that preceding
proof, rather than from any assumption that kernel dimensions stay
constant. The regular representation contains every irreducible;
averaging inner products gives orthogonal decompositions and Schur's
lemma gives \[
        \mathcal K(\mathcal H_{\mathbb C}^G)^G
         \cong\bigoplus_{V\in\widehat G}\mathcal K(M_V)\otimes1_V.
 \tag{4.7}
\] There are finitely many \(V\), since the regular representation of
the finite group is finite dimensional before adding infinite trivial
multiplicity. The scalar-cycle comparison sends \(V\) to the rank-one
projection in its multiplicity block. Matrix stability and the
compact-operator K-theory calculation give one independent integer for
each \(V\) in degree zero and zero in degree one. Thus (4.6) identifies
the actual equivariant homotopy group with \(R(G)\). The product of the
finite zero-operator cycles for \(V,W\) is the zero-operator cycle on
\(V\otimes W\): its creation errors are compact and its positivity test
is zero. Hence this is a ring identification. Equivalently, an odd
invariant self-adjoint Fredholm operator is a compact perturbation of an
invariant involution by spectral calculus, including an arbitrary sign
on its finite kernel; that involution is degenerate. \(\square\)

\subsection{5. Other universal models and geometric
cycles}\label{other-universal-models-and-geometric-cycles}

The following assertions give context. They are stated with their exact
scope and are not used in the cutoff construction, Theorem 4.3, the
examples or the building proof.

\textbf{Stated, unused: recognition for numerable proper spaces.} A
space is numerable for the family of compact subgroups if it has an
invariant open cover admitting equivariant maps to homogeneous spaces
\(G/H\), \(H\) compact, and an invariant subordinate partition of unity.
Within that category, a space \(E\) is a full universal proper model if
and only if every compact \(H\) fixes a point and the two projections
\(E\times E\to E\) are equivariantly homotopic. This is {[}Lück,
Definition 2.1 and Theorem 2.5(ii){]}. Its numerability hypothesis is
part of the statement. In this lesson, existence of maps into
\(\mathcal U\), Euclidean spaces, trees and the local-field building has
been proved directly, rather than inferred from this recognition
assertion.

\textbf{Stated, unused: homogeneous and Rips models.} If \(G\) is almost
connected, meaning \(G/G^0\) is compact, it has a maximal compact
subgroup \(K\), unique up to conjugacy, and \(G/K\) is a numerable
universal proper model {[}Lück, Theorem 4.3{]}. This includes connected
real reductive Lie groups and their symmetric-space models. If a
finitely generated discrete group \(\Gamma\), with finite symmetric
generating set \(S\), satisfies \[
 d(x,y)+d(z,t)\le\max\{d(x,z)+d(y,t),d(x,t)+d(y,z)\}+2\delta
\] for its word metric, the barycentric subdivision of its Rips complex
\(P_d(\Gamma,S)\) is a universal proper \(\Gamma\)-CW model whenever the
integer \(d\ge16\delta+8\) {[}Lück, Theorem 4.16{]}. Its simplices are
finite sets of vertices of pairwise distance at most \(d\). Here
``finite model'' means finitely many cell orbits; the group itself can
be infinite. The existence of maximal compact subgroups and the Rips
contractibility argument are not proved here.

\textbf{Stated, unused: the geometric comparison.} For a finite CW-pair
\((X,A)\), a Baum--Douglas cycle consists of a smooth compact
spin\({}^c\) manifold \(M\) with boundary, a smooth Hermitian complex
vector bundle \(V\) on \(M\), and a continuous map \(f:M\to X\) with
\(f(\partial M)\subset A\). The relations are direct sum/disjoint union,
spin\({}^c\) bordism and modification by an even-dimensional
spin\({}^c\) vector bundle. Parity is the dimension of each component,
with disjoint union separating even and odd components. The Dirac
construction induces a natural isomorphism from this geometric group to
analytic K-homology for finite CW-pairs {[}Baum--Higson--Schick,
Definitions 5.1--5.7 and Theorem 6.2{]}. The free preprint proves this
precise nonequivariant theorem. It is not an assertion here about
arbitrary equivariant spaces. No part of the comparison theorem is
needed to define (4.1)--(4.3).

\subsection{6. Exercises and solutions}\label{exercises-and-solutions}

\textbf{Exercise 6.1.} Construct a squared-normalized cutoff for the
translation action of \(\mathbb Z\) on \(\mathbb R\).

\textbf{Solution.} Put \(h(x)=\max(1-|x|,0)\) and \(c(x)=\sqrt{h(x)}\).
It is continuous, nonnegative and supported in \([-1,1]\). Write
\(x=k+t\), \(k\in\mathbb Z\), \(0\le t<1\). If \(0<t<1\), the only
positive summands in \(\sum_m h(x-m)\) are \(1-t\) and \(t\); at \(t=0\)
the only positive summand is one. Thus \(\sum_m c(x-m)^2=1\). Counting
measure is the chosen Haar normalization. For \(\mathbb Z^n\), the
product \(\prod_jc(x_j)\) has the same property because the finite orbit
sum factors.

\textbf{Exercise 6.2.} Verify the two compact-fixed-point and
projection-homotopy properties for the Cayley tree of \(F_2\), and prove
the mapping property for proper locally compact sources.

\textbf{Solution.} The reduced-word argument in Example 3.4 proves
\(F_2\) torsion free, so its only compact subgroup is the trivial group.
The tree is nonempty and contracts along the unique paths to any chosen
vertex. For the diagonal action, \((p,q,t)\mapsto\gamma_{p,q}(t)\) is a
continuous equivariant homotopy from the first projection to the second.
Finally, for any proper source \(X\), (3.4) is the unique minimizer of a
finite sum of squared tree distances on every compact neighborhood of
\(X\). Lemma 3.1 proves its existence and continuity; the
change-of-variable argument proves equivariance. The same path homotopy
joins any two maps. This verifies universality directly, including the
mapping property.

\textbf{Exercise 6.3.} Prove that any two full universal proper models
are equivariantly homotopy equivalent. Explain the corresponding
statement for the locally compact-stage formulation.

\textbf{Solution.} Let \(E,F\) be models in the category on which their
full universal property is imposed. Existence gives equivariant maps
\(u:E\to F\), \(v:F\to E\). Homotopy uniqueness in \(E\) gives
\(vu\simeq_G\mathrm{id}_E\); uniqueness in \(F\) gives
\(uv\simeq_G\mathrm{id}_F\). Hence they are equivariantly homotopy
equivalent. In particular this argument applies to the locally compact
models whenever both are objects of the standing proper-space category.
If an ambient model is only specified by its universal property for
locally compact sources, global maps from that ambient space are not
part of that weaker definition. Theorem 4.3 instead proves the canonical
equivalence of all locally compact-stage K-homology limits, using
precisely the map and homotopy properties available. Neither conclusion
applies \(C_0\) to a non-locally-compact ambient space.

\textbf{Exercise 6.4.} For a finite \(G\), identify both groups
\(RK_i^G(EG)\) and the ring multiplication in degree zero.

\textbf{Solution.} The point is a final \(G\)-compact stage. Thus the
groups are \(KK_i^G(\mathbb C,\mathbb C)\). After invariant averaging,
an even cycle has the index (4.6). The compact comparison (4.7) proves
that irreducible representation classes form an independent integer
basis and that every class has this form. Tensoring the finite
zero-operator cycles gives the product, so \(RK_0^G(EG)=R(G)\) as rings.
The degree-one group of (4.7) is zero; alternatively, the invariant
involution replacement at the end of Proposition 4.4 makes every odd
cycle degenerate. Hence \(RK_1^G(EG)=0\).

\subsection{Algebraic and metric foundations for the building
model}\label{algebraic-and-metric-foundations-for-the-building-model}

The following sections supply the proof of Theorem 3.6. Their order is
by dependency: split root coordinates and compact point groups, affine
chamber geometry and valuation descent, then rational relative roots and
arbitrary-group descent. The complete central factor is retained at each
stage. The labels BF, AB and GB distinguish the three sets of
mathematical formulas and proof locators.

\subsection{Split root filtrations and compact chamber
groups}\label{split-root-filtrations-and-compact-chamber-groups}

\subsubsection{BF.0. Scope and exact algebraic
inputs}\label{bf.0.-scope-and-exact-algebraic-inputs}

Let k be any nonarchimedean local field, O its valuation ring, \(\pi\) a
uniformizer, and \(\omega\) ( \(\pi\) )=1. Put \textbar c\textbar=
\(q^{-\omega(c)}\) , where q is the finite residue-field cardinality.
Thus O and its fractional ideals are compact, O is complete, and k has
its valuation topology. These are part of the local-field hypothesis.
The compactness can also be seen by choosing residue representatives:
compatible finite truncations modulo \(\pi^n\) identify O with their
inverse limit; the valuation metric gives the same topology. A product
of finitely many O-balls is compact by a diagonal subsequence argument.

Let G be a connected split reductive k-group, T a split maximal torus,
X= \(X^*\) (T), \(X^\vee\) = \(X_{*}\) (T), \(\Phi\) its reduced root
system and \(a^\vee\) its coroots. The exact earlier programme inputs
are:

\begin{itemize}
\tightlist
\item
  \href{../AG-RG/AG-RG-03.html}{Roots and reductive groups of rank one},
  Lemma 3.1, Theorem 4.1 and Theorem 7.1: the integral big-cell open
  immersion, additive root parametrizations, the perfect opposite-root
  pairing, the rank-one \(SL_{2}\) homomorphism and its Gauss identity.
\item
  \href{../AG-RG/AG-RG-04.html}{Root data, Weyl chambers and the Bruhat
  decomposition}, Theorem 2.1, Sections 3--5 and Theorem 4.1: the
  reduced root datum, finite reflection Weyl group W= \(N_{G}\) (T)/T
  with root-reflection representatives, and scheme coordinates on each
  positive unipotent group in every root order; multiplication U\^{}-
  times T times U\^{}+ is an open immersion. Section 5 proves the
  torus-weight commutator formula over arbitrary rings.
\item
  \href{../AG-RG/AG-RG-05.html}{Pinnings and the classification of split
  reductive groups}, Theorem 5.1 and Theorem 10.1: the pinned group is
  obtained from a pinned split reductive group \(G_{Z}\) over Z. Theorem
  5.1 proves the pinning-preserving comparison; its proof also proves
  root transport and normalizer conjugation. The opposite-root pairs in
  these sections are paired by the rank-one perfect pairing.
\item
  \href{../AG-GS/AG-GS-01.html}{Group schemes, actions and Hopf
  algebras}, Theorem 2.1, Theorem 4.1 and Lemma 5.1: the affine Hopf
  dictionary, character weight decomposition over a ring, and finite
  subcomodules over a field.
\item
  The earlier commutative-algebra programme supplies the elementary ring
  inputs: \href{../AG-CA/AG-CA-03.html}{Noetherian and Artinian rings},
  Theorem 2.1, proves finite-type and localization Noetherianity;
  \href{../AG-CA/AG-CA-12.html}{Regular sequences, depth and
  Cohen--Macaulay modules}, Theorem 6.1, proves that regular local rings
  are domains; \href{../AG-CA/AG-CA-18.html}{Smooth algebras over a
  field and the Jacobian criterion}, Theorem 2.1, proves that smooth
  local rings are regular. The short DVR and finite-lattice arguments
  actually needed are also given in BF.4.
\end{itemize}

The root-valuation conventions are those of Bruhat--Tits I, Definition
6.2.1 and Examples 6.2.3(a)--(b). Integral root coordinates are also
discussed in Demazure, SGA 3, Exposé XXIII, Sections 1.1--1.2, and
Conrad, Reductive group schemes, Lemma 6.2.8. The preceding lessons
named above prove the algebraic inputs; the filtration, compactness and
chamber arguments are proved below.

Choose root coordinates \(x_{a}\) : \(\mathbf G_{a} \to U_{a}\) from
integral root frames of \(G_{Z}\) , with each negative frame paired with
its positive frame. Changing an integral frame only multiplies it by a
sign. Consequently a Weyl representative \(n_{w}\) formed from the
integral rank-one representatives carries \(x_{a}\) (u) to \(x_{wa}\) (
\(\epsilon\) u), \(\epsilon\) =±1: it is an isomorphism of integral root
lines, hence its scalar is a unit of Z. The earlier pinned comparison
transports these coordinates to G.

Write V= \(X^\vee\) tensor R. Its central part is \(V_{Z}\)
=intersection ker(a), and its derived part is the span
\(V_{\mathrm{der}}\) of the coroots. These form a direct sum. Indeed
average any inner product over the finite W. Each reflection \(s_{a}\)
fixes ker(a) and has minus-one line R \(a^\vee\) , so \(a^\vee\) is
perpendicular to ker(a). Therefore \(V_{\mathrm{der}}\) is perpendicular
to \(V_{Z}\) ; their dimensions are complementary since the roots span
the dual of the coroot space. W acts trivially on \(V_{Z}\) . A point of
the enlarged apartment is x=(v,z) in \(V_{\mathrm{der}}\) plus \(V_{Z}\)
.

\subsubsection{BF.1. Split valuations, all residue
characteristics}\label{bf.1.-split-valuations-all-residue-characteristics}

For v in \(V_{\mathrm{der}}\) define

\[
 \varphi_{v,a}(x_a(u))=\omega(u)+a(v),\qquad
 \varphi_{v,a}(1)=\infty,
\]

and for any real r,

\[
 U_{a,v,r}=x_a\bigl(\pi^{\lceil r-a(v)\rceil}O\bigr).
 \tag{BF.1}
\]

The values of nonidentity elements are Z+a(v). Additivity of \(x_{a}\)
and the nonarchimedean triangle inequality show that every upper level
set is a subgroup. Only the identity has infinite value. The displayed
fractional ideal is compact and open in k; \(x_{a}\) is an algebraic
coordinate isomorphism, so these groups are compact and open in
\(U_{a}\) and shrink to its identity as r tends to infinity.

Here and below ``open root subgroup'' means open in its root group, not
open in G.

For linearly independent roots a,b choose a positive system containing
both. In integral positive-root coordinates the commutator has the form

\[
 [x_a(u),x_b(t)]
 =\prod_{pa+qb\in\Phi,\ p,q>0}
 x_{pa+qb}(C_{pq}u^pt^q),\qquad C_{pq}\in\mathbb Z.
 \tag{BF.2}
\]

To recall why the exact integrality needed here follows from the earlier
coordinate proof, form the commutator morphism over Z. Its root
coordinates are polynomials in u,t. Under torus conjugation a monomial
\(u^p\) t\^{}q has character pa+qb. Setting either variable to zero
makes the commutator the identity, so p,q are positive. Independence of
a,b gives at most one exponent pair for each root. Coefficients of this
Z-morphism are integers. This proves BF.2 without assuming they are
units or dividing by any root-string constant. If u and t have
respective values at least r and s, then

\[
 \omega(C_{pq}u^pt^q)+(pa+qb)(v)\ge pr+qs.
\]

Thus the commutator is contained in the subgroup generated by
\(U_{pa+qb,v,pr+qs}\) . A coefficient which vanishes or has positive
valuation only improves this bound. This includes types B,C,F in residue
characteristic two and type \(G_{2}\) in residue characteristics two and
three.

The rank-one matrix calculation is

\[
 n_a(u)=x_{-a}(-u^{-1})x_a(u)x_{-a}(-u^{-1})
 =a^\vee(u)n_a(1),\quad u\ne0,
\]

because its \(SL_{2}\) matrix is

\[
 \begin{pmatrix}1&0\\-u^{-1}&1\end{pmatrix}
 \begin{pmatrix}1&u\\0&1\end{pmatrix}
 \begin{pmatrix}1&0\\-u^{-1}&1\end{pmatrix}
 =\begin{pmatrix}0&u\\-u^{-1}&0\end{pmatrix}.
 \tag{BF.3}
\]

Both opposite factors have value - \(\omega\) (u)-a(v)=-
\(\varphi_{v,a}\) ( \(x_{a}\) (u)). These are also the only two opposite
factors which can complete \(x_{a}\) (u) to a reflection representative:
in the \(SL_{2}\) product the two diagonal entries must vanish. The
rank-one central quotient does not change that condition or the root
coordinates, by \href{../AG-RG/AG-RG-03.html}{Roots and reductive groups
of rank one} Theorem 7.1.

Every normalizer element with Weyl image \(s_{a}\) is t \(n_{a}\) (1).
Conjugation carries \(x_{-a}\) (c) to \(x_{a}\) (-a(t)c); the change
between its two valuations is constant in c.~Finally

\[
 t x_a(c)t^{-1}=x_a(a(t)c),
\]

so torus conjugation adds \(\omega\) (a(t)) to the value. These prove
all the reduced valuation axioms of Bruhat--Tits I Definition 6.2.1: at
least three values, subgroup levels, constant opposite-reflection
change, the commutator bound, and the opposite-factor condition. The
doubled-root axiom is absent for a reduced root system. No
characteristic was excluded.

\subsubsection{BF.2. The apartment, origins and the sign of
translations}\label{bf.2.-the-apartment-origins-and-the-sign-of-translations}

The family \{ \(\varphi_{v}\) :v in \(V_{\mathrm{der}}\) \} is an affine
space: the difference \(\varphi_{v'}\) - \(\varphi_{v}\) on \(U_{a}\)
minus its identity is a(v'-v), and roots separate \(V_{\mathrm{der}}\) .
Let w be the Weyl image of n.~Define

\[
 (n\varphi)_a(u)=\varphi_{w^{-1}a}(n^{-1}un).
 \tag{BF.4}
\]

Translating \(v\) changes the value by \(a(wv)\). Thus
\(n(\varphi_v)=n(\varphi_0)+wv\). For \(t\in T(k)\), let
\(\widetilde\nu(t)\in V\) satisfy

\[
 \chi(\widetilde\nu(t))=\omega(\chi(t))\quad(\chi\in X),
\]

and let \(\nu\) (t) be its derived projection. BF.4, using \(t^{-1}\)
rather than t, gives

\[
 t:v\mapsto v-\nu(t),\qquad n_w:v\mapsto wv.
 \tag{BF.5}
\]

In particular BF.3 gives

\[
 n_a(u):v\mapsto v-(a(v)+\omega(u))a^\vee.
 \tag{BF.6}
\]

Its fixed wall is a(v)+ \(\omega\) (u)=0. This also directly checks the
sign of the affine reflection.

The affine family is independent of pinning. To prove this, prescribe
the ratios \(c_{\alpha}\) in \(k^*\) of new to old simple frames. The
simple roots are a basis of the root lattice Q, so they define a
homomorphism c:Q \(\to k^*\) . Over a finite field extension there is t
in T with \(\alpha\) (t)= \(c_{\alpha}\) for all simple \(\alpha\) : put
the inclusion Q \(\subset\) X into integer Smith normal form and adjoin
the finitely many roots required to solve the resulting monomial
equations. This permits inseparable roots; no separability is asserted
or needed. Torus conjugation and the uniqueness of the pinned comparison
imply that the ratio for any root a is c(a), up to a harmless integral
sign arising from the integral frame choices. Hence \(\omega\)
(c(a))=a(d) for one d in \(V_{\mathrm{der}}\) . The new coordinates give
\(\varphi'_v\)= \(\varphi_{v-d}\) as functions on the same root groups.
Their entire affine families are equal; only their origins differ.

\subsubsection{BF.3. The enlarged normalizer apartment is universal
proper}\label{bf.3.-the-enlarged-normalizer-apartment-is-universal-proper}

On V= \(V_{\mathrm{der}}\) plus \(V_{Z}\) use the action

\[
 t n_w:x\mapsto wx-\widetilde\nu(t).
 \tag{BF.7}
\]

This is a group action. Products of the integral Weyl representatives
have the same finite Weyl effect, and their residual torus factor is
integral with integral inverse, hence has zero valuation. Torus
conjugation transforms \(\widetilde\nu\) equivariantly. Thus the action
identifies

\[
 N/T(O)=X^\vee\rtimes W.
 \tag{BF.8}
\]

For detail, a basis of X identifies T(k) with ( \(k^*\) )\^{}rank X.
Coordinate valuations are arbitrary integers, their kernel is ( \(O^*\)
)\^{}rank X=T(O), and each cocharacter \(\lambda\) is realized by
\(\lambda\) ( \(\pi\) ). Every normalizer element is t \(n_{w}\) , by
the earlier Weyl comparison. A normalizer element acts trivially in BF.7
exactly when w=1 and its torus valuations vanish. This proves BF.8
including its kernel assertion. T(O) is compact and open in N; the
quotient is discrete.

Choose a W-invariant Euclidean norm on V. For compact sets K,L
\(\subset\) V, any affine element \(\lambda\) w with ( \(\lambda\) w)K
intersecting L has \(\lambda\) in the bounded set L-WK. A lattice has
finitely many points in a bounded set: in a lattice basis each integer
coordinate is bounded. The transporter in N is therefore a closed subset
of a finite union of compact T(O)-cosets and is compact. This is
properness of the N-action. Its orbit space is compact because the full
translation lattice \(X^\vee\) has a compact fundamental parallelepiped.
A point stabilizer projects to a finite subgroup of the discrete affine
group and is a finite union of T(O)-cosets; it is compact and open.

A compact subgroup H \(\subset\) N has finite image in the discrete
quotient. Average any point over that finite affine group. The result is
H-fixed. The entire fixed set is a nonempty affine subspace and is
contractible by straight interpolation. The two projections V times V
\(\to\) V are N-homotopic by the same interpolation.

Here is a direct universal-property proof. A numerable proper N-space
has an invariant locally finite partition of unity ( \(\lambda_{i}\) )
with support( \(\lambda_{i}\) ) contained in invariant open sets
\(Y_{i}\) , each with an equivariant map \(p_{i}\) : \(Y_{i} \to\) N/
\(H_{i}\) for a compact \(H_{i}\) . Choose \(z_{i} \in V^{H_{i}}\) . The
map \(nH_{i} \to nz_{i}\) is well defined and continuous. Then

\[
 f(y)=\sum_i\lambda_i(y)\,p_i(y)z_i
\]

is a continuous N-map: a point outside \(Y_{i}\) has a neighborhood
disjoint from support( \(\lambda_{i}\) ), so its term extends
continuously by zero, and the sum is locally finite, and affine maps
preserve a convex combination whose weights sum to one. Two such maps
are N-homotopic by straight interpolation, as are any two N-maps to V.
This proves universality in the numerable proper category. For proper
N-CW spaces the same conclusion follows cell by cell: each
compact-isotropy cell N/H times \(D^m\) maps into the nonempty
contractible affine space \(V^H\) ; radial extension in that affine
space extends boundary maps and homotopies. Thus V is the usual enlarged
model for the universal proper normalizer space, with its full central
translation directions retained.

\subsubsection{BF.4. A local integral representation
lemma}\label{bf.4.-a-local-integral-representation-lemma}

The following representation lemma supplies the linear algebra for
compactness of the point groups.

\textbf{Lemma.} A split reductive group scheme H over the complete
local-field DVR O admits a closed immersion into GL(L) for a finite free
O-module L. Every specified split torus in H acts on L through an
O-direct sum of character-weight modules.

\textbf{Proof.} Let A=O{[}H{]}. It is a finitely generated flat
O-algebra. Its generic and special fibres are geometrically integral.
Over an algebraically closed field, smoothness gives regular local
rings, and the exact earlier ring proofs in BF.0 make these domains.
Distinct irreducible components therefore cannot intersect: at an
intersection the local ring would have two minimal primes. There are
finitely many components, so each is open and closed. Connectedness
leaves one component, and the local domain assertion also gives
reducedness. Thus a smooth connected affine algebraic group is integral.
Reductivity supplies precisely smoothness and geometric connectedness
here. Flatness embeds A into its integral generic fibre, so A is a
domain, and A/ \(\pi\) A is a domain. Consequently the localization
\(A_{(\pi)}\) is a Noetherian local domain with maximal ideal generated
by \(\pi\) . It is a DVR. Indeed its intersection of powers of \(\pi\)
is zero: their intersection J is an ideal, hence finitely generated, and
J= \(\pi\) J by division. For generators of J write this equality as v=
\(\pi\) Mv. The determinant of I- \(\pi\) M is a unit in the local ring;
its adjugate then forces every generator to be zero. Thus J=0. Repeated
division then assigns every nonzero element a finite \(\pi\) -order, and
a remainder outside the maximal ideal is a unit. This defines a
valuation extending \(\omega\) on k.

Choose finitely many algebra generators of A. By the earlier
finite-comodule lemma \href{../AG-GS/AG-GS-01.html}{Group schemes,
actions and Hopf algebras}, Lemma 5.1 they lie in a finite-dimensional
k-subcomodule E \(\subset\) A{[} \(\pi^{-1}\) {]}; include 1. Set L=E
\(\cap\) A. It is an O-lattice in E. Here is the needed boundedness
proof: the DVR valuation just defined gives a genuine norm on E. In a
k-basis, the triangle inequality gives an upper bound by a constant
times the coordinate maximum norm. It also gives continuity. The
coordinate unit sphere is compact, by compactness of O, so that positive
continuous norm has a positive minimum there. Thus its unit ball is
contained in a bounded coordinate lattice. Since A lies in that
valuation's unit ball, L is bounded; clearing the finitely many basis
denominators makes it full. It is therefore finite free over the DVR.

L is saturated in A. Both A and A/L are torsion-free, hence flat over O:
each is a filtered union of its finitely generated torsion-free
submodules, which are free by DVR row elimination; tensoring an
injection remains injective first on each free submodule and then in the
filtered union. For l \(\in\) L, its coaction belongs to A tensor A, and
after inverting \(\pi\) its first factor belongs to E=L tensor k. Its
image in (A/L) tensor A is therefore \(\pi\) -torsion and zero, since
that tensor product is flat. Exactness after tensoring proves Delta(L)
\(\subset\) L tensor A. Thus L is an integral H-comodule.

In an O-basis write Delta( \(l_{j}\) )= \(\sum_{i} l_{i}\) tensor
\(a_{ij}\) . Counit and antipode give an invertible coefficient matrix
and a group-scheme map H \(\to\) GL(L). Applying counit in the first
factor writes every \(l_{j}\) as an O-linear combination of the
\(a_{ij}\) . All chosen algebra generators of A belong to L, so the
matrix-coefficient algebra surjects onto A. This is a closed immersion.
Finally the split-torus weight projections are obtained by comparing
Laurent monomials in the coaction, precisely as in
\href{../AG-GS/AG-GS-01.html}{Group schemes, actions and Hopf algebras},
Theorem 4.1. They split L into finitely many projective, hence free,
weight modules. \(\square\)

\subsubsection{BF.5. Compact point groups and their normalizer
intersections}\label{bf.5.-compact-point-groups-and-their-normalizer-intersections}

For x=(v,z) \(\in\) V let \(P_{v}\) be generated by T(O) and
\(U_{a,v,0}\) for every root. Let \(N_{x}\) be the stabilizer for BF.7
and put \(K_{x}\) = \(P_{v} N_{x}\) . Then \(P_{v}\) and \(K_{x}\) are
compact open subgroups of G(k), with

\[
 nP_vn^{-1}=P_{(nx)_{der}},\quad nK_xn^{-1}=K_{nx},
 \quad P_v\cap N\subset N_x,\quad K_x\cap N=N_x.
 \tag{BF.9}
\]

\textbf{Proof.} Apply BF.4 to the integral pinned model over O. Choose
an O-basis of each character-weight module \(L_{\lambda}\) . On L tensor
k define

\[
 \|\xi\|_x=\max_\lambda q^{-\lambda(x)}\,|\xi_\lambda|_{max}.
 \tag{BF.10}
\]

The expansion rho( \(x_{a}\) (u))= \(\sum_{i\ge 0} u^i R_{i}\) has
integral matrices and \(R_{0}\) =1. Torus equivariance says that
\(R_{i}\) maps the \(\lambda\) -weight module to the ( \(\lambda\)
+ia)-weight module. This follows by comparing Laurent characters over O,
rather than by differentiating; it remains true in positive
characteristic. If \(\omega\) (u)+a(v) \(\ge\) 0, each term has norm
bound

\[
 q^{-i\omega(u)}q^{-ia(x)}\le1,
\]

because a(z)=0. The inverse is \(x_{a}\) (-u), with the same bound. Thus
every generating root element is an isometry of BF.10. T(O) is an
isometry as well. The isometry group of this finite-dimensional norm is
compact: entries of an operator and its inverse are bounded by fixed
constants in the chosen basis; the pairs (M, \(M^{-1}\) ) form a closed
subset of a finite product of compact matrix balls. The closed immersion
in BF.4 is a topological embedding on k-points, since its inverse
coordinate functions are regular on its image. Its intersection with
this compact isometry group is compact.

\(P_{v}\) contains an open neighborhood of the identity. The big-cell
coordinates identify a sufficiently small product of root neighborhoods
and a torus neighborhood contained in T(O) with such a neighborhood; all
root levels here contain sufficiently small fractional ideals. Therefore
\(P_{v}\) is an open subgroup. An open subgroup is closed, since its
complement is a union of open cosets. \(P_{v}\) is consequently a closed
subset of the compact isometry intersection and is compact.

Write n=t \(n_{w}\) . The integral map rho( \(n_{w}\) ) identifies the
weight lattices by integral invertible matrices, so it preserves their
coordinate maximum norms. On a vector of weight \(\lambda\) , rho(n)
sends that lattice isometrically to its w- \(\lambda\) lattice and
multiplies it by (w \(\lambda\) )(t). The resulting BF.10 norm ratio is

\[
 q^{-\omega((w\lambda)(t))-(w\lambda)(x)+\lambda(x)}.
\]

If n \(\in P_{v}\) , it is an isometry, so this ratio is 1 for every
occurring weight \(\lambda\) . Equivalently \(\lambda\) (x)= \(\lambda\)
( \(w^{-1}\) (x+ \(\widetilde\nu\) (t))). The occurring weights span X
tensor R: otherwise a positive-dimensional subtorus would act trivially
in the closed faithful representation. This forces wx- \(\widetilde\nu\)
(t)=x, hence n \(\in N_{x}\) . Notice that this proves the full central
condition, not merely the derived condition.

BF.1--BF.2 show that normalizer conjugation transports every level group
and T(O) to the corresponding level at nx, proving the first equality of
BF.9. In particular \(N_{x}\) normalizes \(P_{v}\) . Thus \(K_{x}\) is a
subgroup. The compactness of \(N_{x}\) was proved in BF.3; explicitly,
in the Weyl component t \(n_{w}\) its fixed-point equation is
\(\widetilde\nu\) (t)=wx-x. Each such fibre is empty or a coset of the
compact kernel T(O). There are finitely many w, so \(N_{x}\) is compact.
Consequently \(K_{x}\) is compact as the product \(P_{v} N_{x}\) and is
open because it contains \(P_{v}\) . Normalizer conjugation gives its
covariance. If p \(n_{0} \in K_{x} \cap\) N with \(n_{0} \in N_{x}\) ,
then p \(\in P_{v} \cap\) N \(\subset N_{x}\) ; hence p
\(n_{0} \in N_{x}\) . The reverse inclusion is immediate. Pinning
independence follows because the generating root level sets are
functions of the valuation point, as proved in BF.2. \(\square\)

\subsubsection{BF.6. Split chamber factorization in every root
order}\label{bf.6.-split-chamber-factorization-in-every-root-order}

Remove the affine hyperplanes a(v)+m=0, m \(\in\) Z, from
\(V_{\mathrm{der}}\) . Let C be one of the resulting open chambers. The
integer

\[
 f_C(a)=\lceil-a(v)\rceil
\]

is constant for v \(\in\) C: it can change only on a removed hyperplane.
For every root, \(f_{C}\) (a)+a(v)\textgreater0 and \(f_{C}\) (a)+
\(f_{C}\) (-a)=1. Define \(I_{C}\) to be generated by T(O) and \(x_{a}\)
( \(\pi^{f_{C}(a)}\) O). For any positive system and any orders of its
positive and negative roots, multiplication is a homeomorphism

\[
 \left(\prod_{a\in-\Phi^+}\pi^{f_C(a)}O\right)
 \times T(O)\times
 \left(\prod_{a\in\Phi^+}\pi^{f_C(a)}O\right)
 \longrightarrow I_C.
 \tag{BF.11}
\]

In particular \(I_{C}\) is contained in the big cell; its root-block
factors are exactly its intersections with U\^{}- and U\^{}+.

\textbf{Proof.} We first prove that every finite word in the prescribed
root generators lies in the big cell with precisely the required root
bounds.

Fix roots \(a_{1}\) ,\ldots, \(a_{m}\) , with repetitions permitted, and
form over Z the word morphism

\[
 F:\mathbb A^m\to G_Z,\qquad
 (t_1,\ldots,t_m)\mapsto x_{a_1}(t_1)\cdots x_{a_m}(t_m).
\]

Let the split integral torus act on \(t_{i}\) with character \(a_{i}\) .
F is equivariant for conjugation. Let Omega=U\^{}- T U\^{}+ be the
integral big cell. Its inverse root coordinates have their corresponding
root characters, while torus characters on its T-factor are
conjugation-invariant. Let J be the radical ideal of the closed
complement of \(F^{-1}\) (Omega) in R=Z{[} \(t_{1}\) ,\ldots, \(t_{m}\)
{]}. This ideal is homogeneous. For a completely algebraic check, the
action on the product with the split torus is the Laurent-ring
automorphism \(t_{i} \mapsto t_{i} X^{a_{i}}\) of R{[} \(X_{1}^{±1}\)
,\ldots, \(X_{d}^{±1}\) {]}. Equivariance of F and conjugation
invariance of Omega make that closed complement invariant. Its extended
ideal is radical, because the Laurent polynomial ring over R/J is
reduced, so invariance of the closed set preserves that ideal. For f
\(\in\) J its image is \(\sum_{\chi} f_{\chi} X^\chi\) . Unique Laurent
coefficients imply \(f_{\chi} \in\) J for every character \(\chi\) .
This proves homogeneity without division or averaging. The zero section
lies in \(F^{-1}\) (Omega). Pulling J back along that section gives the
unit ideal of Z, so J contains a polynomial D with D(0)=1. Taking its
weight-zero part retains membership in J and the value 1. Thus choose D
of weight zero. The principal open D(D) lies in \(F^{-1}\) (Omega).

For the a-root coordinate, pullback along F on that principal open has
form \(N_{a}\) / \(D^h\) , with \(N_{a}\) an integral polynomial of
character a and \(N_{a}\) (0)=0. To justify homogeneity, use the
character grading of Z{[} \(t_{i}\) , \(D^{-1}\) {]}, where D has degree
zero; the homogeneous component of any representative of the specified
character is another representative. A torus character \(\chi\) and its
inverse pull back similarly with weight zero and value1 at zero.

Now substitute \(u_{i}\) with \(\omega\) ( \(u_{i}\) )+ \(a_{i}\)
(v)\textgreater0. Each nonconstant weight-zero monomial has valuation

\[
 \sum_i j_i\omega(u_i)
 =\sum_i j_i(\omega(u_i)+a_i(v))>0.
\]

Consequently D(u) \(\in\) 1+ \(\pi\) O and is a unit. F(u) therefore
lies in Omega(k). Each monomial of \(N_{a}\) has weight
\(\sum_{i} j_{i} a_{i}\) =a, and hence

\[
 \omega(\text{monomial})+a(v)
 \ge\sum_i j_i(\omega(u_i)+a_i(v))>0.
\]

Integral coefficients have nonnegative valuation; the denominator is a
unit. The nonarchimedean inequality gives the same strict bound for the
root coordinate. Its valuation is an integer, so that bound is exactly
\(\omega\) (root coordinate) \(\ge\) ceil(-a(v))= \(f_{C}\) (a). For a
torus character and its inverse, the constant term is1 and every other
term has positive valuation. Thus all torus coordinates are units, so
its torus factor belongs to T(O).

This argument uses a finite polynomial on a finite word; no convergence,
perfect-residue assumption, general Iwahori theorem, or
ramified-extension descent is hidden in it. A torus-unit factor in a
word can be moved to its left using torus conjugation: this only
multiplies root parameters by units. Inverses of root generators are
\(x_{a}\) (-u), which have the same bound. Therefore the entire
generated group \(I_{C}\) has exactly the ordered factorization asserted
in BF.11. Conversely every such factor is one of its generators.
Uniqueness follows from the earlier integral big-cell open immersion and
the root-coordinate isomorphisms in every order. Both maps are
continuous, since the inverse is given by regular big-cell coordinates.
Thus BF.11 is a homeomorphism.

If \(v_{0} \in\) closure C, then \(f_{C}\) (a)+a( \(v_{0}\) ) \(\ge\) 0
by taking limits. Hence every \(I_{C}\) generator lies in \(P_{v_{0}}\)
, so \(I_{C} \subset P_{v_{0}}\) . By BF.5 \(P_{v_{0}}\) is compact.
\(I_{C}\) is open: its coordinate product contains a neighborhood of the
identity by the big-cell chart. The quotient \(P_{v_{0}}\) / \(I_{C}\)
is therefore both compact and discrete and is finite. This proves the
claimed finite index, including points on several walls. \(\square\)

\subsubsection{BF.7. The nonreduced unitary rank-one
valuation}\label{bf.7.-the-nonreduced-unitary-rank-one-valuation}

Let L/k be any separable quadratic extension, with involution u
\(\mapsto\) bar(u). Extend \(\omega\) to L with its restriction to k
normalized; its value group is \(e^{-1}\) Z, where e is the ramification
index. The following proof covers ramified and unramified extensions,
residue characteristic two and field characteristic two.

First, \(\omega\) (bar(u))= \(\omega\) (u). To supply the field fact
used here, take a k-basis of L. Its coordinate maximum norm bounds
\textbar u\textbar{} from above by the ultrametric inequality. Thus
\textbar u\textbar{} is continuous in those coordinates; on the compact
coordinate unit sphere it is positive and has a positive minimum.
Scaling by powers of \(\pi\) proves the reverse bound. Every k-linear
map on L is consequently bounded. In particular
\textbar bar(u)\textbar{} \(\le\) B\textbar u\textbar{} for some B.
Apply this to \(u^n\) and take nth roots, then let n tend to infinity;
this gives \textbar bar(u)\textbar{} \(\le\) \textbar u\textbar.
Applying the same argument to bar(u) gives equality. The norm comparison
also proves that all closed valuation balls in L are compact and have
their usual coordinate topology.

Choose \(\lambda\) with \(\lambda \ne\) bar( \(\lambda\) ), put
\(\delta\) = \(\lambda\) -bar( \(\lambda\) ), and theta= \(\lambda\) /
\(\delta\) . Then

\[
 \bar\delta=-\delta,\qquad \theta+\bar\theta=1,
 \qquad L^0=\ker\operatorname{Tr}_{L/k}=\delta k.
 \tag{BF.12}
\]

Indeed the trace is onto since Tr(theta)=1; its kernel is
one-dimensional and contains \(\delta\) . These statements remain valid
in characteristic two, where \(\delta\) is fixed by the involution and
\(L^0\) =k. The pairs satisfying v+bar(v)=-u bar(u) are exactly

\[
 H=\{(u,-\theta u\bar u+\delta c):u\in L,\ c\in k\}.
 \tag{BF.13}
\]

This is also a coordinate homeomorphism L times k \(\to\) H. It uses
division by the nonzero \(\delta\) , not division by two.

Put

\[
 J=\begin{pmatrix}0&0&1\\0&1&0\\1&0&0\end{pmatrix},
 \quad G=\{g\in SL_3(L):\bar g^tJg=J\},
 \quad S=\{\operatorname{diag}(s,1,s^{-1}):s\in k^*\}.
\]

Write a(s)=s. For (u,v) \(\in\) H use exactly the coordinates

\[
 x_+(u,v)=\begin{pmatrix}1&-\bar u&v\\0&1&u\\0&0&1\end{pmatrix},
 \qquad
 x_-(u,v)=\begin{pmatrix}1&0&0\\u&1&0\\v&-\bar u&1\end{pmatrix}.
 \tag{BF.14}
\]

Matrix multiplication verifies their unitary equations. Conversely, the
unitary equation on a triangular unipotent matrix forces precisely these
coordinates and the trace relation. Thus \(U_{a}\) = \(x_{+}\) (H),
\(U_{-}\) a= \(x_{-}\) (H), \(U_{2a}\) = \(x_{+}\) (0, \(L^0\) ), and
\(U_{-}\) 2a= \(x_{-}\) (0, \(L^0\) ). Over L, the algebra L
\(\otimes_{k}\) L is L times L with the involution exchanging the
factors; the unitary equation identifies the second matrix with J times
the inverse transpose of the first times J. The determinant condition
therefore identifies the base-changed group with \(SL_{3}\) . In
particular this is the smooth connected reductive unitary group.

For clarity its k-split rank is one: characters of any k-split torus
acting faithfully on \(L^3\) pair under the invariant nondegenerate
Hermitian form with their negatives. In dimension three there is at most
one nonzero pair, so these characters span a space of dimension at most
one. S attains that rank. Its weights are a,0,-a; the upper u tangent
coordinates have weight a and the trace-zero v coordinates weight 2a,
both nonzero. This assertion compares torus characters, not their
derivatives in k, so characteristic two does not identify a with 2a. The
negative groups give their negatives, and differences of the three
representation weights give no other nonzero root. Hence the relative
root system is \{a,2a,-a,-2a\}.

For either sign, direct multiplication gives

\[
 \begin{aligned}
 x_\epsilon(u,v)x_\epsilon(w,z)
   &=x_\epsilon(u+w,v+z-\bar u w),\\
 x_\epsilon(u,v)^{-1}&=x_\epsilon(-u,\bar v),\\
 [x_\epsilon(u,v),x_\epsilon(w,z)]
   &=x_\epsilon(0,\bar w u-\bar u w),
 \end{aligned}
 \tag{BF.15}
\]

where the commutator is \(ghg^{-1} h^{-1}\) . Expanding the trace
verifies closure. The doubled-root group is central; it is the entire
centre since for u \(\ne\) 0 choosing w=u \(\lambda\) makes the last
parameter nonzero. It is also the derived group: with u=1, the map w
\(\mapsto\) bar(w)-w is a nonzero k-linear map onto \(L^0\) , and BF.13
supplies a root element for every w. The quotient by that centre is the
additive group of L.

If v=0, the trace relation forces u=0. For every pair in H,

\[
 2\omega(u)=\omega(u\bar u)
 =\omega(v+\bar v)\ge\omega(v).
 \tag{BF.16}
\]

At the identity both sides have their extended infinite values. Define

\[
 \begin{aligned}
 \varphi_{\pm a}(x_\pm(u,v))&=\tfrac12\omega(v),\\
 \varphi_{\pm2a}(x_\pm(0,v))&=\omega(v),
 \end{aligned}
 \tag{BF.17}
\]

with infinity only at the identity. A short-root level r is the set
\(\omega\) (v) \(\ge\) 2r. BF.16 then gives \(\omega\) (u) \(\ge\) r.
The cross term in BF.15 has valuation at least 2r for two such elements,
proving that this is a subgroup. Inversion preserves the value.
Long-root levels are additive balls in \(L^0\) , so are subgroups as
well. The short-root level is the closed subset H of the product of the
two compact balls \(\omega\) (u) \(\ge\) r and \(\omega\) (v) \(\ge\)
2r; it is compact and open in H. As r increases both coordinates tend to
zero, so these sets give exactly the coordinate neighborhood basis. The
long-root assertion follows identically in \(L^0\) . Nonidentity finite
value sets are discrete; they need not be all of (2e)\^{}\{-1\}Z. The
elements x\_±(0, \(\delta \pi^n\) ) give infinitely many values, while
the exact long-root set is \(\omega\) ( \(\delta\) )+Z.

The doubled-root identity is exactly \(\varphi_{2a}\) =2varphi\_a on
\(U_{2a}\) , and similarly negatively. BF.15--BF.16 also give

\[
 [U_{a,r},U_{a,s}]\subset U_{2a,r+s},\qquad
 [U_{-a,r},U_{-a,s}]\subset U_{-2a,r+s}.
 \tag{BF.18}
\]

Commutators with the corresponding doubled-root group vanish. These
exhaust the valuation commutator axiom: opposite-ray pairs are excluded
in that axiom, and a+a=2a is its only possible positive-integer root sum
on one ray, apart from the negative counterpart.

We now prove the opposite-factor axiom for every possible normalizer
factorization, including its uniqueness. The S-weight lines show that
its centralizer is diagonal and its normalizer either fixes or exchanges
the two outer lines. An automorphism of the one-dimensional split torus
sends its character to ±a: an invertible group morphism on k{[}t,
\(t^{-1}\) {]} has t \(\mapsto t^{±1}\) . Solving the determinant and
unitary equations consequently gives every normalizer matrix as one of

\[
 D(t)=\operatorname{diag}(t,\bar t/t,1/\bar t),\quad t\in L^*,
 \qquad
 M(v)=\begin{pmatrix}0&0&v\\0&-\bar v/v&0\\1/\bar v&0&0\end{pmatrix},
 \quad v\in L^*.
 \tag{BF.19}
\]

The outer column pairings of M(v) are 1, its middle column has norm 1,
and its determinant is 1, verifying membership directly.

For \(x_{+}\) (u,v) \(\ne\) 1 put V=bar(v). The two opposite factors are

\[
 y'=x_-(-u/v,1/V),\qquad
 y''=x_-(-u/V,1/V).
 \tag{BF.20}
\]

They belong to H because 1/V+1/v=-u bar(u)/(vV). Direct multiplication
first gives

\[
 x_+(u,v)y''=
 \begin{pmatrix}0&0&v\\0&-V/v&u\\1/V&\bar u/v&1\end{pmatrix},
\]

and then y' \(x_{+}\) (u,v)y'\,'=M(v). Both factors have value -
\(\omega\) (v)/2=- \(\varphi_{a}\) ( \(x_{+}\) (u,v)). They are the only
possible factors even if their product is merely required to belong to
N: in \(x_{-}\) (p,q) \(x_{+}\) (u,v) \(x_{-}\) (r,s), the (1,3) entry
is v \(\ne\) 0, so a normalizer product must be antidiagonal. Vanishing
of (1,2) forces r=-u/V, then (1,1) forces s=1/V, (2,3) forces p=-u/v and
(3,3) forces q=1/V. This proves the axiom for every permissible product.
For a long-root input u=0, V=-v; both opposite factors are \(x_{-}\)
(0,-1/v), with value - \(\omega\) (v). Conjugation by J proves both
negative-root cases. None of these computations divides by two.

The normalizer conjugations, with all parameters shown, are

\[
 \begin{aligned}
 D(t)x_+(w,z)D(t)^{-1}
   &=x_+(\bar t^2w/t,t\bar t z),\\
 D(t)x_-(w,z)D(t)^{-1}
   &=x_-(\bar t w/t^2,z/(t\bar t)),\\
 M(v)x_+(w,z)M(v)^{-1}
   &=x_-(-\bar v w/v^2,z/(v\bar v)),\\
 M(v)x_-(w,z)M(v)^{-1}
   &=x_+(-\bar v^2w/v,v\bar v z).
 \end{aligned}
 \tag{BF.21}
\]

Every right-hand pair satisfies the trace relation: the norm multiplier
on its first coordinate is exactly the multiplier on its second. The
short-root valuation changes for D(t) are respectively + \(\omega\) (t)
and - \(\omega\) (t); for M(v) they are respectively - \(\omega\) (v)
and + \(\omega\) (v). Long-root changes are twice these. Thus the
required normalizer difference is constant on each nonidentity root
group.

On the real apartment with a(h)=h, the action is

\[
 D(t):h\mapsto h-\omega(t),\qquad
 M(v):h\mapsto-h-\omega(v).
 \tag{BF.22}
\]

BF.21 verifies filtration covariance directly. This is a group action:
matrix multiplication gives

\[
 D(t)D(s)=D(ts),\quad D(t)M(v)=M(tv),
 \quad M(v)D(t)=M(v/\bar t),\quad M(v)M(w)=D(v/\bar w),
\]

and applying \(\omega\) checks all four affine compositions. For the
reflection in BF.20 put \(r=\varphi_a(x_+(u,v))\). Its action is
\(h\mapsto-h-2r\), with fixed wall \(a(h)+r=0\). The doubled-root wall
is \(2a(h)+\varphi_{2a}=0\) and gives the same reflection. For a
negative root the wall is \(-a(h)+\varphi_{-a}=0\), with the analogous
doubled-root formula.

At any other origin add a(h)=h to the positive short-root valuation,2h
to its double, and their negatives to the opposite roots. Subgroup and
topology assertions merely relabel their levels. In BF.18 the two
short-root shifts add to the doubled-root shift. In BF.20 the opposite
shifts are exactly the negative of the input shift. The normalizer
differences remain constant. Thus all valuation axioms and topology hold
at every real origin, with no ramification or characteristic
restriction. The real factor 1/2 in BF.17 never represents division by
two in k or L. \(\square\)

\subsection{Affine apartments and unramified
descent}\label{affine-apartments-and-unramified-descent}

\subsubsection{AB.0. Finite separable extensions of a local
field}\label{ab.0.-finite-separable-extensions-of-a-local-field}

\textbf{Theorem.} Let \(k\) be a nonarchimedean local field, with
complete discrete valuation ring \(O\), uniformizer \(\pi\), and
normalized valuation \(\omega_k\). Every finite separable extension
\(L/k\) has a unique real-valued valuation extending \(\omega_k\). Its
valuation ring is the integral closure \(B\) of \(O\) in \(L\). This
ring is a complete DVR, finite free of rank \([L:k]\) over \(O\). If
\(\Pi\) is its uniformizer and \(\pi=\varepsilon\Pi^e\) with
\(\varepsilon\in B^*\), then

\[
 \omega_L(\Pi)=1/e,\qquad \omega_L(L^*)=e^{-1}\mathbb Z.
 \tag{AB.1}
\]

Writing \(f=[B/(\Pi):O/(\pi)]\), one has \([L:k]=ef\). The extended
valuation topology is the finite-dimensional \(k\)-vector topology. In
particular \(L\) is complete, its valuation balls are compact, and every
\(k\)-automorphism of \(L\) preserves \(\omega_L\).

\textbf{Proof.} We first prove the finiteness assertion without assuming
a valuation on \(L\). Elements integral over a ring form a ring: if
\(x,y\) satisfy monic equations, the powers of \(x,y\) below those two
degrees span a finite module stable under both elements. The determinant
identity for the corresponding multiplication operators gives monic
equations for \(x+y\) and \(xy\). The same argument proves transitivity
of integrality by adjoining the finitely many coefficients of a monic
equation first. The ring \(O\) is integrally closed in \(k\): a monic
equation for an element of negative valuation has its leading term of
strictly smaller valuation than every other term, which cannot sum to
zero.

Choose a \(k\)-basis \(b_1,\ldots,b_n\) of \(L\) consisting of integral
elements. Such a basis is obtained by multiplying an arbitrary basis by
sufficiently large powers of \(\pi\), clearing the coefficients of its
finitely many minimal polynomials. The trace pairing on \(L\) is
nondegenerate. Here is the field argument: since \(k\) is infinite and
the extension is separable, a linear combination of a basis can be
chosen outside the finitely many hyperplanes on which two distinct
embeddings agree. It has \(n\) distinct conjugates and hence generates
\(L\). For its power basis, the matrix of evaluations at these
conjugates is a Vandermonde matrix with nonzero determinant; the trace
matrix is its transpose times itself and is invertible.

If \(x\in B\), then \(\operatorname{Tr}_{L/k}(xb_i)\in O\). Indeed each
conjugate of an integral element is integral, their sum is integral, and
a trace lies in \(k\), where integral elements belong to \(O\). The
nondegenerate trace pairing therefore places \(B\) inside the finite
free trace-dual lattice of \(\sum_i O b_i\), while that latter lattice
is contained in \(B\). A submodule of a finite free module over the DVR
is finite free, without assuming a generator list. Project a submodule
of \(O^n\) to its first coordinate. Its nonzero image ideal has a
generator of minimum valuation; lift that generator to the submodule.
Its free rank-one span splits off, since the first coordinate of every
other element is a unique multiple of the chosen generator. The
remaining kernel lies in O\^{}(n \(-\) 1), so induction proves finite
freeness. If the image ideal is zero, apply the induction directly to
that kernel. Thus \(B\) is finite free of rank \(n\), and is complete
for its \(\pi\)-adic topology. It is a domain and is integrally closed
in \(L\), by transitivity of integrality. Also \(B[\pi^{-1}]=L\).

We next show that \(B\) is local. Every maximal ideal contains \(\pi\).
Otherwise its nonzero residue field, a finite \(O\)-module, would
satisfy \(M=\pi M\). On finite generators this gives a matrix equation
with coefficient matrix \(1-\pi A\); its determinant is an \(O\)-unit
and annihilates \(M\), a contradiction. The finite ring \(B/\pi B\) is
an Artinian algebra over the finite residue field. If it had two maximal
ideals, it would have a nontrivial idempotent. Explicitly, its
nilradical is nilpotent: choose finitely many nilpotent generators and
use the pigeonhole principle on products. Chinese remainders for powers
of the finitely many distinct maximal ideals then express the ring as a
product of two nonzero rings and supply that idempotent.

Every idempotent modulo \(\pi B\) lifts to \(B\). If \(z^2-z\in\pi B\),
then \(2z-1\) is invertible: its square is \(1+4(z^2-z)\), and
\(1+\pi B\) consists of units by the convergent geometric series.
Newton's update

\[
 z\longmapsto z-(z^2-z)/(2z-1)
\]

squares the order of the error. Completeness gives an idempotent with
the specified reduction. This argument is valid in characteristic two. A
domain has only the idempotents zero and one, so \(B/\pi B\), and
therefore \(B\), is local.

For completeness we prove the DVR assertion too. A nonzero prime of
\(B\) contains \(\pi\): otherwise its localization would be a prime of
the field \(B[\pi^{-1}]\), and injectivity of that localization would
make it zero. Primes containing \(\pi\) are maximal, since the quotient
is a finite-dimensional residue algebra. Thus the sole nonzero prime is
the maximal ideal \(\mathfrak m\). Choose \(0\ne a\in\mathfrak m\). Then
\(\sqrt{(a)}=\mathfrak m\); finitely many generators give
\(\mathfrak m^d\subset(a)\) for some \(d\). For minimal \(d\), take
\(b\in\mathfrak m^{d-1}\setminus(a)\) and put \(u=b/a\). Then
\(u\notin B\) but \(u\mathfrak m\subset B\). If
\(u\mathfrak m\subset\mathfrak m\), the determinant identity on this
finite nonzero ideal makes \(u\) integral over \(B\), contrary to
normality. Hence \(u\mathfrak m=B\), and \(\mathfrak m=u^{-1}B\) is
principal. Write its generator as \(\Pi\). The intersection
\(J=\bigcap_d\Pi^dB\) satisfies \(J=\Pi J\) by cancellation in the
domain. It is finite, so the same determinant identity, with unit
determinant modulo \(\mathfrak m\), makes it zero. Repeated division by
\(\Pi\) now writes every nonzero element uniquely as a unit times a
nonnegative power of \(\Pi\). This is precisely a DVR.

Since \(\pi\) is a nonzero nonunit, it is \(\varepsilon\Pi^e\) for a
positive integer \(e\). Define \(\omega_L\) to be the \(\Pi\)-order
divided by \(e\). Every \(O\)-unit remains a \(B\)-unit, so its
restriction is \(\omega_k\). Any other extending valuation is
nonnegative on every integral element, by the same leading-term
argument; it has value zero on \(B^*\) and therefore value \(1/e\) on
\(\Pi\). It coincides with \(\omega_L\), proving uniqueness.

The filtrations \(\pi^dB\) and \(\Pi^dB\) are cofinal. In the finite
free \(O\)-basis the former is exactly the coordinate topology.
Completeness and compactness follow from completeness and compactness of
\(O\). Finally the \(\Pi\)-filtration of \(B/\pi B\) has \(e\)
quotients, each its residue field of dimension \(f\) over \(O/\pi O\).
On the other hand its dimension is \(n\), by freeness. Thus \(n=ef\). An
automorphism pulls the valuation back to an extending valuation;
uniqueness proves invariance. This establishes the exact extension
prerequisite used in BF.7. \(\square\)

The primary arithmetic source consulted is J. S. Milne, \emph{Algebraic
Number Theory}, the free author-hosted text, Proposition 2.29 and
Theorem 7.38. The proof above supplies the lattice, idempotent, DVR,
normalization and topology arguments locally; its valuation assertion is
not merely a citation to that theorem.

\subsubsection{AB.1. The split quotient and its proper
topology}\label{ab.1.-the-split-quotient-and-its-proper-topology}

Use the split local-field data of BF.0--BF.6, including the enlarged
apartment \(V=V_{\mathrm{der}}\oplus V_Z\), normalizer \(N\), compact
open groups \(K_x\), and their covariance and normalizer intersection.
Define

\[
 (g,x)\sim(h,y)
 \quad\Longleftrightarrow\quad
 y=nx\text{ and }g^{-1}hn\in K_x\text{ for some }n\in N.
 \tag{AB.2}
\]

\textbf{Theorem.} This relation is an equivalence relation. Its quotient
\(\mathcal B\) has a continuous transporter-proper \(G\)-action and is a
second countable locally compact Hausdorff space with compact orbit
space. The map \(j(x)=[1,x]\) is injective and restricts to a
homeomorphism on every compact subset of \(V\). Its point stabilizer is
\(G_{[g,x]}=gK_xg^{-1}\). A root element \(x_a(u)\) fixes every \(x\)
with \(a(x)+\omega(u)\ge0\).

\textbf{Proof of the set assertions.} Reflexivity uses \(n=1\). If
\(k=g^{-1}hn\in K_x\) and \(y=nx\), then
\(h^{-1}gn^{-1}=nk^{-1}n^{-1}\in K_y\), proving symmetry. For a second
relation, \(z=my\) and \(h^{-1}\ell m=k'\in K_y\), the required element
is \(mn\), because

\[
 g^{-1}\ell mn=(g^{-1}hn)(n^{-1}k'n)\in K_x.
\]

Covariance makes the last membership valid. Left translation preserves
the relation. If \([1,x]=[1,y]\), its normalizer witness belongs to
\(K_x\cap N=N_x\); hence \(y=x\). The same computation, with the
apartment coordinate unchanged, gives the displayed stabilizer. A root
element of nonnegative value belongs to \(P_x\subset K_x\), proving the
half-apartment assertion.

Here are the topology details; closedness of a relation alone would not
suffice. First the family \(K_x\) has the following local property:

\[
 x\text{ sufficiently close to }x_0\quad\Longrightarrow\quad K_x\subset K_{x_0}.
 \tag{AB.3}
\]

For each of the finitely many roots, the integer \(\lceil-a(x)\rceil\)
is locally constant unless \(a(x_0)\) is integral; at such a point its
nearby values are its value there or that value plus one. Thus all
nearby root level groups lie in their groups at \(x_0\), and
\(P_x\subset P_{x_0}\). By BF.3 only finitely many affine normalizer
labels can fix a point of a fixed compact neighborhood of \(x_0\). Each
label's fixed locus is closed. Shrink the neighborhood to exclude those
which do not fix \(x_0\); then \(N_x\subset N_{x_0}\). This proves AB.3.
In particular the graph of \(K_x\) is closed. Over a compact set
\(D\subset V\), the union \(K_D=\bigcup_{x\in D}K_x\) is compact: AB.3
first bounds it in finitely many compact groups, and the closed graph
over \(D\) then proves that this union is closed in that bound.

Choose a reduced closed chamber \(\overline C\) and a compact
parallelepiped \(D_Z\) for the lattice \(X_*(T)\cap V_Z\); put
\(D=\overline C\times D_Z\). The chamber and the lattice translates
cover \(V\) locally finitely, and \(\overline C\) is compact. The
complete elementary chamber proof is AB.2 below. The central lattice has
full rank in \(V_Z\), since that subspace is defined by integer root
equations. By BF.6 the open compact group \(I=I_C\) fixes every point of
\(D\). Consequently the set quotient is also the quotient of

\[
 S=(G/I)\times D.
 \tag{AB.4}
\]

This is a countable disjoint union of compact metrizable sets.
Countability follows because \(G\) has a countable base and its disjoint
open \(I\)-cosets each contain a different basis member.

The equivalence relation on \(S\) is closed and has proper projections.
Indeed only finitely many affine normalizer labels move a point of \(D\)
into \(D\), by BF.3. Their torus-unit kernel is already in every
\(K_x\), so choose finitely many representatives \(n\). For a fixed
representative \(g\) of a \(G/I\)-coordinate, all equivalent coordinates
\(hI\) have

\[
 h\in gK_Dn^{-1}.
\]

This is a finite union of compact sets, which meets only finitely many
open \(I\)-cosets. Closedness now follows from the closed graph just
proved, checking the finitely many normalizer representatives. The
preimage under either relation projection of a compact subset of \(S\)
is a closed subset of finitely many coordinate copies of \(D\times D\),
hence compact. The fibers are finite, because both the possible coset
coordinates and the affine images of the specified point are finite.

We record why this gives the required quotient topology. A relation with
proper projections has closed saturation of every closed set. To check
this here, a convergent sequence of saturated points lies in a compact
set after adding its limit; properness supplies convergent related
witnesses, and closedness supplies the limit relation. Thus the quotient
map \(q:S\to\mathcal B\) is closed. Its fibers are compact. Two distinct
fibers have disjoint neighborhoods with no related pair, by closedness
of the relation and compactness of their product. Replacing each
neighborhood \(U\) by the saturated open set
\(S\setminus q^{-1}q(S\setminus U)\) separates their quotient points.
Thus the quotient is Hausdorff.

A closed map with compact fibers pulls compact sets back to compact
sets: for an open cover of the preimage, cover each fiber by finitely
many members; closedness gives a quotient neighborhood whose entire
preimage is in that finite union, and compactness of the target set
reduces these neighborhoods to finitely many. Every fiber here lies in
finitely many open coordinate copies of \(D\), a compact set. The same
saturated-open construction gives a quotient neighborhood with compact
closure, proving local compactness. Finally the sets
\(\mathcal B\setminus q(S\setminus U)\), where \(U\) ranges over finite
unions of a fixed countable base of \(S\), form a countable quotient
base: a compact fiber inside an open preimage has such a finite base
cover. This proves second countability.

The topology from AB.4 is exactly the topology from \(G\times V\). In
one direction the map from \(G\times D\) factors through the open coset
map \(G\to G/I\). In the other direction, restrict to each
\(G\times nD\) in the locally finite closed normalizer cover of
\(G\times V\). On that set the map is, after replacing \(x\) by
\(n^{-1}x\), the continuous map \((g,x)\mapsto q(gnI,n^{-1}x)\). A
subset is closed if its restrictions to a locally finite closed cover
are closed, so the map from \(G\times V\) is continuous for the AB.4
topology. The two quotient topologies coincide.

The action is jointly continuous. At any quotient point, closedness of
\(q\) gives a neighborhood all of whose lifts lie in finitely many
prescribed coordinate copies of \(D\). An open subgroup of \(G\)
contained in the stabilizers of these finitely many \(G/I\)-coordinates
fixes that entire neighborhood. This proves continuity near \((1,b)\);
translation proves it everywhere. If compact subsets
\(A,B\subset\mathcal B\) have lifts with coset representatives
\(a_i,b_j\), then an element \(g\) with \(gA\cap B\ne\varnothing\)
satisfies

\[
 g\in\bigcup_{i,j,n} b_j n K_D a_i^{-1}.
\]

This is compact. The transporter is closed, because witnesses in the
compact sets have convergent subnet limits and the action is continuous.
It is therefore compact. Every orbit meets \(j(D)\), proving compactness
of the orbit space. The continuous injection \(j\) on any compact
apartment set is a homeomorphism onto its image by Hausdorffness.
\(\square\)

\subsubsection{AB.2. Affine reflection geometry from the root
datum}\label{ab.2.-affine-reflection-geometry-from-the-root-datum}

Let \(\Phi\) be the reduced root system of the split group in BF, on the
derived apartment \(V_D\). For \(a\in\Phi\) and \(m\in\mathbb Z\), set

\[
 H_{a,m}=\{v:a(v)+m=0\},\qquad
 s_{a,m}(v)=v-(a(v)+m)a^\vee.
 \tag{AB.5}
\]

Choose an invariant Euclidean inner product. Its existence follows by
averaging any inner product over the finite Weyl group, whose finiteness
and root-datum reflection action are proved in
\href{../AG-RG/AG-RG-04.html}{Root data, Weyl chambers and the Bruhat
decomposition}, §§2--3. The displayed map is the orthogonal reflection
in its wall. The arrangement is locally finite, since there are finitely
many roots and only finitely many integers \(m\) whose walls meet a
given bounded set. Its chambers are open convex intersections of wall
half-spaces. A chamber is bounded: for every root \(a\), its values on a
chamber lie between two consecutive integers, and the roots span
\(V_D^*\). Its closure is consequently a compact convex polytope. In
rank zero take the chamber to be the one-point apartment.

\textbf{Lemma.} If \(C\) is a chamber, the reflections in its walls
generate \(W_{\mathrm{aff}}=Q^\vee\rtimes W\). This group acts simply
transitively on chambers. Those reflections, denoted \(S\), have exactly
the Coxeter relations \(s^2=1\) and \((st)^{m(s,t)}=1\) for the finite
pair orders; an infinite pair order supplies no relation. For
\(w\in W_{\mathrm{aff}}\), its length is the number of walls separating
\(C\) from \(wC\). If \(v\in\overline C\), its stabilizer is generated
by \(J_v=\{s\in S:s(v)=v\}\). No two distinct points of \(\overline C\)
are in the same \(W_{\mathrm{aff}}\)-orbit.

\textbf{Proof.} The reflection preserves the entire arrangement, because

\[
 b(s_{a,m}v)=s_a(b)(v)-m\langle b,a^\vee\rangle
\]

and the Cartan integer is integral. At a wall of \(C\) it interchanges
the two adjacent chambers. A generic polygonal path between interior
points crosses finitely many walls, crosses them one at a time, and
misses their pairwise intersections. Starting at \(C\), reflect across
each crossed wall. At the current image \(wC\), that reflection is
\(wsw^{-1}\) for a wall reflection \(s\) of \(C\). Thus these
reflections reach every chamber.

We supply the elementary gallery argument that also proves freeness and
the presentation. Realize a closed chamber gallery by a polygonal loop
crossing each prescribed common wall in its relative interior. Fill the
loop by a polygonal disc in \(V_D\). Only finitely many arrangement
hyperplanes meet its compact image. Subdivide and perturb the interior
vertices, keeping the boundary crossings, so that the disc misses
intersections of codimension at least three, meets codimension-two
intersections at finitely many interior points, and is transverse to
individual walls away from those points. These perturbations exist by
avoiding the finitely many proper polynomial conditions on the vertices;
subdivisions permit small perturbations that preserve the boundary. For
one-dimensional \(V_D\), an interval path gives just backtracking
cancellations. In higher dimensions, sweeping the disc by polygonal arcs
from the boundary loop to a constant arc changes the crossed-wall word
in only two ways. A tangency inserts or removes two consecutive
crossings of the same wall. Passing a codimension-two intersection
replaces one route around that intersection by the other route. This
description can also be obtained by cutting small disjoint discs around
the intersection points: the remaining wall curves are disjoint arcs,
across which consecutive crossings cancel.

At a codimension-two intersection, use its perpendicular two-dimensional
plane. The walls through the intersection form a finite reflection
arrangement: the corresponding reflections fix that point, and their
linear parts are in the finite group \(W\). Consecutive wall lines bound
a sector. If their angle is \(\theta\), their reflections generate a
finite dihedral group, and its lines divide the plane into sectors of
angle \(\pi/m\). Since the two original lines were consecutive in the
full invariant arrangement, no generated line lies inside their sector;
hence \(\theta=\pi/m\). Crossing around the intersection gives precisely
the alternating braid relation of length \(m\), or equivalently
\((st)^m=1\). Pulling the current chamber back to \(C\) expresses this
in its two wall generators. Thus every closed gallery word reduces by
the stated relations. These relations hold as actual reflections. An
element preserving \(C\) therefore has a closed gallery word and is the
identity; the same argument shows that every relation is generated by
the listed ones. This proves both simple transitivity and the Coxeter
presentation, without assuming them as an external building theorem.

Every wall is a wall of some chamber, so its reflection is now a
conjugate of a generator. All arrangement reflections generate exactly
\(Q^\vee\rtimes W\): their products satisfy
\(s_{a,m}s_{a,0}=t_{-ma^\vee}\), giving all coroot translations, while
the reflections with \(m=0\) generate \(W\); conversely each displayed
reflection is in this semidirect product.

A gallery must cross every wall separating its endpoints. A generic
straight segment crosses each separating wall once and crosses no other
wall, so the minimal gallery length is exactly their number. This proves
the length assertion, including \(\ell(sw)=\ell(w)\pm1\).

The chambers whose closures contain a fixed point \(y\) form its local
star. In a sufficiently small ball around \(y\), only the walls through
\(y\) occur. Their finitely many sectors are connected by crossing
consecutive walls, so any two star chambers are joined by a gallery
crossing only those walls. If \(w\) fixes \(v\in\overline C\), apply
this to \(C,wC\). Pulling successive crossings back to \(C\) gives only
generators in \(J_v\). Simple transitivity then gives
\(w\in\langle J_v\rangle\); the reverse containment is immediate. If
\(wv=y\) with \(v,y\in\overline C\), the same local-star argument
supplies an element \(u\) fixing \(y\) and sending \(C\) to \(wC\).
Freeness gives \(u=w\), whence \(v=y\). Point stabilizers are finite:
for each finite linear Weyl element there is at most one translating
vector making it fix that point. \(\square\)

\subsubsection{AB.3. Split affine wall exchange and the generated
group}\label{ab.3.-split-affine-wall-exchange-and-the-generated-group}

Keep the pinned split group, the valuation and the enlarged apartment of
BF. Write \(T^b=T(O)\), let \(I=I_C\) be BF.6's compact open chamber
group, and let \(N_{\mathrm{aff}}\) be the inverse image of
\(W_{\mathrm{aff}}\) in \(N_G(T)(k)\). Thus
\(N_{\mathrm{aff}}/T^b=W_{\mathrm{aff}}\). For a wall \(\alpha=a+m\) of
\(C\), positive on \(C\), choose \(\dot s=n_a(\pi^m)\), whose action is
\(s_{a,m}\). Its square is \(a^\vee(-1)\in T^b\).

\textbf{Lemma.} For every \(n\in N_G(T)(k)\),

\[
 I\dot s I n I\ \subset\ I\dot s n I\ \cup\ I n I.
 \tag{AB.6}
\]

Moreover \(I\cap N_G(T)(k)=T^b\), \(\dot s I\dot s^{-1}\ne I\), and

\[
 G^0:=\langle I,N_{\mathrm{aff}}\rangle
     =\langle T^b,U_a(k):a\in\Phi\rangle
\]

is an open normal subgroup of \(G(k)\). The full group is \(T(k)G^0\).
Its core under normalizer conjugation is
\(\bigcap_{n\in N}nIn^{-1}=T^b\).

\textbf{Proof.} On \(C\) the affine root \(\alpha\) takes values between
zero and one. BF.6 therefore gives the levels \(f_C(a)=m\),
\(f_C(-a)=1-m\). Crossing this wall changes only these two root levels:
a wall for any other root does not separate the two adjacent chambers.
Choose a positive root system containing \(a\), and order its positive
factors with \(a\) last. The arbitrary-order factorization in BF.6
writes

\[
 I=J U_{a,m},\qquad J\subset I\cap\dot s^{-1}I\dot s.
 \tag{AB.7}
\]

Indeed the torus and every factor except \(U_{a,m}\) have levels that
remain admissible after reflection. Consequently it suffices to examine
\(\dot s x_a(u)n\), with \(\omega(u)\ge m\). If \(\omega(u)>m\),
conjugation by \(\dot s\) puts this root factor in
\(U_{-a,1-m}\subset I\), giving the first cell in (AB.6). If
\(\omega(u)=m\) and \(n^{-1}\alpha\) is positive on \(C\), conjugation
by \(n^{-1}\) puts \(x_a(u)\) in \(I\), giving that same cell. An affine
root has constant sign on \(C\), so the remaining case is that
\(n^{-1}\alpha\) is negative.

Put \(c=\pi^m\). Multiplication of the two-by-two matrices in the pinned
rank-one map gives the identity

\[
 n_a(c)x_a(u)
 =x_a(-c^2/u)\,a^\vee(-c/u)\,x_{-a}(1/u).
 \tag{AB.8}
\]

This is valid in every residue characteristic, including two. The first
parameter has valuation \(m\), the torus parameter is a unit, and the
last factor has the level of \(-\alpha\). In the remaining case
\(n^{-1}(-\alpha)\) is positive, so that last factor conjugates into
\(I\). The result lies in \(InI\). This proves (AB.6) with its signs and
integer thresholds.

BF.5 gives \(K_v\cap N=N_v\). Every element of \(I\) fixes every
\(v\in C\times V_Z\). An element of \(I\cap N\) consequently has affine
action equal to the identity on an open set and hence everywhere. BF.3
identifies its kernel as \(T^b\). Conversely \(T^b\subset I\).
Reflection conjugation carries \(U_{a,m}\) to \(U_{-a,-m}\); the latter
has a parameter of valuation \(-m\) and is not in \(I\), whose
opposite-root level is \(1-m\), by BF.6's uniqueness and root
intersections. Thus the conjugated groups differ.

Let \(R=\langle T^b,U_a(k)\rangle\). Every \(n_a(u)\) is a product of
opposite root elements by the pinned rank-one formula, so
\(I,N_{\mathrm{aff}}\subset R\). Conversely conjugating the root tails
in \(I\) by powers of a coroot translation lowers their levels without
bound, since \(\langle a,a^\vee\rangle=2\). Those tails exhaust
\(U_a(k)\), proving \(R=G^0\). It is open since it contains \(I\). The
preceding programme proof \href{../AG-RG/AG-RG-04.html}{Root data, Weyl
chambers and the Bruhat decomposition}, Theorem 6.2, writes every
element of the split group in a finite Bruhat cell. Its unipotent
coordinates are root elements and its Weyl representative is a product
of \(n_a(1)\). Therefore \(G=T(k)G^0\). The torus normalizes \(R\), and
the root generators belong to it, so it is normal in \(G\).

Finally an element of the intersection of all normalizer conjugates
belongs to \(I\). Write its unique BF.6 negative-root, torus and
positive-root coordinates. Torus conjugation preserves those coordinates
and scales each root coordinate by the corresponding root character. For
each nonzero parameter, a power of a cocharacter with nonzero pairing
lowers its valuation below the fixed chamber threshold, contradicting
membership in all torus conjugates of \(I\). Thus every root parameter
is zero. The element lies in \(T^b\), and this subgroup is normal in
\(N\), giving equality. \(\square\)

\subsubsection{AB.4. Double cosets, proved from the wall
exchange}\label{ab.4.-double-cosets-proved-from-the-wall-exchange}

We prove the group-theoretic argument needed for the building, instead
of citing a BN-pair theorem. Suppose a group \(H\) is generated by
subgroups \(B,N_0\), with \(T=B\cap N_0\) normal in \(N_0\). Suppose
\(W_0=N_0/T\) is generated by nonidentity involutions \(S\), and, for
representatives of \(s\in S\),

\[
 B\dot s B\dot w B\subset B\dot w B\cup B\dot s\dot w B,
 \qquad \dot s B\dot s^{-1}\ne B.
 \tag{AB.9}
\]

The double coset \(D_w=B\dot wB\) is independent of its representative.
Denote word length in \(S\) by \(\ell\).

\textbf{Lemma.} The double cosets \(D_w\) partition \(H\), their labels
are distinct, and

\[
 D_sD_w=
 \begin{cases}
 D_{sw},&\ell(sw)>\ell(w),\\
 D_w\cup D_{sw},&\ell(sw)<\ell(w).
 \end{cases}
 \tag{AB.10}
\]

The length changes by exactly one. Products along a reduced word equal
the double coset of its label. For \(J\subset S\),

\[
 P_J:=\langle B,\dot s:s\in J\rangle
       =\coprod_{u\in W_J}D_u,
 \qquad D_wP_J=\bigcup_{u\in W_J}D_{wu}.
 \tag{AB.11}
\]

\textbf{Proof.} First \(D_s^2\subset B\cup D_s\), and it contains \(B\),
because \(\dot s^{-1}\in D_s\). If it were contained in \(B\), then
\(\dot sB\dot s^{-1}\subset B\); since \(\dot s^2\in T\), conjugating
again yields the reverse inclusion. This contradicts (AB.9). Hence
\(D_s^2=B\cup D_s\). The union of all \(D_w\) is closed under
multiplication by the generators \(B\) and \(D_s\), using (AB.9), and
under inverses. It is therefore \(H\).

For distinctness, suppose \(D_v=D_w\), and induct on
\(\min(\ell(v),\ell(w))\). If this is zero, equality with \(B\) places
the normalizer representative in \(B\cap N_0=T\), so both labels are the
identity. Otherwise choose the shorter label \(v=sv'\), with
\(\ell(v')=\ell(v)-1\). Since \(\dot s^{-1}\in D_s\),
\(D_{v'}\subset D_sD_v=D_sD_w\subset D_w\cup D_{sw}\). Double cosets
either coincide or are disjoint, so \(D_{v'}\) equals one of the two
terms. Induction gives \(v'=w\) or \(v'=sw\). The first contradicts the
choice of the shorter length; the second gives \(v=w\).

We next prove the ascending product rule by induction on \(\ell(w)\),
assuming only \(\ell(sw)\ge\ell(w)\). The identity case is immediate.
Write \(w=w't\) reduced, with \(t\in S\). Since
\(\ell(sw)\le\ell(sw')+1\), one has \(\ell(sw')\ge\ell(w')\). Induction,
followed by the right-hand version of (AB.9), obtained by inversion,
gives

\[
 D_sD_w\subset D_sD_{w'}D_t
       =D_{sw'}D_t\subset D_{sw'}\cup D_{sw}.
\]

It also lies in \(D_w\cup D_{sw}\). The labels \(sw\) and \(w\) differ.
Also \(sw'\ne w\): equality would give \(sw=w'\), contradicting
\(\ell(sw)\ge\ell(w)\). Distinctness therefore leaves precisely
\(D_{sw}\), which is contained in the product by multiplication of
representatives. This proves the ascending rule.

If \(\ell(sw)\le\ell(w)\), the ascending rule applied to \(sw\) gives
\(D_sD_{sw}=D_w\). Multiplying once more gives
\(D_sD_w=(B\cup D_s)D_{sw}=D_{sw}\cup D_w\). Equality of the two lengths
would give both formulas and contradict distinctness. Since a generator
changes word length by at most one, its change is exactly one.
Right-hand products follow by inversion; repeated ascending products
prove the reduced-word assertion. The union of labels in \(W_J\) is a
subgroup by the product formulas and equals the generated \(P_J\).
Iterating the right product formula shows that \(D_wP_J\) is contained
in the union in (AB.11); conversely multiplication of representatives
puts each \(D_{wu}\) in that product. \(\square\)

Apply this lemma with \(H=G^0,B=I,N_0=N_{\mathrm{aff}}\), using
AB.2--AB.3. It gives a bijection
\(W_{\mathrm{aff}}\to I\backslash G^0/I\). For the full normalizer
quotient \(\widetilde W=N/T^b\), put
\(\Omega=\{w\in\widetilde W:wC=C\}\), where stabilization refers to the
reduced chamber, allowing central translations. Each label is uniquely
\(\omega w\), \(\omega\in\Omega,w\in W_{\mathrm{aff}}\), by AB.2's
simple transitivity. Such \(\omega\) preserves all chamber root levels,
so its representative normalizes \(I\); \(W_{\mathrm{aff}}\) is normal
in \(\widetilde W\) by the lattice and finite-root action. Since
\(G=T(k)G^0\), every element is in a cell \(InI\).

We verify distinctness also for these extended labels; this is needed to
identify \(G^0\cap N\). The preceding distinctness induction applies to
a comparison of an affine label \(v\) with any label
\(w\in\widetilde W\): induct on \(\ell(v)\), using (AB.6) for the
possibly nonaffine label. The identity case still uses \(I\cap N=T^b\).
In the inductive step equality gives \(v'=w\) or \(v'=sw\). In the first
case \(w=v'\) is already affine, so the proved affine distinctness
contradicts the shorter length; in the second case \(w=sv'=v\). Thus
equality forces \(w=v\). Thus \(G^0\cap N=N_{\mathrm{aff}}\). If
\(D_{\omega u}=D_{\omega'v}\), left multiplication by \(\omega^{-1}\),
which normalizes \(I\), compares the affine label \(u\) with
\(\omega^{-1}\omega'v\). The comparison just proved forces
\(\omega^{-1}\omega'\in W_{\mathrm{aff}}\). It also belongs to
\(\Omega\), whose intersection with \(W_{\mathrm{aff}}\) is trivial by
AB.2. Thus \(\omega=\omega'\) and \(u=v\). We have proved

\[
 \widetilde W\ \xrightarrow{\sim}\ I\backslash G/I.
 \tag{AB.12}
\]

This argument uses the freely available primary author exposition of
Todd Trimble on double-coset exchange as a mathematical source. The
entire argument actually used is written above; its omitted
Coxeter-criterion proof is not used, since AB.2 proves the required
geometry directly.

\subsubsection{AB.5. Common apartments, face groups and the enlarged
product}\label{ab.5.-common-apartments-face-groups-and-the-enlarged-product}

\textbf{Proposition.} Any two points of the split quotient of AB.1 lie
in an apartment. Any two chambers do also. For \(v\in\overline C\),

\[
 P_v=\coprod_{w\in W_v}I\dot wI,
 \qquad W_v=\langle J_v\rangle.
 \tag{AB.13}
\]

There is a homeomorphism

\[
 \mathcal B\simeq\mathcal B^0\times V_Z,\qquad
 \mathcal B^0=((G^0/I)\times\overline C)/\sim,
 \tag{AB.14}
\]

where \((gI,v)\sim(hI,v)\) exactly when \(g^{-1}h\in P_v\). Each reduced
apartment has its Euclidean affine-Coxeter tiling.

\textbf{Proof.} Translate point coordinates into
\(\overline C\times V_Z\) using \(N_{\mathrm{aff}}\). If their group
coordinates are \(g,h\), write \(g^{-1}h=i_1ni_2\) by (AB.12). The group
\(I\) fixes both point coordinates. Therefore the two points are in the
apartment \(gi_1V\): the first is \([gi_1,x]\), the second
\([gi_1,ny]\). The identical argument with chamber interiors proves
common apartments for chambers.

By definition \(P_v\) is generated by \(T^b\) and root elements of
nonnegative value, so \(I\subset P_v\subset G^0\). If \(p\in P_v\) lies
in the affine Bruhat cell \(I\dot wI\), then \(\dot w\in P_v\), and BF.5
implies \(w(v)=v\). AB.2 identifies such labels as
\(W_v=\langle J_v\rangle\). Conversely each wall representative for
\(J_v\) is a product of opposite root elements of value zero at \(v\),
by BF.1's rank-one formula. It lies in \(P_v\). This proves (AB.13).

Every point of \(\mathcal B\) is represented by \((g^0,v,z)\) with
\(g^0\in G^0,v\in\overline C,z\in V_Z\), since the full group is
\(\Omega G^0\) and \(\Omega\) preserves the reduced chamber. Consider
two such representatives. The equivalence relation gives an \(n\in N\)
carrying \((v,z)\) to \((v',z')\), with \((g^0)^{-1}h^0n=p n_0\),
\(p\in P_v,n_0\in N_{(v,z)}\). Hence
\(n n_0^{-1}\in G^0\cap N=N_{\mathrm{aff}}\). Its action carries
\((v,z)\) to \((v',z')\), so AB.2 gives \(v=v'\), while its central
translation is zero and gives \(z=z'\). Its finite stabilizer label is
in \(W_v\), whose representatives lie in \(P_v\) by (AB.13). Thus the
group-coordinate relation is precisely \((g^0)^{-1}h^0\in P_v\). The
converse is immediate. This proves the set bijection (AB.14).

For topology, \(P_v\) is compact, contains the open group \(I\), and has
finite index over it. The level description and finite stabilizers give
the same local upper-containment and closed-relation argument as AB.1
for the quotient defining \(\mathcal B^0\). It is a closed quotient with
finite fibers, locally compact Hausdorff and second countable. The map
(AB.14) is continuous. It is proper: a compact set in \(\mathcal B\)
has, by AB.1, lifts in finitely many \(G/I\)-coordinate copies of \(D\).
For each of their finitely many representatives write \(g=\omega g^0\).
Conjugation by \(\omega\) preserves \(I,G^0\), preserves
\(\overline C\), and translates \(V_Z\) by one fixed vector. Thus all
the corresponding preimages lie in finitely many chamber copies times
one bounded closed central set, a compact set. The preimage is closed
there. A proper continuous bijection between these locally compact
Hausdorff spaces is a homeomorphism: images of closed sets are closed,
checked on compact neighborhoods. This proves (AB.14).

The map from \(G\) to the central translations factors through
\(G/G^0=N/N_{\mathrm{aff}}\). It is the central projection of the torus
valuation translation of BF.3; the Weyl group acts trivially on \(V_Z\).
Consequently the full group acts on (AB.14) by its chamber action on
\(\mathcal B^0\) and these central translations. \(\square\)

\subsubsection{AB.6. The metric from affine chamber
data}\label{ab.6.-the-metric-from-affine-chamber-data}

This paragraph is a general geometric theorem, with explicit inputs. Let
\(W_{\mathrm{aff}}\) act on a finite-dimensional Euclidean apartment
with the locally finite alcove tiling proved in AB.2. Let \((H,B,N_0)\)
satisfy the exchange hypotheses of AB.4, with this affine Weyl group,
and let its face groups be the subgroups \(P_{J_v}\) in (AB.11). Form
\(X=((H/B)\times\overline C)/\sim\), using those face groups. No residue
field, compactness or valuation hypothesis is assumed in this paragraph.
Apartments are the translates of the chamber-tiling apartment supplied
by \(N_0\). These inputs have all been proved for \(\mathcal B^0\)
above.

Every two chambers lie in an apartment: if \(g^{-1}h=b_1nb_2\), both are
in \(gb_1A\), just as in AB.5. Face equivalence and (AB.11) also prove
that the base apartment is embedded. More explicitly, for a fixed
\(v\in\overline C\), labels \(w,w'\) give the same face point exactly
when \(w^{-1}w'\in W_{J_v}\), which is exactly the equality of the
Euclidean points \(wv,w'v\), by AB.2.

Define a chain distance: take the infimum of the sum of Euclidean
lengths of finitely many successive segments, each lying in a closed
chamber. There is always such a chain because two points lie in a common
apartment and a straight segment meets finitely many chambers. To show
that this is a metric and that apartments retain their Euclidean
distance, we prove the required retraction rather than presupposing a
building metric.

For the base chamber, a chamber \(gB\) has a unique Bruhat label \(w\).
Send its point with chamber coordinate \(v\) to the Euclidean point
\(wv\) in the base apartment. This is independent of all face
representatives: by (AB.11), the possible labels of \(gP_{J_v}\) are
exactly \(wW_{J_v}\), and each of these labels gives the same point
\(wv\). This map \(r_C\) is an isometry on each closed chamber and
therefore does not increase chain lengths. It fixes the base apartment.
The same construction, conjugated, gives a retraction centered at every
chamber onto every apartment containing that chamber. Indeed any
apartment containing the base chamber is \(bA\) for some \(b\in B\): in
\(gA\), that chamber has coordinate \(nB\), hence \(gn\in B\). Bruhat
labels are unchanged by left multiplication by \(b\), so the retraction
restricts to an isometry on every such apartment.

If \(x,y\) are in an apartment \(A\), choose a center chamber in \(A\).
Applying its retraction to any chain gives a length at least \(|x-y|\).
The straight segment in \(A\) gives the reverse inequality. Thus

\[
 d(x,y)=|x-y|\quad\text{whenever }x,y\in A.
 \tag{AB.15}
\]

Any distinct two points lie in a common apartment and have positive
distance there. Thus the chain distance is a metric. Straight apartment
segments are geodesics. All the retractions are 1-Lipschitz in this
metric. We shall need the following stronger equality. If \(p\) lies in
the closed center chamber, and \(z\) is any point, take an apartment
containing that chamber and a chamber containing \(z\). The retraction
is an isometry there and fixes \(p\). Hence

\[
 d(p,z)=|p-r_C(z)|\qquad(p\in\overline C).
 \tag{AB.16}
\]

All assertions here follow from the displayed affine data and the
locally proved double-coset identities. There is no additional apartment
axiom assumed implicitly.

\subsubsection{AB.7. CAT(0), completeness and comparison of the two
topologies}\label{ab.7.-cat0-completeness-and-comparison-of-the-two-topologies}

\textbf{Theorem.} The metric in AB.6 is geodesic and CAT(0). It is
complete even when the number of chambers at a panel is infinite. For
the split group over a local field it is proper and induces precisely
the quotient topology in AB.5. The enlarged product with the Euclidean
central factor has all these properties, and the full group acts
isometrically.

\textbf{Proof.} Let \(c(t)\), \(0\le t\le1\), be the straight segment
from \(x\) to \(y\) in a common apartment \(A\). For each \(t\), choose
a chamber of \(A\) whose closure contains \(c(t)\), and retract onto
\(A\) with that chamber as center. Write \(r\) for this retraction. By
(AB.16), \(d(z,c(t))=|rz-c(t)|\). The Euclidean squared-distance
identity and 1-Lipschitz property give

\[
 d(z,c(t))^2
 \le(1-t)d(z,x)^2+t d(z,y)^2-t(1-t)d(x,y)^2.
 \tag{AB.17}
\]

This is the argument of the freely accessible original Bruhat--Tits I,
§3.2.1, with the distance equality required for its center proved
explicitly in AB.6. It works at walls and vertices as well as interior
points.

Any midpoint of another geodesic from \(x\) to \(y\) has distances
\(d(x,y)/2\) to both endpoints. Substituting it for \(z\) in (AB.17),
with \(t=1/2\), shows that it coincides with \(c(1/2)\). Repeating on
dyadic subsegments and using continuity proves uniqueness of the
geodesic. Apartment intersections are convex, since both apartments
contain their identical unique segment between any two common points.

For completeness of the curvature assertion, let \(p=c_{xy}(t)\),
\(q=c_{xz}(s)\), and put \(A=d(x,y), B=d(x,z), C=d(y,z)\). Apply (AB.17)
first along \(xy\), then along \(xz\) to bound \(d(y,q)^2\), while
\(d(x,q)=sB\). It yields

\[
 d(p,q)^2\le t^2A^2+s^2B^2-ts(A^2+B^2-C^2).
 \tag{AB.18}
\]

The right side is exactly the squared distance between the corresponding
points in the Euclidean comparison triangle, by its cosine identity. Any
two boundary points are on two edges with a common vertex; relabeling
proves the full CAT(0) triangle comparison, including degenerate
triangles.

Here is a completeness proof which makes no local-finiteness assumption
on the chamber multiplicities. The chamber coordinate defines a type map
\(\lambda:X\to\overline C\). It is well defined at faces and is an
isometry on each closed chamber, hence is 1-Lipschitz for the chain
metric. Fix \(y\in\overline C\). In the standard apartment there are
finitely many walls meeting the closed unit ball about \(y\). Set

\[
 r(y)=\tfrac12\min\bigl(\{1\}\cup
 \{\operatorname{dist}(y,H):H\cap\overline B(y,1)\ne\varnothing, y\notin H\}\bigr)>0.
 \tag{AB.19}
\]

For a point \(p\) of type \(y\), take a common apartment for a chamber
containing \(p\) and a chamber containing any \(q\) with
\(d(p,q)\le r(y)\). Formula (AB.15) identifies this distance with the
Euclidean distance there. The ball crosses no wall not passing through
\(p\). Thus \(q\) lies in the closure of a chamber containing the
minimal face of \(p\): its signs at every wall not through \(p\) are
those of that face, while the walls through \(p\) only select a
neighboring chamber. Each such chamber contains \(p\), and its type map
is injective. In particular distinct points of type \(y\) are at
distance greater than \(r(y)\).

If \(p_n\) is Cauchy, then \(y_n=\lambda(p_n)\) converges to some \(y\)
in the compact polytope \(\overline C\). Choose a closed chamber
containing \(p_n\), and in that same chamber choose its point \(q_n\) of
type \(y\). One has \(d(p_n,q_n)=|y_n-y|\to0\). The sequence \(q_n\) is
Cauchy, and its uniform separation just proved makes it eventually
constant. Hence \(p_n\) converges. This proves completeness with
arbitrary panel multiplicity; it does not assert that affine chamber
data for a particular nonsplit group have already been constructed.

For the split local-field construction, every wall group
\(P_{\{s\}}=I\cup I\dot sI\) is contained in the compact group \(P_v\)
at a point of that wall, so it has finitely many right \(I\)-cosets. A
reduced product in (AB.10) therefore shows that each \(I\dot wI/I\) is
finite. Fix \(p\) in the closed base chamber. Any chamber meeting a
metric ball of radius \(R\) about \(p\) lies with the base chamber in an
apartment. Retracting onto the base apartment sends it isometrically to
\(w\overline C\), meeting that same ball. The diameter of
\(\overline C\) is finite and the Euclidean tiling is locally finite, so
there are only finitely many such labels \(w\). For each label there are
only finitely many chambers, by the finite coset assertion. Thus each
bounded ball is covered by finitely many closed chambers. Those chambers
are compact for the metric, by (AB.15), so a closed ball is compact.
This proves properness directly.

The quotient-to-metric identity is continuous, since its restrictions to
the chamber-coordinate copies are isometries and the quotient has that
defining topology. A neighborhood ball of radius one is contained in a
finite union \(K\) of closed chambers. This \(K\) is compact for the
quotient topology and Hausdorff for the metric, so the continuous
identity on \(K\) is a homeomorphism. At its center, metric convergence
is therefore quotient convergence. These neighborhoods prove the reverse
continuity everywhere, establishing equality of the two topologies.

For the enlarged product define

\[
 d_{\mathrm{enl}}((p,z),(q,z'))^2=d(p,q)^2+\|z-z'\|^2.
 \tag{AB.20}
\]

Product segments are geodesics. Adding (AB.17) to its Euclidean equality
gives the same strong convexity inequality, and (AB.18) proves CAT(0).
Completeness and properness follow from those of the two factors; their
metric topology is the product topology of (AB.14). The group preserves
apartment Euclidean distances: on the derived factor this follows from
its chamber action and the normalizer's affine Weyl action, and on the
central factor from BF.3's translations. The chain metric, and therefore
(AB.20), is invariant. \(\square\)

\subsubsection{AB.8. The split enlarged building is a universal proper
space}\label{ab.8.-the-split-enlarged-building-is-a-universal-proper-space}

The action is proper and cocompact by AB.1 and AB.14, and AB.7 supplies
the compatible proper CAT(0) metric. We give the fixed-set and
universal-map argument, since contractibility of the building alone
would not prove universality.

\textbf{Proposition.} For every compact subgroup \(H\subset G\),
\(\mathcal B^H\) is nonempty and contractible. For every proper
\(G\)-space carrying an invariant partition of unity subordinate to
slice maps to \(G/H_i\), with \(H_i\) compact, there is a \(G\)-map to
\(\mathcal B\), unique up to \(G\)-homotopy. Thus \(\mathcal B\) is a
model for the numerable universal proper \(G\)-space.

\textbf{Proof.} A compact subgroup has a bounded orbit \(E\). The
function \(F(b)=\sup_{e\in E}d(b,e)^2\) is continuous and tends to
infinity when \(b\) tends to infinity, since \(E\) is bounded.
Properness gives a minimum. Formula (AB.17), at the midpoint, makes its
minimizer unique. The function is \(H\)-invariant, so the minimizer is
fixed. The fixed set is closed and convex, because isometries preserve
the unique segments. Contract it along the unique segments to one fixed
point.

The segment operation is continuous. Indeed CAT(0) comparison for two
segments with the same initial point gives distance at fraction \(t\) at
most \(t\) times the distance of their endpoints. Inserting the
intermediate segment with one endpoint from each of two segments and
applying the triangle inequality yields
\(d(c_{xy}(t),c_{x'y'}(t))\le(1-t)d(x,x')+t d(y,y')\). Together with the
constant-speed formula this gives joint continuity in endpoints and
parameter.

The model itself is numerable. Point stabilizers \(K_p\) are compact
open. Its orbit \(Gp\) is discrete and closed: properness makes the map
\(G/K_p\to\mathcal B\) proper, and a compact subset of the discrete
space \(G/K_p\) is finite. Thus choose \(\delta>0\) with no other orbit
point within distance \(\delta\) of \(p\). If \(\varepsilon<\delta/3\),
the translates of \(B(p,\varepsilon)\) intersect only when their
translating cosets in \(G/K_p\) agree, since an intersection implies
\(d(p,gp)<2\varepsilon\). This gives an equivariant slice map from its
saturated open set to \(G/K_p\). The orbit space has the metric
\(\bar d(Gx,Gy)=\inf_gd(x,gy)\): distinct closed orbits have positive
distance, since a sequence \(gy\to x\) in a bounded ball has a
convergent transporter subsequence by properness; its zero-distance
limit would belong to that orbit. Metric balls are precisely the images
of saturated unions of ordinary balls, so this metric gives the quotient
topology. This quotient, denoted \(Q\), is compact by AB.1. Finitely
many slice neighborhoods cover it; denote their smaller neighborhoods by
\(U_i\). Shrink them to a finite open cover with closures still inside
the original neighborhoods: a sufficiently small uniform ball radius
does this on a compact metric space. The functions
\(\min(1,\operatorname{dist}(x,Q\setminus U_i))\), with value one if
that complement is empty, divided by their positive sum, give a finite
partition of unity whose closed supports are in the original ones.
Pulling it back gives the invariant subordinate partition. Thus the
model meets the numerability requirement as well as properness.

For the universal map, choose a point \(b_i\in\mathcal B^{H_i}\) for
each slice map \(p_i:Y_i\to G/H_i\). It defines a continuous equivariant
point map \(z_i:Y_i\to\mathcal B\) by \(gH_i\mapsto gb_i\). Let
\(\lambda_i\) be the subordinate invariant locally finite partition of
unity, with support contained in \(Y_i\). At \(y\), minimize
\(\sum_i\lambda_i(y)d(b,z_i(y))^2\). There are finitely many active
terms. Their total weight is one; boundedness of their point set gives
coercivity, so a minimum exists by properness. The midpoint inequality
gives uniqueness. At a fixed point shrink a neighborhood to exclude the
finitely many nearby closed supports not containing that point. Each
remaining support is inside its map domain, so the corresponding
finitely many point maps are defined and bounded on a further
neighborhood. A fixed comparison point bounds the minimum values. The
inequality \(d(b,z_i)\ge d(b,o)-d(o,z_i)\), summed with total weight
one, places all minimizers in one compact ball. Any convergent subnet of
minimizers minimizes the limiting finite sum, by continuity of its
weights and active point maps; zero-weight terms tend to zero on this
bounded ball. Uniqueness forces the limit, proving continuity.
Invariance of the weights and uniqueness of the minimizer prove
equivariance. Two equivariant maps are homotopic by the continuous
unique-segment operation. \(\square\)

AB.0--AB.8 prove the full split local-field block, including the
enlarged central factor required for properness. AB.6--AB.7 also provide
an abstract geometric theorem usable over an infinite-residue valued
field once its exact affine chamber data have separately been proved.
They do not supply nonsplit root filtrations or connected facet models
by assertion.

\subsubsection{AB.9. Unramified integral
descent}\label{ab.9.-unramified-integral-descent}

Here k is a nonarchimedean local field, R its complete DVR, and f=
\(F_{q}\) its residue field. AB.0 proves the finite-extension
arithmetic. The regular-domain and local Jacobian criteria used below
are proved in the preceding lessons named in BF.0:
\href{../AG-CA/AG-CA-12.html}{Regular sequences, depth and
Cohen--Macaulay modules}, Theorem 6.1, and
\href{../AG-CA/AG-CA-18.html}{Smooth algebras over a field and the
Jacobian criterion}, Lemma 1.1 and Theorem 2.1. All lifting, openness
and descent arguments used here are supplied below. References are
Conrad, \emph{Lang's theorem and unirationality}, Section 2, and
Bruhat--Tits II, Sections 5.1.5--5.1.9.

\paragraph{AB.9.1. FINITE UNRAMIFIED EXTENSIONS AND INTEGRAL
TRACE}\label{ab.9.1.-finite-unramified-extensions-and-integral-trace}

Let E/k be finite unramified (a finite separable extension with
ramification index one) and B= \(R_{E}\) . Write m={[}E:k{]}. AB.0 gives
that \(\pi\) is also a uniformizer of B, B is complete and finite free,
and B/ \(\pi\) B= \(F_{q^m}\) . Let bar( \(\alpha\) ) generate
\(F_{q^m}\) over f, and let bar(F) be its monic minimal polynomial, of
degree m. Such a generator exists: the finite multiplicative group of a
field is cyclic. To check the latter claim, in each primary part of a
finite abelian subgroup choose an element of maximal order and multiply
these elements. Its order is the exponent h of the group. Every element
is a root of \(X^h\) -1, so the root bound gives the group order at most
h; equality follows. A generator of the multiplicative group generates
the field as well.

Lift bar(F) to a monic F in R{[}T{]}. It is irreducible over k. Indeed
if F=PQ with monic P,Q in k{[}T{]}, multiply P and Q by the smallest
nonnegative powers \(\pi^r\) , \(\pi^s\) making their coefficients
integral and primitive. Their reductions have nonzero product over f, so
their product is primitive. The content of \(\pi^{r+s}\) F is
\(\pi^{r+s}\) ; consequently r+s=0. The factors were already integral
and monic, and their reductions would contradict irreducibility of
bar(F).

The residue polynomial is separable because finite fields are perfect:
the q-power map is bijective, and an irreducible polynomial with zero
derivative would be a pth power. Choose an arbitrary lift \(a_{0}\) of
bar( \(\alpha\) ) in B. Then F( \(a_{0}\) ) is in \(\pi\) B and F'(
\(a_{0}\) ) is a unit. Newton's iteration a -\textgreater{} a-F(a)/F'(a)
at least doubles the \(\pi\) -order of the error and keeps the
derivative a unit. Completeness gives a root \(\alpha\) in B with the
specified reduction. A second root with that reduction would have
difference times a unit equal to zero, by polynomial division and the
derivative test, so the lift is unique.

Since deg(F)=m={[}E:k{]}, E=k( \(\alpha\) ). In fact B=R{[} \(\alpha\)
{]}. For x in B, write its residue as a polynomial of degree less than m
in bar( \(\alpha\) ), subtract the corresponding element of R{[}
\(\alpha\) {]}, and divide by \(\pi\) . Repeating gives x as a \(\pi\)
-adic sum of elements in the free R-module with basis 1, \(\alpha\)
,\ldots, \(\alpha^{m-1}\) . That module is complete and is contained in
B; the partial sums converge to x in B. Thus B = R{[}T{]}/(F), F'(
\(\alpha\) ) in \(B^*\) . This finite free algebra is etale: in a
square-zero lifting problem the equation F(a)=0 has a unique correction
because F'(a) is invertible. Iterating square-zero steps handles
nilpotent ideals; this is formal etaleness, and the displayed finite
presentation is the etale criterion.

Each of the m residue conjugates bar( \(\alpha\) )\^{}\{ \(q^i\) \} has
a unique root lift \(\alpha_{i}\) of F in B. Hence E is the splitting
field of a separable polynomial of degree m and E/k is Galois. Reduction
identifies its group \(\Gamma\) with Gal( \(F_{q^m}\) / \(F_{q}\) ),
generated by q-Frobenius. These m Frobenius powers are distinct: if a
power \(q^i\) with 0\textless i\textless m fixed the entire field,
\(X^{q^i}\) -X would have \(q^m\) roots. Thus the group is cyclic of
order m. The differences \(\alpha_{i}\) - \(\alpha_{j}\) are units, so
the Chinese remainder theorem proves the Galois-cover identity B
\(\otimes_{R}\) B -\textgreater{} \(\prod_{\gamma\in \Gamma}\) B, x
tensor y \textbar-\textgreater{} (x \(\gamma\) (y))\_ \(\gamma\) , is an
isomorphism. B is faithfully flat over R because it is free of positive
rank. This proves the complete finite etale Galois-cover claim.

The trace reduces to the residue trace. One can see this from the Galois
sum Tr(x)= \(\sum_{\gamma} \gamma\) (x), or from a multiplication matrix
in the displayed integral basis. The residue trace is onto f: its
polynomial is T(X)=X+ \(X^q\) +\ldots+X\^{}\{ \(q^{m-1}\) \}. It is a
nonzero polynomial of degree less than \(q^m\) and therefore cannot
vanish on all of \(F_{q^m}\) . It is f-linear and has values in f, so
its nonzero image is all of the one-dimensional f-space f.~Choose bar(c)
of trace one and any lift \(c_{0}\) in B. The element t=Tr( \(c_{0}\) )
lies in R and is congruent to 1 modulo \(\pi\) , so t is a unit. Then c=
\(c_{0}\) /t satisfies Tr(c)=1. This never divides by m. In particular
it remains valid when p divides m.

\paragraph{AB.9.2. ADDITIVE SEMILINEAR
COCYCLES}\label{ab.9.2.-additive-semilinear-cocycles}

Let M be any B-module, with \(\gamma\) (bu)= \(\gamma\) (b) \(\gamma\)
(u). No finiteness, flatness or freeness of M is required. For
\(z_{\sigma \tau}\) = \(z_{\sigma}\) + \(\sigma\) ( \(z_{\tau}\) ), put
v=- \(\sum_{\gamma\in \Gamma} \gamma\) (c) \(z_{\gamma}\) , where c is
from AB.9.1. Reindex \(\delta\) = \(\sigma \gamma\) and use \(\sigma\) (
\(z_{\gamma}\) )= \(z_{\sigma \gamma}\) - \(z_{\sigma}\) . Directly,
\(\sigma\) (v) =- \(\sum_{\delta} \delta\) (c)( \(z_{\delta}\) -
\(z_{\sigma}\) ) =v+( \(\sum_{\delta} \delta\) (c)) \(z_{\sigma}\) =v+
\(z_{\sigma}\) . Thus \(z_{\sigma}\) = \(\sigma\) (v)-v with exactly the
stated sign. Applying the same calculation to \(F_{q^m}\) / \(F_{q}\)
and a residue trace-one element proves the residue-module assertion,
including degrees divisible by p.

\paragraph{AB.9.3. SMOOTH INTEGRAL POINTS AND CONGRUENCE
QUOTIENTS}\label{ab.9.3.-smooth-integral-points-and-congruence-quotients}

Let H be a smooth affine group scheme of finite presentation over R. Fix
a point modulo \(\pi^j\) , j\textgreater=1. Its reduction modulo \(\pi\)
lies in a standard smooth chart. Choose arbitrary lifts of the chart
coordinates modulo \(\pi^{j+1}\) ; the equations have errors in
\(\pi^j\) / \(\pi^{j+1}\) . Correct the c coordinates indexed by an
invertible Jacobian minor by solving the c linear equations for the
negative error, leaving the other coordinates fixed. Quadratic
correction terms vanish because ( \(\pi^j\) / \(\pi^{j+1}\) )\^{}2=0.
The chart's localization denominators remain units because they were
units at the residue point. This constructs a point modulo \(\pi^{j+1}\)
lifting the given point. Starting from any specified point modulo
\(\pi^n\) and repeating constructs compatible points at every higher
level.

If H=Spec(A), compatible maps A -\textgreater{} R/ \(\pi^j\) define a
map A -\textgreater{} \(\lim_{j}\) R/ \(\pi^j\) =R by taking the value
on each a in A. Addition, multiplication and constants are respected at
every level, hence in R; uniqueness follows from \(\bigcap_{j} \pi^j\)
R=0. Consequently H(R)= \(\lim_{j}\) H(R/ \(\pi^j\) ), and the
constructed compatible lifts show that every reduction map is onto. This
is an actual lifting proof, not an assumption that smooth points lift
over an arbitrary residue field.

Set \(K_{n}\) =ker(H(R)-\textgreater H(R/ \(\pi^n\) )), n\textgreater=1,
and I= \(\pi^n\) R/ \(\pi^{n+1}\) R. A point of H(R/ \(\pi^{n+1}\) )
reducing to the identity modulo \(\pi^n\) is a map \(\epsilon\) +d,
where \(\epsilon\) is the identity's evaluation and d:A -\textgreater{}
I satisfies d(ab)=bar( \(\epsilon\) (a))d(b)+bar( \(\epsilon\) (b))d(a),
d(R)=0. Conversely every such derivation gives a point: \(I^2\) =0 makes
the displayed equation exactly the multiplicativity requirement. Since
\(\pi\) I=0, this is the tangent space at the identity of \(H_{f}\) with
coefficients I, canonically Lie( \(H_{f}\) ) \(\otimes_{f}\) I. The
comultiplication giving the group law has linear term \(d_{1}\) +
\(d_{2}\) ; all products of their values vanish in \(I^2\) . Thus this
kernel is the additive group of that vector space. The surjectivity just
proved identifies it with \(K_{n}\) / \(K_{n+1}\) . The derivation
description uses only identity evaluation and ring maps, so it commutes
with base change and ring automorphisms. In particular for B/R it is the
semilinear \(F_{q^m}\) -module Lie( \(H_{f}\) ) \(\otimes_{f}\) (
\(\pi^n\) B/ \(\pi^{n+1}\) B).

\paragraph{AB.9.4. THE FINITE-FIELD COCYCLE STEP, WITH ITS LANG
PROOF}\label{ab.9.4.-the-finite-field-cocycle-step-with-its-lang-proof}

A smooth connected affine group J over a finite field is geometrically
connected. Here is the needed descent check. Over the algebraic closure
its identity component is open and closed: smooth local rings are
regular domains, so different irreducible components cannot meet; there
are only finitely many components. The component containing the rational
identity is invariant under every field automorphism. Its defining
idempotent in the affine coordinate ring descends, since the
coefficients of an invariant element of A \(\otimes_{f}\) bar(f) belong
to f (expand in a finite linearly independent set of elements of A). A
proper identity component would therefore disconnect J over
f.~Connectedness excludes this.

Write F for q-Frobenius on \(J_{bar}\) (f). It has zero differential.
For any x define the right twisted action \(T_{g}\) (x)= \(g^{-1}\)
xF(g). Its orbit map \(\psi_{x}\) sends the identity to x; after left
translation by \(x^{-1}\) its differential there is -Ad( \(x^{-1}\) ),
hence invertible. The identity \(\psi_{x}\) (g h)=
\(\psi_{\psi_{x}(g)}\) (h) shows that its differential is invertible
everywhere. The Jacobian criterion for morphisms between smooth
varieties of equal dimension makes these orbit maps etale, hence open.
On closed points their images are the disjoint orbits and cover
\(J_{bar}\) (f). Distinct open images cannot intersect, since a nonempty
intersection has a closed point; their union is all, since its closed
complement would otherwise have a closed point. Geometric connectedness
forces a single orbit. In particular g \textbar-\textgreater{}
\(g^{-1}\) F(g) is onto on bar(f)-points. This proves precisely the Lang
assertion used here.

For clarity, openness in this argument is the elementary etale-chart
fact. A standard etale chart is a localization of C{[}t{]}/(P), with P
monic and P' invertible there. The algebra before localization is finite
free. For a principal open D(u) in such an algebra, a fibre meets D(u)
exactly when multiplication by u on that finite-dimensional fibre is not
nilpotent. Equivalently at least one nonleading coefficient of its
characteristic polynomial is nonzero, by Cayley--Hamilton. Its image is
therefore the union of the principal opens defined by these
coefficients. Thus these charts, and consequently etale morphisms, are
open. The tangent criterion is the same local Jacobian criterion: in
smooth local coordinates an invertible differential gives the
dimension-zero invertible-Jacobian presentation.

Now let a be a multiplicative \(\Gamma\) -cocycle with values J(
\(F_{q^m}\) ). For the Frobenius generator \(\sigma\) , put A=
\(a_{\sigma}\) . The cocycle identity gives A F(A)\ldots{} \(F^{m-1}\)
(A)=1. Lang supplies bar(b) in J(bar(f)) with
A=bar(b)\^{}\{-1\}F(bar(b)). Telescoping the preceding product gives
bar(b)\^{}\{-1\} \(F^m\) (bar(b))=1. Therefore bar(b) is actually in J(
\(F_{q^m}\) ). The values on all powers of \(\sigma\) are determined by
A, so \(a_{\gamma}\) =bar(b)\^{}\{-1\} \(\gamma\) (bar(b)) for every
\(\gamma\) . No passage to a larger residue field is needed at the end,
and no restriction on m was used.

\paragraph{\texorpdfstring{AB.9.5. UNRAMIFIED INTEGRAL \(H^1\) ,
INCLUDING p-DIVISIBLE
DEGREES}{AB.9.5. UNRAMIFIED INTEGRAL H\^{}1 , INCLUDING p-DIVISIBLE DEGREES}}\label{ab.9.5.-unramified-integral-h1-including-p-divisible-degrees}

Let H/R be smooth affine of finite presentation with connected special
fibre. Let \(a_{\sigma \tau}\) = \(a_{\sigma} \sigma\) ( \(a_{\tau}\) )
take values in H(B). Its residue cocycle is trivial by AB.9.4. By AB.9.3
lift a residue solution to \(b_{1}\) in H(B), so \(r^{(1)}_{\sigma}\) =
\(b_{1} a_{\sigma} \sigma\) ( \(b_{1}\) )\^{}\{-1\} is in \(K_{1}\) .
These r-values still form a cocycle for the original, untwisted
\(\Gamma\) action, as direct multiplication shows.

Inductively suppose \(b_{n}\) has \(r^{(n)}_{\sigma}\) in \(K_{n}\) . In
the quotient \(K_{n}\) / \(K_{n+1}\) , AB.9.3 turns the cocycle into an
additive semilinear cocycle \(z_{\sigma}\) . By AB.9.2 choose \(v_{n}\)
with \(z_{\sigma}\) = \(\sigma\) ( \(v_{n}\) )- \(v_{n}\) . Lift
\(v_{n}\) to \(h_{n}\) in \(K_{n}\) and set \(b_{n+1}\) =
\(h_{n} b_{n}\) . Then \(r^{(n+1)}_{\sigma}\) =
\(h_{n} r^{(n)}_{\sigma} \sigma\) ( \(h_{n}\) )\^{}\{-1\}, whose
additive class is \(v_{n}\) + \(z_{\sigma}\) - \(\sigma\) ( \(v_{n}\)
)=0. Thus its values lie in \(K_{n+1}\) . Also \(b_{n+1}\) and \(b_{n}\)
have the same image modulo \(\pi^n\) . The compatible images define b in
H(B), by the inverse-limit equality of AB.9.3. For each \(\sigma\) and
every n, b \(a_{\sigma} \sigma\) (b)\^{}\{-1\}=1 mod \(\pi^n\) .
Separatedness of B makes the equality exact. Consequently \(a_{\sigma}\)
= \(b^{-1} \sigma\) (b), and \(H^1\) ( \(\Gamma\) ,H(B)) consists of one
element. The trace-one computation, rather than division by \textbar{}
\(\Gamma\) \textbar, is the entire congruence step. Hence the proof
includes every finite unramified degree, also when p divides it.

There is also an infinite-unramified version for discrete continuous
cohomology. Let \(R_{\mathrm{ur}}\) be the union of all finite
unramified integer rings and \(\Gamma_{\mathrm{ur}}\) =Gal( \(k^ur\)
/k). Every point of H( \(R_{\mathrm{ur}}\) ) is defined over a finite
subextension because H is finitely presented. A continuous cocycle from
the compact \(\Gamma_{\mathrm{ur}}\) into this discrete group has finite
image. Its value is 1 on some open normal subgroup, by continuity at the
identity. Choose one finite unramified E containing the definitions of
all image points and with \(\Gamma_{E}\) inside that subgroup. The
cocycle factors through \(\Gamma_{\mathrm{ur}}\) / \(\Gamma_{E}\) ,
takes values in H( \(R_{E}\) ), and the finite result supplies an
integral solution. This also proves the asserted continuous \(H^1\)
vanishing; compactness here is used only for finite image.

\paragraph{AB.9.6. FIXED COSETS}\label{ab.9.6.-fixed-cosets}

Suppose H has identified generic fibre G. The affine coordinate ring of
H is torsion-free and embeds into its localization k{[}G{]}, so its
integral points embed into G(k). Put \(P_{E}\) =H(B), P=H(R). If
\(xP_{E}\) is \(\Gamma\) -fixed, then \(a_{\sigma}\) = \(x^{-1} \sigma\)
(x) is in \(P_{E}\) and is a cocycle. AB.9.5 gives \(a_{\sigma}\) =
\(b^{-1} \sigma\) (b) with b in \(P_{E}\) . The representative x
\(b^{-1}\) is fixed by \(\Gamma\) , since \(\sigma\) (x \(b^{-1}\) )=x
\(a_{\sigma} \sigma\) (b)\^{}\{-1\}=x \(b^{-1}\) . It therefore belongs
to G(k); this last assertion follows by applying \(\Gamma\) -invariance
to the finitely many affine coordinates of G. Thus G(k)/P
-\textgreater{} (G(E)/ \(P_{E}\) )\^{} \(\Gamma\) is onto. If g,h in
G(k) have equal images, every H-coordinate of \(g^{-1}\) h lies in B
intersection k=R, so \(g^{-1}\) h lies in H(R). This proves injectivity.
The construction commutes with left G(k)-translation.

The quotient topologies are discrete. Choose finitely many R-algebra
generators of the affine ring of H. Inside G(k), integrality of each of
their values is an open condition, because R is open in k; hence P is
open. H(R) is also a closed subset of \(R^d\) cut out by its equations,
so it is compact. The same applies to \(P_{E}\) . A bijection of the
indicated discrete quotients is a homeomorphism.

For \(k^ur\) choose a finite unramified Galois E containing the finitely
many coordinates of x in G( \(k^ur\) ). If its coset is fixed, then for
\(\sigma\) in Gal(E/k) the element \(x^{-1} \sigma\) (x) lies in G(E)
intersection H( \(R_{\mathrm{ur}}\) )=H( \(R_{E}\) ): this equality
follows coordinate by coordinate from \(R_{\mathrm{ur}}\) intersection
E= \(R_{E}\) and AB.0's extension valuation. Apply the finite argument.
Injection is the same intersection argument over k. The target is given
its discrete coset topology (also obtained from the valuation topology,
since the integer ring of \(k^ur\) is open).

\paragraph{AB.9.7. AFFINE SEMILINEAR DESCENT WITHOUT INVOKING A DESCENT
BLACK
BOX}\label{ab.9.7.-affine-semilinear-descent-without-invoking-a-descent-black-box}

Let D be a B-algebra with semilinear \(\Gamma\) -action. Set A=
\(D^\Gamma\) . First prove module descent explicitly, since it will also
descend all the group maps. Choose the basis \(x_{i}\) =
\(\alpha^{i-1}\) of B and its trace-dual basis \(y_{i}\) . The latter is
integral: the residue trace pairing is nondegenerate (if u is nonzero,
pair it with bar(c)/u), so the integral trace Gram determinant is a
unit. These bases satisfy \(\sum_{i} x_{i} \gamma\) ( \(y_{i}\) )=
\(\delta_{\gamma,1}\) . For example let U=( \(\gamma\) ( \(x_{i}\) ))
and V=( \(\gamma\) ( \(y_{i}\) )). The unit Vandermonde determinant
makes U invertible over B, and trace duality says \(U^t\) V=I. Thus U
\(V^t\) =I, whose identity-row entries are the formula.

For any semilinear B-module M put \(M_{0}\) = \(M^\Gamma\) . The natural
map \(\Phi\) :B \(\otimes_{R} M_{0}\) -\textgreater{} M, b tensor u
\textbar-\textgreater{} bu is an isomorphism. Its explicit inverse as an
R-linear map is Psi(m)= \(\sum_{i} x_{i}\) tensor (
\(\sum_{\gamma} \gamma\) ( \(y_{i}\) m)). The second factors are
invariant. The displayed \(\delta\) identity gives \(\Phi\) Psi(m)=m.
For an invariant u, trace duality gives Psi \(\Phi\) (b tensor u)=
\(\sum_{i} x_{i}\) tensor Tr( \(y_{i}\) b)u=b tensor u. Both identities
hold for arbitrary modules, without assuming they are finite, free or
torsion-free. Since \(\Phi\) is B-linear, its inverse is B-linear too.
Applied to D this proves B \(\otimes_{R}\) A = D as algebras.

One further useful identity removes any exactness ambiguity. For a
semilinear module the operator \(P_{c}\) (m)= \(\sum_{\gamma} \gamma\)
(c m) is an R-linear projection onto its invariants: on an invariant
element it is multiplication by Tr(c)=1. If M=B \(\otimes_{R}\) N with
action only on B, its image lies exactly in 1 tensor N, because
\(P_{c}\) ( \(\sum_{j} b_{j}\) tensor \(n_{j}\) )=1 tensor \(\sum_{j}\)
Tr(c \(b_{j}\) ) \(n_{j}\) . Hence (B \(\otimes_{R}\) N)\^{} \(\Gamma\)
=N for every R-module N.

Assume now D is of finite presentation over B. Take finitely many
B-algebra generators \(d_{j}\) and write each as
\(\sum_{i} x_{i} a_{ij}\) , with \(a_{ij}\) in A, using Psi. Let
\(A_{0}\) be the R-algebra generated by these finitely many \(a_{ij}\) .
Then B tensor \(A_{0}\) -\textgreater{} D is onto, and its map into B
tensor A is injective because B is free. Faithful freeness annihilates
no nonzero quotient, so A/ \(A_{0}\) =0. Thus A is of finite type. Since
R is noetherian, the polynomial presentation of A has finitely generated
kernel, making A finitely presented. Here the elementary Hilbert basis
argument is enough: leading coefficients of ideals in one polynomial
variable generate a finite ideal in the coefficient ring; subtract
finitely many chosen leading terms to reduce degrees, then use induction
in the bounded-degree coefficient module. Repeat for the finitely many
variables.

If D is smooth over B, it is B-flat and hence R-flat. A is an R-direct
summand by \(P_{c}\) , so A is R-flat. Equivalently in this DVR
situation it is torsion-free, and a torsion-free R-module is a filtered
union of its finite free submodules; tensoring shows it is flat. The
generic and special fibres of Spec(A) become the smooth fibres of
Spec(D) after the respective field extensions E/k and \(f_{E}\) /f.~Over
an algebraic closure they are the same schemes, so both are
geometrically smooth.

Here is the DVR fibre-to-total-space step explicitly. Present A=R{[}
\(t_{1}\) ,\ldots, \(t_{d}\) {]}/I. At a point of the special fibre
choose equations bar( \(F_{1}\) ),\ldots,bar( \(F_{c}\) ) generating the
localized fibre ideal with an invertible Jacobian minor; the fibre is
smooth. Flatness gives I intersection \(\pi\) R{[}t{]}= \(\pi\) I, so
these equations lift from I/ \(\pi\) I to equations \(F_{i}\) in I. In
the local polynomial ring at this point, the finite module I/( \(F_{1}\)
,\ldots, \(F_{c}\) ) is its own \(\pi\) -multiple. Nakayama's
determinant argument forces it to vanish: on finite generators the
matrix 1- \(\pi\) C has unit determinant and annihilates the module.
After shrinking, I is generated by these \(F_{i}\) and their Jacobian
minor is invertible. This is a standard smooth chart over R. At a
generic-fibre point the same chart argument takes place after inverting
\(\pi\) . These are all points of Spec(A), proving smoothness.

Finally let Spec(D) be a B-group scheme and require the semilinear
action to respect its group law, identity and inverse. Under D=B tensor
A one has D \(\otimes_{B}\) D = B \(\otimes_{R}\) (A \(\otimes_{R}\) A).
The projection identity above identifies its invariants with A
\(\otimes_{R}\) A. Equivariance therefore restricts comultiplication,
counit and antipode to maps A -\textgreater{} A tensor A, A
-\textgreater{} R, A -\textgreater{} A. Their Hopf identities hold
because they hold after the faithfully free extension to B. Hence
Spec(A) is a smooth affine finitely presented R-group scheme descending
the given scheme and action. The construction also proves uniqueness:
invariant maps recover exactly the maps before base change.

If the original special fibre is connected, so is the descended one: a
nontrivial idempotent in A/ \(\pi\) A would remain nontrivial after the
faithful field extension \(f_{E}\) /f, contradicting connectedness of D/
\(\pi\) D. In the smooth-group case it is geometrically connected by the
identity component argument of AB.9.4. Since \(\pi\) is \(\Gamma\)
-fixed and D is torsion-free, invariants commute with inverting \(\pi\)
. Thus if the generic semilinear descent is G, the generic fibre of
Spec(A) is G; the scheme is an integral model of G.

\subsubsection{AB.10. The valuation lattice of an arbitrary local-field
torus}\label{ab.10.-the-valuation-lattice-of-an-arbitrary-local-field-torus}

\textbf{Proposition.} Let \(T\) be any \(k\)-torus and \(L/k\) a finite
Galois splitting field. Put \(\Lambda=X_*(T_L)\),
\(\Gamma=\operatorname{Gal}(L/k)\), and \(e=e(L/k)\). The valuation
vector

\[
 \langle\chi,\nu(t)\rangle=\omega_L(\chi(t)),\qquad t\in T(k),
 \tag{AB.21}
\]

lies in \(\Lambda^\Gamma\otimes\mathbb R\), and

\[
 \Lambda^\Gamma\subset\nu(T(k))\subset e^{-1}\Lambda^\Gamma.
 \tag{AB.22}
\]

Its kernel \(T^b\) is compact open. The invariant lattice is the
cocharacter lattice of the maximal split subtorus \(S\subset T\). Thus
the image is a full discrete lattice in \(V(S)\), including all split
central directions. These assertions impose no condition on the
ramification or residue characteristic.

\textbf{Proof.} AB.0 provides the unique extending valuation and its
value group. The characters are a free lattice; multiplicativity of
their evaluations makes (AB.21) a homomorphism from that lattice to
\(\mathbb R\), hence a vector in \(\Lambda\otimes\mathbb R\). Choose a
character basis. Its coordinates have valuations in \(e^{-1}\mathbb Z\),
so the vector is in \(e^{-1}\Lambda\). Galois conjugation of the
character evaluation does not change its valuation, by AB.0. Since \(t\)
is rational, the vector is fixed by \(\Gamma\), giving the second
containment in (AB.22).

For clarity the invariant-lattice assertion does not require a rational
splitting of the torus. A cocharacter of \(T_L\) is a monomial
homomorphism from \(\mathbf G_{m,L}\); its coordinate formulas descend
to \(k\) exactly when that homomorphism is Galois invariant. One may
check descent coefficient by coefficient in the affine coordinate rings:
invariants of \(A\otimes_k L\) are \(A\), by expanding in a finite
linearly independent set of elements of \(A\), and the invariance of a
morphism makes its ring map take values there. The invariant lattice is
saturated, because if \(n\lambda\) is invariant then
\(n(\gamma\lambda-\lambda)=0\) in the torsion-free lattice. Choose a
basis for this saturated sublattice and extend it to a basis of
\(\Lambda\): Euclidean row and column operations reduce a lattice
inclusion to Smith diagonal form, and saturation forces every nonzero
diagonal factor to be one. The associated monomial map from a split
torus to \(T_L\) is a closed immersion, its coordinate map being
surjective on characters. It descends by the preceding coefficient
argument and gives a split subtorus \(S\) with that cocharacter lattice.
Every split subtorus has invariant cocharacters and is contained in it,
so it is maximal. For each invariant cocharacter \(\lambda\), its value
\(\lambda(\pi)\in T(k)\) has valuation vector \(\lambda\), proving the
first containment in (AB.22). The middle group is consequently a
finite-index over-lattice of \(\Lambda^\Gamma\), and is discrete and
full.

In splitting coordinates, \(T^b\) is exactly
\(T(k)\cap(O_L^*)^{\operatorname{rank}T}\). The coordinate isomorphism
\(T(L)\simeq(L^*)^r\) is a homeomorphism, with inverse given by Laurent
monomials. The subset \(T(k)\) is closed: it is the fixed subset of the
finitely many continuous Galois automorphisms on \(T(L)\); its topology
is the original \(k\)-point topology, since affine coordinates and their
inverse formulas give the finite-dimensional subspace topology of AB.0.
Thus this intersection is compact and open in \(T(k)\). In particular an
anisotropic local-field torus is compact. \(\square\)

The proposition supplies the torus lattice and compact kernel. Relative
roots and reflection representatives are proved in GB.1--GB.5. For the
original nonsplit group, the centralizer lattice and the full rational
normalizer action are proved in GB.19 and GB.22--GB.23.

\subsubsection{AB.11. Passage to arbitrary reductive
groups}\label{ab.11.-passage-to-arbitrary-reductive-groups}

The two-cell identity and the quasi-split root-coordinate construction
are proved in GB.0--GB.6. Smooth root models and connected facet models
are constructed in GB.7--GB.10. GB.11--GB.13 give the exact
component-cocycle test, torus lifting and stable apartments. The
arithmetic of the unramified field is GB.9 and GB.14; its torus
vanishing, regular-class representative and quasi-splitting are proved
in GB.16--GB.18. For descent to the original field, GB.19 proves the
anisotropic centralizer and central lattice, GB.20 the rational
filtration and half-apartment bounds, GB.21 the residue/coface
correspondence, GB.22 the maximal rational apartment and wall-exchange
assertions, and GB.23 the full stabilizers, topology, metric and
properness. These statements provide the hypotheses for the affine
geometry above, while retaining the full Euclidean centre and the finite
component obstruction.

\subsection{Relative roots and arbitrary-group
descent}\label{relative-roots-and-arbitrary-group-descent}

\subsubsection{GB.0. The complete filtered rank-one two-cell
identity}\label{gb.0.-the-complete-filtered-rank-one-two-cell-identity}

Use the valuation extended from the original local field, without
renormalizing intermediate fields. Let \(E/k\) be finite separable. In
the first case take \(H=\operatorname{SL}_2(E)\), with its usual upper
and lower root coordinates and \(\varphi_\pm(x_\pm(u))=\omega(u)\). In
the second case let \(F/E\) be separable quadratic, take
\(H=\operatorname{SU}_3(F/E)\), and use exactly BF.7's matrices

\[
 x_+(u,v)=\begin{pmatrix}1&-\bar u&v\\0&1&u\\0&0&1\end{pmatrix},\quad
 x_-(u,v)=\begin{pmatrix}1&0&0\\u&1&0\\v&-\bar u&1\end{pmatrix},\quad
 v+\bar v=-u\bar u.
 \tag{GB.1}
\]

In this second case \(\varphi_\pm(x_\pm(u,v))=\omega(v)/2\); the
doubled-root group is already included in each short-root group. The
diagonal bounded torus \(T_0^b\) consists respectively of
\(\operatorname{diag}(t,t^{-1})\) with \(\omega(t)=0\), and
\(D(t)=\operatorname{diag}(t,\bar t/t,1/\bar t)\) with \(\omega(t)=0\).
Let \(r\) be an attained nonidentity value, and put

\[
 A=U_{+,r},\quad A^+=U_{+,r+},\quad
 B=U_{-,-r+},\quad D=U_{-,-r}.
 \tag{GB.2}
\]

Choose \(u_0\in A\setminus A^+\) and its BF rank-one normalizer
completion \(s=m(u_0)\). In SL \(_2\) this is \(n(c_0)\), with
\(\omega(c_0)=r\); in SU \(_3\) it is \(M(v_0)\), with
\(\omega(v_0)=2r\), where

\[
 M(v)=\begin{pmatrix}0&0&v\\0&-\bar v/v&0\\1/\bar v&0&0\end{pmatrix}.
 \tag{GB.3}
\]

\textbf{Theorem.} In both cases,

\[
 K_r=\langle T_0^b,A,D\rangle
    =A sT_0^b A\ \cup\ B T_0^b A
 \tag{GB.4}
\]

is a compact subgroup. Its normalizer element acts as reflection in the
wall \(a(h)+r=0\) and has square in \(T_0^b\). The statement includes
every attained short-root level of every ramified SU \(_3\) chart in
every residue characteristic.

\textbf{Proof.} For real weights \(w_i\), set
\(\|z\|_w=\max_i\exp(-\omega(z_i)-w_i)\). A matrix and its inverse
preserve this norm exactly when both have the entry inequalities
\(\omega(g_{ij})\ge w_j-w_i\): necessity follows by testing basis
vectors, sufficiency by the nonarchimedean triangle inequality. In SL
\(_2\) use weights \((-r/2,r/2)\); in SU \(_3\) use \((-r,0,r)\). The
unitary inverse is \(J\bar g^{\mathsf T}J\), where \(J\) is the
antidiagonal form matrix of BF.7, so the same inequalities for \(g\)
imply those for its inverse. The SL \(_2\) inverse formula gives the
identical conclusion. Thus the set of matrices in \(H\) satisfying these
inequalities is a norm-isometry subgroup, denoted \(K\). It is compact:
each entry is in a fixed compact valuation ball by AB.0, and the
determinant and unitary equations cut out a closed subset of their
finite product. The inequalities are

\[
 \begin{array}{c|ccc}
 \operatorname{SU}_3&\text{diagonal}&(1,2),(2,3)&(1,3)\\
 \text{bound}&0&r&2r
 \end{array}
 \quad
 \begin{array}{c|cc}
 &(2,1),(3,2)&(3,1)\\
 \text{bound}&-r&-2r.
 \end{array}
 \tag{GB.5}
\]

They place \(T_0^b,A,D\) in \(K\); the trace inequality
\(\omega(u)\ge\omega(v)/2\) is BF.7's locally proved identity. The
opposite factors completing \(u_0\) have value \(-r\) by BF.7, or the
pinned SL \(_2\) formula, so \(s\) is in the generated subgroup. We now
prove that every element of \(K\) has the displayed two-cell expression.
This also proves subgroup closure of (GB.4), rather than assuming it
from the union formula.

First take SL \(_2\) and write
\(g=\begin{pmatrix}a&b\\c&d\end{pmatrix}\). If \(\omega(c)>-r\), then
\(\omega(bc)>0\); the equation \(ad-bc=1\) forces \(a,d\) to be units.
Direct Gaussian multiplication gives

\[
 g=x_-(c/a)\operatorname{diag}(a,a^{-1})x_+(b/a)\in B T_0^b A.
 \tag{GB.6}
\]

If \(\omega(c)=-r\), left multiplication by \(x_+(-a/c)\in A\) gives
\(\begin{pmatrix}0&-1/c\\c&d\end{pmatrix}=n(-1/c)x_+(d/c)\). The last
parameter has valuation at least \(r\), and \(n(-1/c)s^{-1}\in T_0^b\),
since both normalizer parameters have valuation \(r\). Thus
\(g\in A sT_0^b A\). These are all cases.

Now take SU \(_3\) and write its first column as
\((a,b,c)^{\mathsf T}\). Its unitary isotropy equation is

\[
 a\bar c+c\bar a+b\bar b=0.
 \tag{GB.7}
\]

If \(\omega(c)>-2r\), then \(a\) must be a unit. Indeed, if
\(\omega(a)>0\), (GB.7) forces \(\omega(b)>-r\). All three first-column
entries would then have strictly greater valuation than their entry
bounds. The weighted norm of \(ge_1\) would be strictly smaller than
that of \(e_1\), contradicting norm preservation. With \(a\) a unit,
(GB.7) again gives \(\omega(b)>-r\). The lower pair \((p,q)=(b/a,c/a)\)
satisfies \(q+\bar q=-p\bar p\) and has value greater than \(-r\).
Multiplication by \(x_-(p,q)^{-1}\) kills the two lower first-column
entries. The unitary equations then force the result to be upper
triangular, with diagonal \(D(a)\), and after removing that diagonal its
upper coordinates are \(u=-\overline{g_{12}/a}\), \(v=g_{13}/a\). Both
follow from the unchanged first row of that multiplication. The entry
bounds imply \(\omega(u)\ge r\), \(\omega(v)\ge2r\). Hence
\(g\in B T_0^b A\).

In the remaining case \(\omega(c)=-2r\), put

\[
 u=-b/c,\qquad v=(\bar u b-a)/c.
 \tag{GB.8}
\]

The inequalities give \(\omega(u)\ge r\), \(\omega(v)\ge2r\). Equation
(GB.7), divided by \(c\bar c\), gives
\(a/c+\bar a/\bar c=-b\bar b/(c\bar c)\). Consequently

\[
 v+\bar v
 =-2b\bar b/(c\bar c)-(a/c+\bar a/\bar c)
 =-b\bar b/(c\bar c)=-u\bar u.
 \tag{GB.9}
\]

This equality is an identity in the field and never divides by two. Thus
\(x_+(u,v)\in A\). It sends the first column to \((0,0,c)^{\mathsf T}\).
The unitary pairings with this column and the determinant-one condition
force the resulting matrix to have the form

\[
 g'=\begin{pmatrix}0&0&1/\bar c\\0&-\bar c/c&e\\c&f&j\end{pmatrix}
    =M(1/\bar c)x_+(u',v'),
 \tag{GB.10}
\]

where \(u'=e/(-\bar c/c)\), \(v'=j/c\); its remaining unitary equations
give \(f=-c\bar u'\) and \(v'+\bar v'=-u'\bar u'\). Since \(x_+(u,v)\)
is a norm isometry, \(g'\) still satisfies (GB.5), so
\(\omega(u')\ge r\), \(\omega(v')\ge2r\). Also
\(\omega(1/\bar c)=2r=\omega(v_0)\), so
\(M(1/\bar c)s^{-1}=D((1/\bar c)/v_0)\in T_0^b\). After moving this unit
diagonal across \(s\), we obtain \(g\in A sT_0^b A\). This completes all
cases and proves (GB.4).

Finally the square formulas are \(n(c_0)^2=\operatorname{diag}(-1,-1)\)
and \(M(v_0)^2=D(v_0/\bar v_0)\), whose parameters are units. The
apartment reflection formula is respectively BF.2 or BF.7, with
\(\omega(c_0)=r\) or \(\omega(v_0)=2r\). It is exactly reflection in
\(a(h)+r=0\). \(\square\)

\textbf{Ambient torus extension.} Suppose a compact bounded torus group
\(T^b\) in the ambient group contains the displayed rank-one unit torus
and conjugates the two root charts by absolute-character factors of
valuation zero. It normalizes \(A,D,B\), and hence their subgroup
\(K_r\). The product \(T^bK_r\) is a compact subgroup. Moving its torus
factors across the root groups and across the reflection changes them
only within \(T^b\), so (GB.4) becomes exactly
\(A sT^b A\cup B T^b A=\langle T^b,A,D\rangle\). AB.10 supplies the zero
absolute-character valuations of the bounded torus. Thus once the
relative chart has been identified in the ambient group, this gives the
full torus form of the filtered two-cell statement.

The actual primary reading behind the rank-one chart is free
Bruhat--Tits II §§4.1.9--4.1.12, with the sign conversion already bound
in BF.7. The new proof does not cite the source's omitted matrix
computation: (GB.7)--(GB.10) contain the computation used here.

\subsubsection{GB.1. Relative roots and their rank-one
charts}\label{gb.1.-relative-roots-and-their-rank-one-charts}

Let \(G/k\) be connected reductive and quasi-split. The preceding
programme proof \href{../AG-RG/AG-RG-06.html}{Automorphisms, forms and
parabolic subgroups}, Proposition 5.1, gives a finite Galois splitting
field \(L/k\) and a pinning on \(G_L\) for which
\(\Gamma=\operatorname{Gal}(L/k)\) acts by based-datum automorphisms,
permuting the simple frames. Central directions and isomorphic
components may be permuted. Let \(T,B\) be the descended torus and
Borel, let \(S\) be the maximal split subtorus of \(T\), and let
\(\widetilde\Phi,\widetilde\Delta\) be the absolute roots and simple
roots. AB.10 identifies the real apartment of \(S\) with the invariant
absolute cocharacter space. All statements below concern this chosen
rational Borel pair.

\textbf{Theorem.} No absolute root restricts to zero on \(S\), and
\(Z_G(S)=T\). Thus \(S\) is a maximal split torus of \(G\). The
nondivisible relative roots form a reduced crystallographic root system.
Its finite Weyl group is \(N_G(S)(k)/T(k)\). For each nondivisible
relative root \(a\), the rank-one derived subgroup is centrally covered
by exactly one of

\[
 \operatorname{Res}_{E/k}\operatorname{SL}_2,
 \qquad \operatorname{Res}_{E/k}\operatorname{SU}_3(F/E),
 \tag{GB.11}
\]

with \(E/k\) finite separable and \(F/E\) separable quadratic. These
covers induce isomorphisms of the positive and negative unipotent
groups, including doubled-root subgroups. Every relative Weyl element
has a rational representative which permutes absolute integral root
coordinates with signs.

\textbf{Proof.} Choose an invariant Euclidean inner product on the
absolute derived root space by averaging over the finite Weyl and
diagram groups. Choose a point in the absolute positive chamber and
average its finite \(\Gamma\)-orbit. All simple-root values of the
average are positive, so it is an invariant regular point. Every
positive absolute root has positive value there; consequently no
absolute root restricts to zero. The restrictions of the simple roots
belonging to one \(\Gamma\)-orbit are equal. Restrictions from distinct
orbits are linearly independent: averaging in the absolute simple-root
basis makes them the independent orbit sums, divided by their orbit
sizes. Thus they form the simple inequalities of the positive chamber in
the invariant space.

We first determine the subsystem supported on one simple orbit \(O\),
without appealing to the classification list of Dynkin diagrams. The
graph of the simple roots has no cycle. Normalize the simple vectors to
length one. On an edge their inner product is at most \(-1/2\), since
the product of the two nonzero Cartan integers is at least one; other
inner products are nonpositive. If a cycle had \(m\) vertices, the
squared norm of their sum would therefore be at most \(m-m=0\),
contradicting positive definiteness. The induced graph on \(O\) is a
finite forest and \(\Gamma\) acts transitively on its vertices, so every
vertex has the same degree. A finite regular forest has degree zero or
one: each nonempty tree has a leaf, and a regular tree of degree one is
a single edge. Thus the subsystem is a product of \(A_1\)'s or a product
of \(A_2\)'s. In the edge case both vertices have equal length, so their
equal off-diagonal Cartan integers must be \(-1\) by positive
definiteness of the two-by-two matrix. This also verifies the asserted
\(A_2\) identification directly.

Let \(w_O\) be the product of the longest Weyl elements of these
components. In the isolated case it is the product of their commuting
simple reflections. In the edge case its factor is
\(s_i s_j s_i=s_j s_i s_j\). It commutes with \(\Gamma\), fixes
pointwise the invariant hyperplane where the common restriction \(a\)
vanishes, and sends \(a\) to \(-a\). It is therefore the orthogonal
reflection in that hyperplane. Its relative coroot is

\[
 a^\vee_{\mathrm{rel}}=
 \begin{cases}
 \sum_{i\in O}\alpha_i^\vee,&A_1\text{ components},\\
 2\sum_{i\in O}\alpha_i^\vee,&A_2\text{ components}.
 \end{cases}
 \tag{GB.12}
\]

Indeed the reflection formulas on an invariant vector have equal
simple-root values within the orbit. In the edge case the longest
element is reflection in \(\alpha_i+\alpha_j\), whose value there is
\(2a\). In either case \(a(a^\vee_{\mathrm{rel}})=2\), and the coroot is
an invariant integral cocharacter.

The absolute-root hyperplanes restricted to the invariant space form a
finite arrangement. A regular invariant point lies in one absolute
chamber, which is \(\Gamma\)-stable. Conversely, averaging an interior
point of a stable absolute chamber produces an invariant interior point.
Thus relative chambers are exactly the intersections with stable
absolute chambers. Reflection in an orbit wall sends the base relative
chamber to its neighbor: its longest component element reverses the
roots supported on that orbit while retaining positivity of every root
with another simple-root coefficient, because that coefficient is
unchanged. A generic finite wall-crossing path now shows that the group
generated by the \(w_O\) is transitive on relative chambers. It acts
freely: if an absolute Weyl element preserving the invariant space
stabilizes the relative positive chamber, it sends a regular positive
invariant point into the same absolute positive chamber, and the actual
absolute chamber theorem of \href{../AG-RG/AG-RG-04.html}{Root data,
Weyl chambers and the Bruhat decomposition}, §3, forces it to be the
identity. The same argument shows that every \(\Gamma\)-fixed absolute
Weyl element is in this generated group, by first matching its relative
chamber and then using freeness.

Any relative-root hyperplane is a wall of a relative chamber. Moving
that chamber to the base by the generated group shows that the absolute
roots on its normal ray are conjugate to the subsystem supported on one
simple orbit. The support has only the roots of \(A_1\) or \(A_2\), so
their restrictions are exactly \(\pm a\), or \(\pm a,\pm2a\). This
proves that the nondivisible relative roots are reduced. Formula
(GB.12), transported by the generated Weyl group, gives integral Cartan
pairings with all other restricted roots. Finiteness, spanning,
reflection invariance and these pairings prove the crystallographic
root-system axioms.

The absolute Bruhat-cell coordinates also prove the centralizer
assertion scheme-theoretically. Over a splitting field, a point
centralizing \(S\) maps in \(G/B\) to an \(S\)-fixed point. In each
absolute Bruhat chart all its affine coordinates have nonzero restricted
weights, so a fixed point has every unipotent coordinate zero and is a
Weyl-labelled point. The remaining Weyl label must fix \(S\) pointwise.
It fixes the regular invariant point above, hence is the identity by the
absolute chamber theorem. On closed geometric points the centralizer is
therefore \(T\). Near \(T\) the integral big-cell chart has root
coordinates of nonzero \(S\)-weight; requiring commutation as a
group-scheme identity sets them all to zero. This excludes infinitesimal
directions as well: a coordinate multiplied by a nontrivial Laurent
character can be fixed only when it is zero, over any coefficient
algebra. Hence the entire finite-type centralizer scheme is \(T\). A
split torus containing \(S\) would lie in \(T\), contrary to maximality
there proved in AB.10. Thus \(S\) is maximal in \(G\).

The pinned normalizer representatives of the orbit reflections are
rational. For isolated components take the product of \(n_i(1)\). For an
edge take \(n_i(1)n_j(1)n_i(1)\). Its equality with the reversed braid
can be checked in the pinned \(\operatorname{SL}_3\) model: both
matrices are \(\begin{pmatrix}0&0&1\\0&-1&0\\1&0&0\end{pmatrix}\). This
computation and the commuting isolated factors make the product fixed by
every diagram permutation. Each representative acts on all absolute
integral root lines by signs, by the actual pinned comparison theorem
and integral normalizer tables in \href{../AG-RG/AG-RG-05.html}{Pinnings
and the classification of split reductive groups}. Since the orbit
reflections generate the fixed Weyl group, every element has such a
rational representative. Conversely a rational normalizer of \(S\)
normalizes \(Z_G(S)=T\), and its absolute Weyl image is
\(\Gamma\)-fixed. We have proved exactly \(N_G(S)(k)/T(k)\) equals the
relative reflection Weyl group, including surjectivity on rational
representatives.

It remains to identify the rank-one groups as groups, not only their
root diagrams. For a simple orbit its supported semisimple group is
centrally covered over \(L\) by the product of the matrix groups
\(\operatorname{SL}_2\) or \(\operatorname{SL}_3\) just found. This is
an actual pinned root-datum argument: quotient the matrix group's centre
to match the intermediate character lattice of the supported derived
group, then use \href{../AG-RG/AG-RG-04.html}{Root data, Weyl chambers
and the Bruhat decomposition} Theorem 7.1 and Lesson 05 Theorem 5.1.
Those earlier proofs construct the central quotient and pinned
isomorphism over arbitrary characteristics. Their big-cell quotient
proof preserves each root group and each positive or negative
root-coordinate product, so the central cover induces actual unipotent
isomorphisms, including when its centre is non-smooth.

The Galois group permutes these components transitively. For the
stabilizer of one component, let \(E\) be its fixed field. In the
\(A_1\) case the based diagram has no nontrivial automorphism, and its
pinning descends to \(\operatorname{SL}_{2,E}\). The transitive product
is its restriction of scalars: over \(L\) this consists of its factors
indexed by the embeddings of \(E\), and invariance determines every
component from the chosen one, coefficient by coefficient. In the
\(A_2\) case the component stabilizer exchanges its two simple roots;
otherwise the original orbit would not contain both vertices of an edge.
The kernel of that action fixes a separable quadratic field \(F/E\). Its
pinned twist is exactly the special unitary matrix group in BF.7: over
\(F\), the nontrivial action is \(g\mapsto J(g^{\mathsf T})^{-1}J\).
Choose the integral simple frames \(x_{12}(t)\) and \(x_{23}(-t)\); this
automorphism exchanges those frames exactly. Pinned uniqueness
identifies this model with the given diagram action. The fixed equation
is \(g^{\mathsf T}J\bar g=J\), with determinant one. The triangular
fixed equations are precisely (GB.1). Thus its transitive product
descends to \(\operatorname{Res}_{E/k}\operatorname{SU}_3(F/E)\), with
the long root its trace-zero central coordinate. Transporting by a
rational relative Weyl representative handles every relative root. This
proves (GB.11), its exact root-group assertions, and all stated
generalities. \(\square\)

The free original Bruhat--Tits II §§4.1.1--4.1.12 supplies the primary
rank-one/orbit problem and matrix conventions. The finite-forest
argument above replaces its reference to the diagram classification; its
pinned quotient and coordinate identifications are bound to the
preceding proofs just specified.

\subsubsection{GB.2. Weighted relative valuations and the enlarged
normalizer}\label{gb.2.-weighted-relative-valuations-and-the-enlarged-normalizer}

For a nondivisible relative root \(a\), use any order of the absolute
root coordinates on its ray. If \(\beta|_S=m_\beta a\), then \(m_\beta\)
is one or two by GB.1. Define

\[
 \varphi_a(u)=\min_\beta\frac{\omega(c_\beta)}{m_\beta},\qquad
 \varphi_a(1)=\infty.
 \tag{GB.13}
\]

For a doubled root use \(\varphi_{2a}=2\varphi_a\) on its group.

\textbf{Theorem.} This definition is independent of coordinate order. It
equals \(\omega(u)\) on the restriction-of-scalars SL \(_2\) chart and
\(\omega(v)/2\) on (GB.1), with \(\varphi_{2a}(0,v)=\omega(v)\). Its
level subgroups satisfy

\[
 U_{a,r}=U_a(k)\cap
 \prod_{\beta|_S=m_\beta a}x_\beta\{c:\omega(c)\ge m_\beta r\}.
 \tag{GB.14}
\]

They are compact open in the root group and give its identity
neighborhood basis. If \(b\) is not a negative multiple of \(a\), then

\[
 [U_{a,r},U_{b,s}]\subset
 \left\langle U_{pa+qb,\,pr+qs}:p,q\in\mathbb Z_{>0},\ pa+qb\text{ a root}\right\rangle.
 \tag{GB.15}
\]

Opposite normalizer completions have the negative input value, and all
normalizer valuation changes are the affine root changes. The enlarged
normalizer acts on \(V(S)=V_D(S)\oplus V_Z(S)\) by
\(tn_w:x\mapsto wx-\nu(t)\), with compact open kernel \(T^b\) and full
discrete torus-translation lattice. The resulting apartment and its
valuation translation class do not depend on the equivariant pinning.

\textbf{Proof.} On an SL \(_2\) ray its conjugate absolute parameters
are the embeddings of one element of \(E\); all have the same valuation
by AB.0. They commute, so the result and order independence are
immediate. On an A \(_2\) component of a unitary ray, the two short
parameters are \(u,\bar u\), with signs, and the long parameter is \(v\)
in one order and \(-\bar v\) in the other. Indeed exchanging the two
short root factors changes the central parameter by \(u\bar u\), and
\(v+u\bar u=-\bar v\). Moving the central root factor does nothing.
These are all possible orders. BF.7 proves \(2\omega(u)\ge\omega(v)\),
so their weighted minimum is exactly \(\omega(v)/2\) in every order.
Different components commute and field embeddings preserve valuation.
This proves (GB.13)--(GB.14), including the long-root normalization. The
subgroup, compactness, openness and neighborhood assertions are the
actual matrix and trace calculations of BF.7, or the additive field-ball
assertion of BF.1, transported through GB.1's root-group isomorphisms.

For the commutator estimate first take linearly independent \(a,b\).
Choose a positive system containing both: a real vector can have
positive values on both precisely because they are not opposite
multiples. Its invariant absolute positive system contains all their
absolute root factors. Collect the commutator in an absolute root order
using BF.1's integral commutator identities. Collection terminates by
absolute root height. Each resulting coordinate is a polynomial with
integer coefficients in the input parameters, homogeneous for torus
weights. Give a parameter over \(a\) degree \((m_\beta,0)\), and a
parameter over \(b\) degree \((0,m_\gamma)\). Linear independence makes
the degree of an output root \(\delta\) uniquely
\((p_\delta,q_\delta)\), where \(\delta|_S=p_\delta a+q_\delta b\). Both
coordinates are positive integers for a nonzero commutator output: each
term vanishes when either input family is zero, and the collection
identities create only positive sums. Its valuation is therefore at
least \(p_\delta r+q_\delta s\). Integer structure constants have
nonnegative valuation in every characteristic.

Regroup these absolute coordinates by nondivisible relative rays. The
block order is \(\Gamma\)-stable; uniqueness of the absolute root
factorization makes each resulting block rational. If a block has a
nonzero short parameter, its root \(c=p_ca+q_cb\) has positive integral
coefficients, and its long coordinates, if any, have degree
\((2p_c,2q_c)\). Formula (GB.13) puts the block in \(U_{c,p_cr+q_cs}\).
If \(c\) has nonintegral coefficients in the basis \(a,b\), no short
absolute coordinate of that degree can occur. The block then has only
its long-root coordinate, of integral degree \((p_{2c},q_{2c})\), and
lies in \(U_{2c,p_{2c}r+q_{2c}s}\). A block with zero short coordinates
and integral degree is already covered by the former bound. These cases
exhaust every ray and prove (GB.15). For positive multiples on one ray,
the only nonzero case is the short-short unitary commutator: BF.7 gives
\([U_{a,r},U_{a,s}]\subset U_{2a,r+s}\). Its long-root group is central;
SL \(_2\) is additive. This proves all same-ray cases too.

For rank-one completions use exactly BF.1 and BF.7 and transport by the
rational relative Weyl representatives of GB.1. They prove existence and
uniqueness of the two opposite factors, each of value \(-\varphi_a(u)\),
and the reflection formula. In the absolute coordinates, conjugation by
\(t\in T(k)\) scales a \(\beta\)-parameter by \(\beta(t)\). AB.10 makes
its valuation \(\beta(\nu(t))=m_\beta a(\nu(t))\). Thus (GB.13) changes
by \(a(\nu(t))\). The pinned Weyl representatives permute coordinates
with signs and preserve their valuations. This proves the stated
normalizer covariance and affine signs on every root group; a doubled
root has twice the short-root change.

The absolute coroot space and the central space are stable under
\(\Gamma\), so taking invariants gives the displayed direct sum. AB.10
gives \(X_*(S)\subset\nu(T(k))\subset e_L^{-1}X_*(S)\), and its compact
open kernel. The rational Weyl representatives identify the full
normalizer as \(T(k)\{n_w\}\), so its affine kernel is exactly \(T^b\),
its image is a lattice with a finite linear Weyl part, and its action is
proper with compact orbit space on the enlarged apartment. As in BF.3's
elementary Euclidean argument, compact subgroups have fixed points by
averaging and their fixed affine spaces are contractible; invariant
partitions give the usual affine equivariant maps. Thus this normalizer
apartment is its numerable universal proper space, including the central
factor.

Finally a change of equivariant pinning multiplies the simple frames by
constants \(c_i\) with Galois-equivariant indices. Pinned transport
multiplies each other root frame by the corresponding product of the
\(c_i\), with its simple-root exponents; this is the torus-weight
relation of the actual pinned comparison proof. Galois invariance of the
extended valuation makes \(\omega(c_i)\) constant on each orbit. There
is a unique invariant derived vector \(h\) with
\(\alpha_i(h)=-\omega(c_i)\); the simple roots are a basis on the
absolute derived apartment. For a fixed group element with old parameter
\(u\), the new parameter is \(u/c_i\); its new valuation is
\(\omega(u)-\omega(c_i)=\omega(u)+\alpha_i(h)\). This is exactly the
valuation-origin translation with the pullback convention of BF.2. For a
ray with \(\beta|_S=m_\beta a\), division by \(m_\beta\) in (GB.13)
gives the same translation \(a(h)\). Hence the affine translation class
is intrinsic. Central directions are retained as the independent
enlargement. \(\square\)

Reference: Bruhat--Tits II, Sections 4.2.2--4.2.11. The
coordinate-order, commutator, normalizer and pinning arguments are
proved above.

\subsubsection{GB.3. Compact face groups and every-order relative
Iwahori
coordinates}\label{gb.3.-compact-face-groups-and-every-order-relative-iwahori-coordinates}

On \(V_D(S)\), remove the hyperplanes \(a(v)+r=0\) for nondivisible
relative roots and attained values \(r\). The arrangement is locally
finite: by AB.0 and GB.2 every such value belongs to
\((2e_L)^{-1}\mathbb Z\). A chamber is a connected component of its
complement. Define

\[
 P_v=\langle T^b,U_{a,-a(v)}:a\text{ nondivisible}\rangle,
 \qquad I_C=P_v\quad(v\in C).
 \tag{GB.16}
\]

The doubled-root generators are already contained in the short-root
ones, so this agrees with using all relative roots. For any positive
relative system and any orders on its positive and negative nondivisible
roots,

\[
 \left(\prod_{a<0}U_{a,-a(v)}\right)\times T^b\times
 \left(\prod_{a>0}U_{a,-a(v)}\right)\xrightarrow{\sim} I_C
 \tag{GB.17}
\]

is a homeomorphism. All \(P_v\) are compact open; \(I_C\subset P_{v_0}\)
has finite index for every \(v_0\in\overline C\). For \(x=(v,z)\) the
groups \(K_x=P_vN_x\) are compact open, normalizer covariant, and
satisfy \(K_x\cap N=N_x\).

\textbf{Proof.} On a chamber, membership of an element in
\(U_{a,-a(v)}\) is constant: its sign \(\varphi_a(u)+a(v)\) can change
only at one of the removed hyperplanes. For a nonidentity member it is
strictly positive. In absolute coordinates this says
\(\omega(c_\beta)+\beta(v)>0\), by GB.2. Expand any finite word in the
relative root generators into absolute root factors over \(L\); move its
unit-valued torus factors to the left. Apply the exact finite-polynomial
big-cell argument proved in BF.6. We spell out why its hypotheses and
bounds are unchanged here. The word-map denominator has constant term
one and all its other terms have torus weight zero; each such
nonconstant monomial has strictly positive valuation from the strict
bounds on its variables. Thus the denominator is a valuation unit. Each
output root coordinate has its corresponding torus weight, so every
output monomial has \(\omega(c_\beta)+\beta(v)>0\); every torus
character and its inverse have unit value. The calculation uses integer
coefficient polynomials on one finite word and works with the fractional
value group of \(L\); it needs neither an absolute alcove containing
\(v\) nor a ramified descent theorem.

Choose an absolute positive system restricting to the chosen positive
relative system. Group the arbitrary-order absolute root factors into
the nondivisible relative rays, with any of the requested orders. GB.1
identifies each ray block as a rational root group. More precisely each
block subgroup is Galois stable, and the block order itself is fixed.
The unique absolute factorization of a rational product is therefore
fixed block by block, even in an A \(_2\) block where coordinate
conjugation is nonlinear. The torus factor is rational as well. The
output strict weighted bounds, divided by the ray multiplicities, put
every block back in the level group of (GB.17), and put the torus in
\(T^b\). Conversely all those factors are generators. The absolute
big-cell uniqueness gives uniqueness of this product. Its inverse
coordinates are regular on that cell, hence continuous; its forward map
is multiplication. This proves the homeomorphism and constancy of
\(I_C\).

The rational big cell used here is itself an open chart: the stable
absolute open immersion descends coefficient by coefficient to
\(U^-\times T\times U^+\to G\). Its coordinate-ring kernel and cokernel
vanish after the faithful extension to \(L\), hence already over \(k\);
its image open is described by the product of the Galois conjugates of
an absolute cell denominator. Thus (GB.17), whose factors contain
identity neighborhoods in their root groups and torus, contains an open
identity neighborhood in \(G(k)\). Consequently \(I_C\) is open.

For compactness, use BF.4's faithful integral representation of the
pinned split model over \(L\), with its absolute weight decomposition.
On the norm with weight vector \(v+z\), every absolute root factor in a
relative generator has \(\omega(c_\beta)+\beta(v)\ge0\). BF.5's
weight-polynomial estimate consequently makes that factor and its
inverse norm isometries. The torus \(T^b\) has unit values for every
absolute character by AB.10, and is an isometry as well. The group
generated by these factors lies in the compact norm-isometry group in
\(\operatorname{GL}_N(L)\). The closed immersion and the closed subset
\(G(k)\subset G(L)\) make its intersection with \(G(k)\) compact. Choose
a chamber with \(v\) in its closure. Its \(I_C\) is contained in
\(P_v\), since the strict inequalities become nonnegative in the limit.
A subgroup containing an open subgroup is open and closed. Thus \(P_v\)
is closed in that compact group and is compact. Openness also follows.
Its cosets modulo \(I_C\) are compact and discrete, hence finite.

Normalizer covariance follows from GB.2. The normalizer's affine image
is a discrete lattice with finite Weyl group, so a point stabilizer
consists of finitely many cosets of \(T^b\); it is compact and open in
\(N\). It normalizes \(P_v\), giving the compact open product \(K_x\).
If a normalizer element is in \(P_v\), its norm action in the faithful
representation preserves the norms of all weight vectors. The integral
pinned Weyl representatives are invertible on their integral weight
lattices, and the torus contribution changes their norms by its
character valuations. Exactly the norm-detection calculation in BF.5
therefore forces its affine action to fix \(v+z\). This holds for every
central \(z\), since roots vanish on that factor. Hence
\(P_v\cap N\subset N_x\), and multiplying by \(N_x\) proves
\(K_x\cap N=N_x\). \(\square\)

\subsubsection{GB.4. Relative walls, the pivot identity and affine
exchange}\label{gb.4.-relative-walls-the-pivot-identity-and-affine-exchange}

\textbf{Lemma.} The relative arrangement is reflection invariant and has
bounded reduced chambers. Its wall reflections form an affine Coxeter
group acting simply transitively on its chambers; closed chambers are
fundamental domains, and point stabilizers are generated by their walls
through the point.

\textbf{Proof.} For every attained \(r\), the rank-one completion of
GB.0 is a rational normalizer element acting as reflection in
\(a(v)+r=0\). GB.2's covariance maps every root level and every
corresponding wall to another level and wall. Thus the entire
arrangement is invariant. For any root \(a\), the integral cocharacter
\(a^\vee_{\mathrm{rel}}\) of GB.1 evaluated at \(\pi\) shifts its root
values by two. From one attained level we therefore get an unbounded
arithmetic sequence in both directions. The values of \(a\) on any
chamber lie between consecutive walls of this sequence, so are bounded.
The roots span the dual derived apartment, making its chamber bounded.
Its closure is a compact polytope by local finiteness.

The affine reflection group is a subgroup of the normalizer's discrete
lattice-with-finite-linear-part image. The complete gallery-disc proof
of AB.2 now applies verbatim in its proved general form: it only uses a
locally finite reflection-invariant hyperplane arrangement, bounded
chambers and finite linear point stabilizers. At a codimension-two
intersection the finite linear stabilizer makes consecutive sectors
dihedral. Backtracking and these local dihedral braids reduce a closed
gallery. This proves transitivity, freeness and the Coxeter
presentation, with actual finite pair orders and no relation for
infinite pair orders. The local-star argument proves the closed
fundamental domain and wall-generated point stabilizers. No
identification with the split coroot lattice is asserted for this
possibly differently spaced arrangement. \(\square\)

We record one additional local matrix identity needed to use the
two-cell theorem for exchange. With \(A,D,s,r\) as in GB.0 and
\(u\in A\setminus A^+\),

\[
 s u\in A T_0^b D.
 \tag{GB.18}
\]

In SL \(_2\) , the lower-right entry of \(n(c_0)x_+(c)\) is \(-c/c_0\),
a unit when \(\omega(c)=r\). Its Gaussian factorization with that
lower-right pivot has upper parameter of valuation at least \(r\), lower
parameter at least \(-r\), and unit diagonal. In SU \(_3\) , the
lower-right entry of \(M(v_0)x_+(u,v)\) is \(v/\bar v_0\), again a unit.
For any weighted norm isometry \(g\) with unit \(d=g_{33}\), the upper
pair in its reverse Gauss factorization is \((g_{23}/d,g_{13}/d)\), its
lower pair is \((-\overline{g_{32}/d},g_{31}/d)\), and its torus
parameter is \(t=1/\bar d\). The last column and last row unitary
equations give both trace relations. Multiplying the upper factor
inverse kills its two upper last-column entries; unitary pairings force
the remainder lower triangular with diagonal \(D(t)\). The entry bounds
(GB.5) place these pairs in \(A,D\), respectively. This proves (GB.18)
with every parameter bound, including characteristic two. The ambient
unit torus can be used by GB.0's last paragraph.

\textbf{Proposition.} Let \(C\) be a relative chamber and \(I=I_C\). For
each of its walls choose \(\alpha=a+r\) positive on \(C\), with \(a\)
nondivisible, and the rank-one representative \(s=m(u_0)\) at that
level. For every \(n\in N\),

\[
 I s I n I\subset I s n I\cup I n I,
 \qquad s I s^{-1}\ne I.
 \tag{GB.19}
\]

\textbf{Proof.} At this wall the chamber's positive-ray group is
\(A=U_{a,r}\) and its negative-ray group is \(B=U_{-a,-r+}\). To check
the second equality note that the positive and negative value sets
coincide by the rank-one charts and are closed under negation by the
unique opposite completion. A chamber neighboring the wall has no
intervening attained parallel level. This proves exactly the strict
opposite level, even when the gaps of attained SU \(_3\) values are
unequal.

Crossing the wall changes only these two ray groups. Choose a positive
system containing \(a\) and order its factors with \(a\) last. GB.3
gives \(I=J A\), with \(J\subset I\cap s^{-1}Is\). As in AB.3 it
suffices to examine \(s u n\), \(u\in A\). If \(u\in A^+\), its
reflection conjugate is in \(B\), giving \(I s n I\). If \(u\) is at the
wall level and \(n^{-1}\alpha\) is positive on \(C\), its conjugate
under \(n^{-1}\) is in \(I\), giving the same cell. If that affine root
is negative, (GB.18) gives \(s u\in A T^b D\), and
\(n^{-1}D n\subset I\) because \(n^{-1}(-\alpha)\) is positive. This
gives \(InI\). These are all cases because an attained affine root has
constant nonzero sign on a chamber. Finally \(sAs^{-1}\) contains a root
element at the opposite level \(-r\), which is excluded by the strict
level of \(B\). The unique factorization of GB.3 proves it is not in
\(I\). Thus conjugation changes the chamber group. \(\square\)

\subsubsection{GB.5. Root generation, full-group cells and the relative
affine
labels}\label{gb.5.-root-generation-full-group-cells-and-the-relative-affine-labels}

Let \(N_{\mathrm{aff}}\) be generated by \(T^b\) and the chamber-wall
representatives of GB.4, and let \(W_{\mathrm{wall}}\) be their actual
affine reflection image. Put \(G^0=\langle I,N_{\mathrm{aff}}\rangle\).

\textbf{Proposition.} One has

\[
 G^0=\langle T^b,U_a(k):a\text{ a relative root}\rangle,
 \qquad G(k)=T(k)G^0=\Omega_CG^0,
 \tag{GB.20}
\]

where \(\Omega_C=\{n\in N:nC=C\}\), allowing central translations. The
subgroup \(G^0\) is open and normal, \(I\cap N=T^b\), and

\[
 N/T^b\xrightarrow{\sim}I\backslash G/I,\qquad
 G^0\cap N=N_{\mathrm{aff}},\qquad
 \bigcap_{n\in N}nIn^{-1}=T^b.
 \tag{GB.21}
\]

Any two points and any two chambers in the relative gluing lie in a
common translated enlarged apartment.

\textbf{Proof.} Denote the root-generated group on the right in (GB.20)
by \(R\). All chamber-wall representatives are opposite-root
completions, so \(I,N_{\mathrm{aff}}\subset R\). Conversely every
reflection in any wall of the arrangement is a conjugate of a
base-chamber wall reflection by GB.4. Its rank-one representative
differs from that conjugate only by the affine kernel \(T^b\). Thus it
is in \(N_{\mathrm{aff}}\). Take a root \(a\) and a nonzero root element
of value \(r\); its conjugates by \(a^\vee_{\mathrm{rel}}(\pi)^n\) have
values \(r+2n\) by GB.2. Their reflection products yield translations by
nonzero multiples of \(a^\vee_{\mathrm{rel}}\), in both directions,
inside \(N_{\mathrm{aff}}\). Conjugation by their powers moves the fixed
chamber tail for \(U_a\) to arbitrarily low levels. These tails exhaust
every element of \(U_a(k)\). Hence \(R=G^0\), including every
doubled-root subgroup. It is open because it contains \(I\).

We prove rational full-group generation explicitly. Apply the actual
split Bruhat proof \href{../AG-RG/AG-RG-04.html}{Root data, Weyl
chambers and the Bruhat decomposition}, Theorem 6.2, over \(L\), to the
rational element \(g\). Its absolute Bruhat label is fixed by
\(\Gamma\), since the Borel is stable and its cells are distinct. GB.1
identifies that label as a relative Weyl label with a rational
representative \(n_w\). The absolute cell has unique coordinates in
\(U_w\times T\times U^+\) after selecting that representative. Its
inversion-root set is Galois stable and is a union of whole relative
rays: all absolute roots on a ray have the same sign before and after
the relative Weyl element. Group the coordinates by these rays.
Uniqueness forces each block, the torus factor and the positive factor
to be rational. Thus \(g\) is a product of rational relative root
elements, a torus element and \(n_w\).

For each simple relative reflection choose any nonidentity root element
and its rational opposite completion. It represents that reflection. The
rational representative used above differs from it by an element of
\(T(k)\). Since \(T(k)\) normalizes all relative root groups and
\(T^b\), the set \(T(k)R\) is a subgroup. It contains these
representatives and all the displayed Bruhat factors, proving
\(G=T(k)R\). It also proves normality of \(R\): it is normalized by
\(T(k)\), and by its own elements. This argument does not assume that a
particular integral unitary Weyl representative is itself an attained
level-zero completion.

The normalizer preserves the arrangement. Its chamber action and GB.4
give \(N=N_{\mathrm{aff}}\Omega_C\); elements of \(\Omega_C\) preserve
the root-level chamber groups and normalize \(I\). Since \(G^0\) is
normal and contains \(N_{\mathrm{aff}}\), this proves the last
full-group formula in (GB.20). If an element of \(I\cap N\) fixes all
points of \(C\times V_Z\), its affine action is the identity and it is
in \(T^b\) by GB.2; conversely \(T^b\subset I\). The fixing assertion
follows from GB.3's \(P_v\cap N\subset N_{(v,z)}\), for every chamber
point and every \(z\).

We now have exactly the hypotheses of the locally proved exchange
theorem AB.4 for \((G^0,I,N_{\mathrm{aff}})\): the quotient is the
Coxeter group of GB.4, (GB.19) is exchange, the square is in \(T^b\),
and conjugation changes \(I\). Its proof gives distinct affine labels
and their multiplication rules. The wall group is normal in the full
affine normalizer image, since conjugation maps every wall reflection to
another. Decompose that image into its chamber stabilizer and the simply
transitive wall group. The explicit extended-label induction of AB.4
therefore applies and proves (GB.21), with this actual relative wall
group in place of the split coroot semidirect product. No additional
extended BN-pair theorem is cited.

For the core equality, use GB.3's unique root-coordinate factorization
of an element of \(I\). A split cocharacter scales every absolute
coordinate on a relative ray by its nonzero integral relative-root
pairing. Powers with the appropriate sign would lower any nonzero root
parameter below its fixed chamber bound. Membership in all those torus
conjugates of \(I\) thus forces all root coordinates to be zero. The
remaining factor is \(T^b\), which is normal in \(N\). This proves the
core assertion. Finally for two point or chamber representatives
normalize their reduced coordinates into \(\overline C\) and write their
relative group element as \(i_1ni_2\), using the full double-coset
partition. The translated apartment \(gi_1V(S)\) contains both, exactly
as in the explicit argument AB.5. \(\square\)

\subsubsection{GB.6. The full quasi-split enlarged local-field
building}\label{gb.6.-the-full-quasi-split-enlarged-local-field-building}

\textbf{Theorem.} For every connected quasi-split reductive group over
the original nonarchimedean local field, the quotient formed from
\(G(k)\), its enlarged apartment and the groups \(K_x\) of GB.3 is
locally compact, second countable and Hausdorff, with continuous proper
cocompact action. It has a compatible proper complete CAT(0) metric,
Euclidean on each apartment. It is a numerable universal proper
\(G(k)\)-space. This includes all residue characteristics, all
ramification of its splitting field and all nonsplit central tori.

\textbf{Proof.} The hypotheses of the set and proper-topology proof AB.1
are now all proved. The normalizer action is the
lattice-with-finite-Weyl action of GB.2; compact open \(K_x\),
covariance and normalizer intersection are GB.3. The root-level groups
have local upper-containment at every point: near a point the only
possible level jumps are its finitely many incident walls, and at the
wall the non-strict level includes both neighboring groups. Hence nearby
\(P_y\subset P_x\). The normalizer fixed-point labels have the same
upper-containment by the proper discrete affine action and closed fixed
hyperplanes. Thus \(K_y\subset K_x\) nearby, the graph is closed, and
their union over a compact apartment set is compact, by AB.1's local
finite covering argument.

A closed reduced chamber is compact by GB.4. The central subspace is
rational in the invariant cocharacter lattice, so a full lattice of
central cocharacters has a compact parallelepiped. The product of these
two sets has locally finite normalizer translates covering the enlarged
apartment. Consequently every step of AB.1's proper finite-fiber
quotient proof applies: it proves exactly the asserted quotient
topology, compact transporters and compact orbit representatives.

For \(v\in\overline C\), \(P_v\subset G^0\). An affine Bruhat
representative in \(P_v\) fixes \(v\) by GB.3, so its label is in the
finite subgroup generated by the walls through \(v\), by GB.4.
Conversely each such wall representative is a product of opposite root
factors of value zero there, and lies in \(P_v\). The precise face
equality is therefore

\[
 P_v=\coprod_{w\in\langle J_v\rangle}I n_wI.
 \tag{GB.22}
\]

This is exactly the face-group input for the chamber realization and
retractions, not merely containment in a point stabilizer. Since
\(G^0\cap N=N_{\mathrm{aff}}\), the point relation and the
closed-chamber fundamental-domain argument give the enlarged
homeomorphism \(\mathcal B\simeq\mathcal B^0\times V_Z(S)\), by the
explicit set and proper-topology proof AB.5. The wall reflections and
their translations have zero central component; the full normalizer
supplies the retained central translations. The central projection
therefore gives the full group action on that factor.

The chain-metric proof AB.6 applies to these proved affine chamber and
face data. It constructs the face-independent chamber retractions, gives
their 1-Lipschitz property and the exact distance equality when an
endpoint is in the closed center chamber, and proves every apartment is
Euclidean and every pair of points has an apartment segment. The proof
AB.7 uses just that equality and Euclidean squared-distance identity to
prove the strong convexity and full CAT(0) comparison. Its type-fiber
proof gives completeness with arbitrary panel multiplicity; here each
compact panel group contains the open \(I\) with finite index, so
reduced double-coset products give finite panel multiplicities and
finitely many chambers meeting each metric ball. The resulting
compact-ball argument gives properness and equality of metric and
quotient topologies. Product with the entire Euclidean center preserves
completeness, properness, curvature and the topology. All the actual
statements and inequalities used are proved in AB.6--AB.7; none depends
on the split coroot-lattice name.

Finally AB.8's explicit proper CAT(0) argument gives compact-subgroup
fixed points and contractible fixed sets, a finite invariant slice
partition from the compact metric orbit space and compact open
stabilizers, and continuous equivariant squared-distance minimizers for
every numerable proper source. Its unique-segment homotopies give
uniqueness. This proves universality and numerability for the full
group. \(\square\)

GB.0--GB.6 prove the quasi-split relative-root, filtration, Iwahori,
rank-one, wall-exchange and metric statements. GB.7--GB.15 give
connected models, unramified arithmetic, stable apartments and exact
component descent. GB.16--GB.18 prove quasi-splitting over a finite
unramified extension, and GB.19--GB.23 prove the rational root and
maximal-apartment descent and the full arbitrary-group building theorem.

\subsubsection{GB.7. Smooth integral models of every filtered relative
root
group}\label{gb.7.-smooth-integral-models-of-every-filtered-relative-root-group}

This construction supplies actual models, including wildly ramified
unitary groups in residue characteristic two. The base in this section
may be a local field, or any henselian discretely valued field for which
the finite separable root-coordinate extensions have finite free
valuation rings. GB.9 proves this latter condition for the maximal
unramified extension of a local field. No completeness or compactness is
needed for the construction.

\textbf{Proposition.} For every nondivisible relative root \(a\) and
every real \(r\), the group \(U_{a,r}\) is the group of integral points
of a smooth affine group scheme \(\mathcal U_{a,r}\) with generic fiber
\(U_a\). Its underlying scheme is affine space, its special fiber is
connected and unipotent, and the inclusion into \(U_a\) uses precisely
the coordinates and levels of GB.1--GB.2. The corresponding statement
holds for \(U_{2a,r}\) when \(2a\) is a root.

\textbf{Proof in the reduced case.} In the notation of GB.1 the root
group is a restriction of scalars of the additive group of a finite
separable extension \(E\). The lattice
\(I_r=\{c\in E:\omega(c)\geq r\}\) is a principal fractional ideal of
its valuation ring. Addition on this finite free base-ring module
defines \(\mathcal U_{a,r}\). This is affine space with additive
multiplication; its generic parameter is the original \(c\), not a
renormalized valuation.

\textbf{Proof in the unitary case.} Work first over \(E\), with the
separable quadratic extension \(F/E\), and keep \(\omega\) normalized on
the original base field. The upper group has parameters \[
 H=\{(u,v)\in F^2:\operatorname {Tr}(v)=-u\bar u\},\qquad
 (u,v)(w,z)=(u+w,v+z-\bar u w).
 \tag{GB.23}
\] The level \(r\) is exactly \(\omega(v)\geq2r\), by BF.7. Put \[
 I=\{v\in F:\omega(v)\geq2r\},\qquad
 \operatorname {Tr}(I)=\eta\mathcal O_E,\qquad
 L=\{u\in F:\omega(u\bar u)\geq\omega(\eta)\}.
 \tag{GB.24}
\] The trace is nonzero because the quadratic extension is separable;
its image on a nonzero fractional ideal is a nonzero fractional
\(\mathcal O_E\)-ideal. Choose \(c\in I\) with
\(\operatorname {Tr}(c)=\eta\) and put \(\lambda=c/\eta\). Thus
\(\operatorname {Tr}(\lambda)=1\), with no division by two. Conjugation
preserves \(\omega\), so \(L\) is a principal fractional
\(\mathcal O_F\)-ideal. The trace inequality gives
\(\omega(\eta)\geq2r\), and hence every \(u\in L\) satisfies
\(\omega(u)\geq r\).

Let \(J=I\cap\ker\operatorname {Tr}\). The trace-zero \(E\)-space has
dimension one, and its intersection with the fractional ideal \(I\) is a
free rank-one \(\mathcal O_E\)-lattice; write \(J=\delta\mathcal O_E\).
Choose an \(\mathcal O_E\)-basis \(b_1,b_2\) of \(L\). Every integral
filtered point is uniquely \[
 u=b_1s+b_2t,\qquad v=-\lambda\,u\bar u+\delta z,\qquad
 s,t,z\in\mathcal O_E.
 \tag{GB.25}
\] Indeed the equation forces \(u\bar u\in\eta\mathcal O_E\), so
\(u\in L\). Conversely \(-\lambda u\bar u=-c(u\bar u/\eta)\) is in \(I\)
and has the required trace. Subtracting it from any permitted \(v\)
leaves a unique element of \(J\).

The polynomial \(u\bar u/\eta\), in \(s,t\), has integral coefficients.
The two square coefficients do because
\(\omega(b_i\bar b_i)\geq\omega(\eta)\). Its mixed coefficient is
\(\operatorname {Tr}(b_1\bar b_2)/\eta\), which is integral by the same
inequality and the trace inequality. Thus (GB.25) is an actual morphism
over \(\mathcal O_E\), with values in the fractional lattice \(I\).

Transporting multiplication through (GB.25) gives \[
 (u,z)(w,z')=(u+w,z+z'+C(u,w)/\delta),\quad
 C(u,w)=\lambda(u\bar w+w\bar u)-\bar u w.
 \tag{GB.26}
\] Here \(\operatorname {Tr}(C)=0\). Also
\(u\bar w+w\bar u\in\eta\mathcal O_E\), so its product with \(\lambda\)
is in \(I\); and \(\omega(\bar u w)\geq\omega(\eta)\geq2r\). Therefore
\(C\in J\). Applying this to each of the four pairs \(b_i,b_j\) proves
that the bilinear polynomial \(C/\delta\) has integral coefficients.
Equation (GB.26), its evident identity, and its inverse define a group
law on \(\mathbf A^3_{\mathcal O_E}\). Associativity holds generically
by the matrix product (GB.23), hence everywhere by injection of its
torsion-free polynomial coordinate ring into its generic fiber. The
central \(z\)-line and the additive \((s,t)\)-quotient give a central
extension of two additive groups. Its special fiber is consequently
connected and unipotent, and the entire scheme is smooth. Its generic
fiber, integral points and valuation are exactly the required \(H\),
\(U_{a,r}\) and \(\omega(v)/2\).

For the long root use the additive lattice
\(\{v\in\ker\operatorname {Tr}:\omega(v)\geq r\}\); it is a free
rank-one \(\mathcal O_E\)-module. The lower root groups use exactly the
same construction in the displayed lower matrix chart of BF.7.

Finally restrict scalars from \(\mathcal O_E\) to the original valuation
ring \(R\). This operation here requires no general Weil-restriction
theorem: choose a finite free \(R\)-basis of \(\mathcal O_E\), expand
each coordinate in that basis, and expand the displayed polynomial law
using the integral multiplication table of \(\mathcal O_E\). The
resulting affine space over \(R\) represents the required functor. Its
smoothness and connected special fiber are explicit; the same central
additive filtration proves unipotence. Its generic fiber is the
restriction of scalars of the original rank-one chart, and its
\(R\)-points are exactly the prescribed lattice points. The central
covering in GB.1 is an isomorphism on each root group, including
inseparable central kernels, so these models apply to the root groups of
the original reductive \(G\). \(\square\)

The source of the trace-lattice construction is the freely accessible
Bruhat--Tits II, 4.3, together with its explicit rank-one coordinates in
4.1.9--4.1.12. All the ideal choices, integral coefficients and scheme
assertions required here have been proved above; a source reference does
not replace any step.

\subsubsection{GB.8. Smoothening a bounded affine group model, with its
components
retained}\label{gb.8.-smoothening-a-bounded-affine-group-model-with-its-components-retained}

Let \(R\) be a henselian DVR with perfect residue field \(f\),
uniformizer \(\pi\), and fraction field \(K\). In the applications to
point stabilizers below, \(f\) is algebraically closed. A bounded subset
of a finite-dimensional \(K\)-space means one contained in a fractional
lattice. Boundedness of a subgroup of a linear algebraic group includes
its inverses and is tested in a faithful closed representation.

\textbf{Lemma (group smoothening).} If \(\mathcal H\) is a flat affine
group scheme of finite presentation over \(R\) with smooth generic
fiber, there is a finite sequence of affine group dilatations
\(\mathcal H^{\mathrm{sm}}\longrightarrow\mathcal H\) with smooth final
model, unchanged generic fiber and exactly unchanged \(R\)-points.

\textbf{Proof.} Write \(M=e^*\Omega_{\mathcal H/R}\). The generic fiber
dimension is \(d\). Translation gives \[
 \Omega_{\mathcal H/R}\simeq\mathcal O_{\mathcal H}\otimes_R M.
 \tag{GB.27}
\] One verifies this without a smoothness assumption: a derivation at a
point \(g\) is translated to a derivation at the identity by the
differential of \(x\mapsto g^{-1}x\); the two translations are inverse
on every square-zero extension. They are regular in \(g\), by the Hopf
multiplication and antipode. This represents the displayed isomorphism.
The finite \(R\)-module \(M\) has generic rank \(d\); let \(\delta\) be
the length of its torsion.

Let \(Y\) be the reduced closure in \(\mathcal H_f\) of the reductions
of \(\mathcal H(R)\). It is a closed subgroup: products of the dense
reduction subset are in that subset, its product with itself is dense in
\(Y\times_fY\), and \(Y\times_fY\) is reduced because \(f\) is perfect.
Pulling back the ideal of \(Y\) by multiplication and inverse therefore
gives zero. A reduced finite-type group over a perfect field is smooth.
Here is the needed argument: a reduced finite-type variety over a
perfect field has a nonempty smooth open subset on each irreducible
component, by the separating-transcendence-basis Jacobian criterion;
translating a smooth geometric point to the identity, and then to each
other geometric point, proves smoothness of the group everywhere. Thus
\(Y\) is smooth.

Dilatate along \(Y\). If \(A=\mathcal O(\mathcal H)\) and \(J\) is its
inverse-image ideal, the dilatation has coordinate ring
\(A[J/\pi]\subset A[1/\pi]\). It is finite type and flat over the DVR,
has unchanged generic fiber, and is universal for maps from flat
\(R\)-schemes whose special fibers land in \(Y\). All these assertions
are also proved explicitly in the preceding programme lesson
\href{../AG-GS/AG-GS-07.html}{Néron models}, Proposition 4.1. The
universal property gives a group structure because \(Y\) is a subgroup.
Every original \(R\)-point reduces into \(Y\), so the new and old
\(R\)-points are exactly the same.

We supply the strict defect decrease, rather than invoke a smoothening
theorem whose proof has been omitted. If \(\delta>0\), put
\(t=\dim_f(M/\pi M)>d\) and \(q=t-d\geq1\). Near the identity embed
\(\mathcal H\) in a smooth affine \(R\)-scheme \(Z\) of relative
dimension \(t\), with \(\Omega_{Z/R}|_Y\to\Omega_{\mathcal H/R}|_Y\) an
isomorphism. To obtain \(Z\), start with any polynomial presentation and
eliminate the relations whose Jacobian minors are units at the identity.
There are precisely enough of these to leave \(t\) variables. Shrink at
the identity; (GB.27) makes the remaining rank \(t\) constant along
\(Y\) there. Since \(Y\) is smooth, choose étale coordinates
\(y_1,\ldots,y_m,z_1,\ldots,z_{t-m}\) on this ambient neighborhood with
the ideal of \(Y\) equal to \(J=(\pi,z_1,\ldots,z_{t-m})\).

For any defining equation \(F\) of \(\mathcal H\) in \(Z\), the
differential \(dF|_Y\) is zero, by that cotangent isomorphism.
Expressing \(F\in J\) in these coordinates gives \[
 F=\pi G+\sum_j z_jG_j,\qquad G_j\in J.
 \tag{GB.28}
\] At every \(R\)-section reducing into this neighborhood of \(Y\),
\(F=0\) and \(J\subset\pi R\), whence \(G\) has zero reduction. Those
reductions are dense in \(Y\) by its definition, so \(G\in J\) too.
Consequently every defining equation is in \(J^2\).

The ambient dilatation has coordinates \(y_i,z'_j=z_j/\pi\).
Substitution into \(F\in J^2\) makes \(F/\pi^2\) regular; it vanishes on
the dilatation of \(\mathcal H\), since it vanishes generically and its
coordinate ring is torsion-free. At its identity section, if the old
differential row is \(\sum b_i\,dy_i+\sum c_j\,dz_j\), the new relation
row is \[
 \sum_i(b_i/\pi^2)\,dy_i+\sum_j(c_j/\pi)\,dz'_j.
 \tag{GB.29}
\] These coefficients are integral by the regular divided equation. The
old relation matrix has generic rank \(q\). Over a DVR, diagonalizing a
matrix by elementary row and column operations shows that the minimum
valuation of its nonzero \(q\)-minors is exactly the torsion length of
its cokernel. Choose a minor attaining the old minimum \(\delta\). The
corresponding minor of (GB.29) divides it by at least \(\pi^q\), since
each of its \(q\) columns has been divided by \(\pi\) or by \(\pi^2\).
The new relation module contains these rows. Its generic rank is still
\(q\), so its minimum minor valuation is at most \(\delta-q\). Therefore
\(\delta_{\mathrm{new}}\leq\delta-1\).

After at most the initial \(\delta\) steps the identity is smooth. For
clarity, the criterion used here is the unit-minor Jacobian criterion:
torsion length zero lets one cut out a smooth ambient scheme of the
generic dimension; the remaining ideal vanishes in its regular local
domain because it vanishes generically and the source is flat. This is
the complete argument of the actual \href{../AG-GS/AG-GS-07.html}{Néron
models}, Lemma 4.2. Translation makes all geometric special-fiber points
smooth; the generic fiber was already smooth. Thus the final group is
smooth over \(R\). This proves the lemma, including termination.
\(\square\)

The argument in (GB.28)--(GB.29) is the group version of the locally
proved defect calculation in that lesson's Lemma 4.3. The freely
accessible Edixhoven--Romagny author paper, Section 5.1--5.6, supplies
the dilatation/smoothening framework.

\textbf{Lemma (integral points detect functions).} Assume now that \(f\)
is algebraically closed. If \(X=\operatorname {Spec}A\) is smooth affine
of finite presentation over \(R\), and \(F\in A[1/\pi]\) has integral
value at every \(X(R)\)-point, then \(F\in A\). Hence a generic morphism
from \(X_K\) to any affine finite-type \(R\)-scheme extends uniquely if
all its coordinate functions have integral values at every
\(X(R)\)-point.

\textbf{Proof.} Express \(F=\pi^{-n}g\), with \(g\in A\), choosing
minimal \(n\geq0\). If \(n>0\), then \(g\notin\pi A\). The nonzero
function \(\bar g\) on the reduced special fiber is nonzero at some
\(f\)-point. Smoothness and henselianity lift that point to an
\(R\)-point, contradicting integrality of \(F\). If the special fiber is
empty, \(\pi\) is a unit in \(A\) and \(n=0\) already. This proves the
criterion. Apply it to finitely many affine coordinate generators; their
relations hold generically and hence in the torsion-free ring \(A\).
Uniqueness follows from that same injection. In particular, two smooth
affine models of the same generic group with the same integral-point
subgroup are uniquely isomorphic through the generic identity.
\(\square\)

\textbf{Proposition (full and connected bounded torus models).} Over
this strictly henselian \(R\), any \(K\)-torus \(T\) has a smooth affine
model \(\mathcal T^b\) whose integral points are exactly the subgroup
\(T(K)^b\) on which every splitting-field character has valuation zero.
There is a smooth affine model \(\mathcal T^0\) with the same generic
fiber and connected special fiber, and \[
 T(K)^0:=\mathcal T^0(R)\ \triangleleft\ T(K)^b,\qquad
 T(K)^b/T(K)^0\simeq\pi_0(\mathcal T^b_f)(f).
 \tag{GB.30}
\] This is a finite group; it is not silently set equal to zero.

\textbf{Proof.} Every bounded subgroup lies in \(T(K)^b\): powers of an
element with a nonzero character valuation are unbounded. Conversely
unit character coordinates and their inverses are bounded in a faithful
splitting-field diagonal representation, so \(T(K)^b\) is bounded.
Choose a faithful \(K\)-representation \(V\). For an initial lattice
\(M\), the module \(\Lambda=\sum_{g\in T(K)^b}gM\) contains \(M\) and is
contained in a fractional lattice. Projection/induction on rank over the
DVR proves that a submodule of a finite free module is finite free:
project to one coordinate, lift a generator of its image ideal, split
off that rank-one summand, and apply induction to the kernel. Thus
\(\Lambda\) is a lattice. Its stabilizer is bounded and contains
\(T(K)^b\), so its intersection with \(T(K)\) is exactly \(T(K)^b\).
Take the schematic closure of \(T\) in \(\operatorname {GL}(\Lambda)\).
Its coordinate ring is the quotient by the kernel of the generic map,
and is torsion-free; it is a finite-type flat affine group model with
precisely those integral points. Smoothen it by the lemma.

Dilatate the resulting smooth model along its special-fiber identity
component. This is the open part consisting of the whole generic fiber
and just that special-fiber component: locally on the component the
ideal is \((\pi)\), while the other closed components are removed. The
affine dilatation description proves that it is affine and of finite
presentation. It is smooth, has connected special fiber, and has exactly
the integral points reducing into the identity component. Each
special-fiber component has an \(f\)-point; smoothness and henselianity
lift it. Reduction therefore gives the exact quotient (GB.30). All
constructions are unique through their integral points by the preceding
lemma. \(\square\)

\textbf{Corollary (any given maximal bounded subgroup).} With the
perfect-residue DVR hypotheses of this section, let \(G/K\) be any
smooth affine algebraic group and let \(P\subset G(K)\) be maximal among
bounded subgroups. Then \(P\) is exactly the integral-point group of a
smooth affine finite-type model of \(G\).

\textbf{Proof.} In a faithful generic representation, the module
\(\Lambda=\sum_{g\in P}gM\) is a finite free invariant lattice by the
projection/induction argument above. The schematic closure of \(G\) in
\(\operatorname {GL}(\Lambda)\) is a flat finite-type affine group
model. Its integral-point group is bounded and contains \(P\), hence
equals \(P\) by the assumed maximality. Group smoothening preserves
exactly these points. The proof requires neither connectedness of \(G\)
nor an assertion that a maximal bounded subgroup exists in every
situation; \(P\)'s maximality is the explicit hypothesis. \(\square\)

\subsubsection{GB.9. The maximal unramified field and its finite
extensions}\label{gb.9.-the-maximal-unramified-field-and-its-finite-extensions}

Put \(K=k^{\mathrm{ur}}\), the union inside a fixed separable closure of
all finite unramified extensions of the original local field \(k\). Its
valuation ring is \(R\), its uniformizer is the original \(\pi\), and
its residue field is the algebraic closure of the finite residue field
of \(k\).

\textbf{Proposition.} \(R\) is a strictly henselian DVR. Every finite
separable extension \(L/K\) has a unique extended valuation, with value
group \(n^{-1}\mathbf Z\), where \(n=[L:K]\); its valuation ring is
finite free of rank \(n\) over \(R\). More precisely \(L=K(\Pi)\), with
\(\Pi\) satisfying an Eisenstein polynomial over some finite unramified
extension of \(k\), and \(\mathcal O_L=R[\Pi]\). This assertion includes
equal-characteristic local fields, for which \(K\) itself is imperfect.

\textbf{Proof.} Every nonzero element of \(R\) lies in a finite
unramified valuation ring and is \(\pi^m\) times a unit. Thus \(R\) is a
DVR with the stated residue field. A polynomial and a simple residue
root involve finitely many coefficients and residue elements. Choose a
finite unramified complete subfield containing all of them; Hensel's
Newton argument there lifts the root. This proves henselianity of \(R\),
and its algebraically closed residue field makes it strictly henselian.
The actual unramified construction and Newton lifting used here are
proved in AB.9.1.

Choose a primitive element \(\alpha\) for \(L/K\); separability gives
one since \(K\) is infinite. Its monic minimal polynomial \(P\) has its
finitely many coefficients in some finite unramified \(E/k\). It is
irreducible over \(E\), because a factorization there would be one over
\(K\). Hence \(M=E(\alpha)\) has degree \(n\) over \(E\), and \(L=MK\).
The finite-extension theorem AB.0 applies to \(M/E\), and gives \(n=ef\)
for its ramification index and residue degree.

Its residue field extension, of degree \(f\), lifts to an unramified
subextension \(E_f/E\) inside \(M\): lift a generator's simple
finite-field minimal polynomial by Hensel's lemma in \(M\), using
AB.9.1. This unramified field is already in \(K\). If \(f>1\), extension
of scalars from \(E\) to \(E_f\subset K\) would lower the degree of
\(\alpha\), contradicting irreducibility of \(P\) over \(K\). Therefore
\(f=1\) and \(e=n\).

Choose a uniformizer \(\Pi\) of \(M\). Its value is \(1/n\) in the
original normalization. Since an extension containing \(\Pi\) has
ramification index divisible by \(n\), its degree over \(E\) is at least
\(n\); hence \(M=E(\Pi)\). The \(n\) conjugates of \(\Pi\) all have
value \(1/n\), by the uniqueness in AB.0. The nonleading coefficients of
its minimal polynomial consequently have positive integral valuation,
while the constant coefficient has valuation exactly one. Thus this
polynomial is Eisenstein.

For every finite unramified \(E'/E\) in \(K\), this same polynomial
remains Eisenstein. It defines a complete degree-\(n\), totally ramified
extension \(ME'/E'\). Its elements have unique expansions
\(\sum_{i=0}^{n-1}c_i\Pi^i\), and the distinct fractional parts of
\(\omega(c_i)+i/n\) give \[
 \omega\Bigl(\sum_{i=0}^{n-1}c_i\Pi^i\Bigr)
   =\min_i\{\omega(c_i)+i/n\}.
 \tag{GB.31}
\] In particular an element is integral if and only if every \(c_i\) is
integral. The valuations of these finite complete extensions agree on
overlaps by AB.0 uniqueness. Their union is \(L\), so they give the
unique valuation on \(L\), and (GB.31) gives its value group and its
valuation ring \(R[\Pi]\). Any other extension restricts to the unique
valuations on all finite complete subextensions, proving uniqueness once
more. The argument uses only finite-field perfectness and separable
finite extensions, and makes no assumption that \(K\) is perfect.
\(\square\)

The same proof works after replacing the original local field by any
finite extension of it. In particular every further finite separable
root-coordinate extension of \(L\) has a finite free valuation ring.
Fractional ideals are finite free over \(R\), with the explicit basis
given by (GB.31). Conjugation preserves the valuation. These are exactly
the algebraic lattice hypotheses in GB.7--GB.8.

\textbf{Conditional geometric extension, with its hypothesis explicit.}
If \(G_K\) is quasi-split, the algebraic constructions GB.0--GB.5 apply
over \(K\): replace compactness of a lattice ball by boundedness, and
use the finite free rings just proved. The two-cell matrix identities,
integral commutators, finite-polynomial Iwahori proof, reflection
arrangement, Coxeter labels, face groups and common apartments do not
use local compactness. The chamber metric and completeness proofs
AB.6--AB.7 allow arbitrary panel multiplicity and therefore construct a
complete CAT(0) enlarged building over \(K\). It is not claimed proper,
since its residue field is infinite. The hypotheses \(G_K\) quasi-split
and the later apartment descent theorem are not inferred from this
arithmetic proposition.

\subsubsection{GB.10. Actual smooth connected facet models over the
unramified
field}\label{gb.10.-actual-smooth-connected-facet-models-over-the-unramified-field}

Assume explicitly in this section that \(G_K\) is quasi-split, where
\(K=k^{\mathrm{ur}}\). Use the building constructed in GB.9. For a facet
\(F\) in an apartment, choose a rational point \(x\) in its relative
interior, including a rational central coordinate. The root levels are
constant there. Write \(K_x\) for the full point stabilizer and \(P_F\)
for the subgroup generated by \(T(K)^b\) and the root levels at \(x\).

\textbf{Theorem.} There are actual smooth affine finite-presentation
\(R\)-models \(\mathcal H_x\) and \(\mathcal G_F\), both with generic
fiber \(G_K\), such that \[
 \mathcal H_x(R)=K_x,\qquad
 \mathcal G_F(R)=P_F^0:=\langle T(K)^0,\ U_{a,-a(x)}\rangle.
 \tag{GB.32}
\] The special fiber of \(\mathcal G_F\) is connected. The subgroup
\(P_F^0\) is independent of the chosen apartment and interior point, and
is normal in the full facet stabilizer. Moreover \[
 K_x/P_F^0\simeq \pi_0((\mathcal H_x)_f)(f),\qquad
 P_F/P_F^0\simeq T(K)^b/T(K)^0.
 \tag{GB.33}
\] The first quotient is finite and includes possible normalizer and
torus components. Neither quotient is declared trivial.

\textbf{Construction of the full model.} Choose a finite Galois
splitting field \(L/K\), with valuation ring \(B\), and the integral
faithful pinned representation of BF.4 over \(B\). View its \(L\)-space
as a \(K\)-space. The induced representation of \(G_K\) is a closed
immersion: after a separable splitting base change, its
restriction-of-scalars target is a product, one projection is the
original faithful immersion, and the diagonal source is its closed
graph. Its weights are \(w_i=\lambda_i(x)\). The weighted norm is
\(\|\sum z_i e_i\|=\max_i\exp(-\omega(z_i)-w_i)\).

For each real \(t\), its ball lattice is \[
 M_t=\{\sum z_i e_i:\omega(z_i)+w_i\geq t\}.
 \tag{GB.34}
\] GB.9 makes this a finite free \(R\)-lattice. There are finitely many
distinct \(M_t\) for \(0\leq t<1\), and \(M_{t+1}=\pi M_t\). Preserving
this finite lattice chain, with inverses, is therefore exactly
preserving the norm. Its automorphism scheme is affine of finite
presentation: in the product of the finitely many
\(\operatorname {GL}(M_t)\), impose that their matrices commute with
each inclusion matrix. Its generic fiber is the common
\(\operatorname {GL}_K\)-space, since all inclusions are generically
invertible. Take the schematic closure of \(G_K\) in it. This is a
finite-type flat affine group model.

Its integral points are precisely the norm-isometries in \(G(K)\), which
are \(K_x\). Indeed the Iwahori fixes that norm and the full Bruhat
partition GB.5 writes every \(g\) as \(i_1ni_2\). Such a \(g\) preserves
the norm if and only if \(n\) does. The weight detector BF.5, with the
full central lattice retained, says precisely that \(n\) fixes \(x\).
Thus \(I N_x I\subset P_xN_x=K_x\), and conversely every defining root
level and every \(N_x\)-element preserves the norm. Smoothen this
closure by GB.8, obtaining \(\mathcal H_x\) with unchanged integral
points.

Dilatate \(\mathcal H_x\) along its special-fiber identity component, as
in GB.8. Call this smooth affine model \(\mathcal G_x\). At first its
integral points are just the kernel of the special-fiber component map;
we now identify that kernel exactly with the subgroup in (GB.32).

\textbf{The connected subgroup is contained in that kernel.} Every root
model \(\mathcal U_{a,-a(x)}\) of GB.7 has all its \(R\)-points in
\(K_x\). Its generic inclusion therefore extends to \(\mathcal H_x\) by
the integral-point criterion GB.8. Its special fiber is connected and
its identity maps to the identity, so its special image lands in the
identity component and the inclusion factors through \(\mathcal G_x\).
The same argument applies to \(\mathcal T^0\). Hence
\(P_F^0\subset\mathcal G_x(R)\).

\textbf{A neighborhood in the connected model has the reverse
inclusion.} In the split pinned group over \(B\), select a principal
open neighborhood \(D(D)\) of the identity inside its integral big cell.
Its complement ideal is stable under torus conjugation. Projecting a
defining function onto its weight-zero summand and rescaling its
identity value gives a function \(D\) of conjugation weight zero with
\(D(e)=1\), whose nonvanishing still implies membership in the big cell.
This selection may be made over \(B\): some projected identity value is
a unit, since their sum is a unit. Each Gaussian root coordinate on this
principal open has a numerator of its corresponding torus weight and a
power of \(D\) as denominator; each Gaussian torus character and its
inverse has weight-zero numerator. Weight projections supply these
homogeneous numerators. The integral faithful closed representation
supplies polynomial lifts in matrix entries and inverse determinant of
those same weights: its coordinate-ring surjection is torus-equivariant,
so projecting any lift to the required weight gives a lift. These are
the exact integral big-cell and grading inputs proved in BF.4--BF.6,
ultimately \href{../AG-RG/AG-RG-03.html}{Roots and reductive groups of
rank one} Lemma 3.1 and \href{../AG-RG/AG-RG-05.html}{Pinnings and the
classification of split reductive groups} Theorem 10.1.

Every matrix monomial of weight \(\beta\) has valuation at least
\(-\beta(x)\) on \(K_x\), by the entry inequalities of the norm.
Weight-zero monomials consequently have integral values. The inverse
determinant has valuation zero because norm isometries and their
inverses have opposite determinant inequalities. Thus \(D\) is integral
on all \(K_x\), and each coordinate numerator has its asserted weighted
lower bound. To apply GB.8 although these functions are \(L\)-valued,
expand them in the finite free integral basis of \(B/R\) from GB.9; all
coefficient functions are \(K\)-regular and integral on
\(\mathcal H_x(R)\), hence extend to \(\mathcal H_x\). Let
\(V\subset\mathcal H_x\) be the principal open where
\(\operatorname {N}_{L/K}(D)\) is a unit. It contains the identity, and
every \(V(R)\)-point has \(D\) a unit in \(B\).

Gaussian decomposition of any \(V(R)\)-point therefore has root
coordinates satisfying \(\omega(c_\beta)+\beta(x)\geq0\). The torus
characters and their inverses have nonnegative valuation, hence
valuation zero. Uniqueness of the decomposition in a root order grouped
into relative-root blocks makes these factors \(K\)-rational. GB.2 and
GB.7 put the root factors in precisely their root models, and the torus
factor in \(T(K)^b\). The generic Gaussian maps from \(V_K\) to these
affine models extend over \(R\) by GB.8, because they have integral
coordinates at every \(V(R)\)-point.

Restrict these maps to the inverse image \(V^0\subset\mathcal G_x\). Its
special fiber is a nonempty open subset of a connected smooth group over
the algebraically closed \(f\), and is irreducible. For completeness, in
a smooth algebraic group the irreducible components are disjoint, since
its regular local rings are domains; translation identifies each with a
coset of the component through the identity. That component is a
subgroup by irreducibility of its product and inverse images. Hence
connectedness implies irreducibility. The special Gaussian torus map on
\(V^0_f\) lands in a single component of \(\mathcal T^b_f\); at the
identity that component is the neutral one. Consequently \[
 V^0(R)\subset P_F^0.
 \tag{GB.35}
\]

For any \(h\in\mathcal G_x(R)\), the two nonempty opens \(V^0_f\) and
\(h^{-1}V^0_f\) intersect. Choose an \(f\)-point in their intersection
and lift it smoothly to \(z\in\mathcal G_x(R)\). Both \(z\) and \(hz\)
belong to \(V^0(R)\), so (GB.35) gives \(h=(hz)z^{-1}\in P_F^0\). This
proves the equality in (GB.32).

\textbf{Components, torus intersection and independence.} Every
component of \((\mathcal H_x)_f\) has an \(f\)-point and that point
lifts smoothly. This proves the first exact quotient in (GB.33). Also
\(T(K)^b\cap P_F^0=T(K)^0\). To see the potentially delicate reverse
containment, take \(t\in T(K)^b\cap\mathcal G_x(R)\), and choose \(z\)
with \(z,tz\in V^0(R)\) by the preceding open-intersection argument.
Their Gaussian torus factors are both in \(T(K)^0\). Left multiplication
by \(t\) conjugates the negative root factor and multiplies the torus
factor by \(t\), while preserving its relative-root order. Thus these
two factors differ by multiplication by \(t\), and \(t\in T(K)^0\).
Since \(P_F\) is generated by \(T(K)^b\) and these same root levels, its
quotient by \(P_F^0\) is exactly the second quotient in (GB.33).

The full stabilizer \(K_x\) is intrinsic to the building. Smooth affine
models with that same integral-point subgroup are uniquely isomorphic by
GB.8, so the full smooth model and its identity-component model are
independent of the chosen weighted representation and apartment. Within
a given facet the root levels and \(T(K)^0\) are constant, so (GB.32) is
independent of the interior rational point too. Rational points exist
because the arrangement has rational discrete levels. Conjugation or a
valuation-preserving field automorphism carries \(K_x\) to \(K_{gx}\) or
\(K_{\sigma x}\); the generic map extends uniquely through these smooth
models by their integral points. It therefore carries the neutral
special component, and hence \(P_F^0\), to the corresponding neutral
component. This proves apartment independence, equivariance and
normality under the full facet stabilizer. We may now write
\(\mathcal G_F\) for \(\mathcal G_x\). \(\square\)

In particular, a Galois-stable facet has a Galois-stable smooth
connected model. Its affine coordinate algebra is an \(R\)-subalgebra of
\(K[G]\); every element has a finite Galois orbit because its finitely
many coefficients lie in a finite unramified field. Thus the semilinear
action is continuous on this algebra, and the explicitly proved affine
descent in AB.9 descends \(\mathcal G_F\) to a smooth connected model
over the original valuation ring. This proves model existence, rather
than assuming it before invoking the descent lemma.

\subsubsection{GB.11. Exact Galois descent and the full-stabilizer
component
cocycle}\label{gb.11.-exact-galois-descent-and-the-full-stabilizer-component-cocycle}

The result here applies to a Galois-stable facet and its models
constructed in GB.10. Galois cohomology uses the usual discrete topology
on algebraic points, so continuous cocycles have finite image. It does
not assert that arbitrary fixed points already lie in descended
apartments.

\textbf{Proposition.} Let
\(\Gamma=\operatorname {Gal}(k^{\mathrm{ur}}/k)\), and let \(F\) be a
\(\Gamma\)-stable facet. Then \[
 H^1_{\mathrm{cts}}(\Gamma,P_F^0)=1.
 \tag{GB.36}
\] For a \(\Gamma\)-fixed point \(x\in F\), put \(D_x=K_x/P_F^0\), with
its actual finite \(\Gamma\)-action. The map \[
 H^1_{\mathrm{cts}}(\Gamma,K_x)\longrightarrow
 H^1_{\mathrm{cts}}(\Gamma,D_x)
 \tag{GB.37}
\] is injective. Consequently a fixed coset \(gK_x\) has a
representative in \(G(k)\) if and only if the component cocycle of
\(g^{-1}\sigma(g)\) is trivial in \(H^1_{\mathrm{cts}}(\Gamma,D_x)\).
The image of the fixed-coset orbit map is precisely \[
 \ker\{H^1_{\mathrm{cts}}(\Gamma,K_x)
                 \longrightarrow H^1_{\mathrm{cts}}(\Gamma,G(K))\}.
 \tag{GB.38}
\] Thus full bounded stabilizers cannot be substituted for connected
facet groups in a Lang argument without this component test.

\textbf{Proof.} GB.10 actually constructs the descended smooth connected
facet model. The integral Lang proof in AB.9 therefore applies and gives
(GB.36), including degrees divisible by the residue characteristic. For
infinite unramified descent a continuous cocycle has finite image and
involves finitely many field coefficients; it is trivial on an open
subgroup. After enlarging a finite unramified field to contain these
coefficients, it is a cocycle over its finite unramified Galois group
and is killed by the finite integral Lang proof. This is exactly the
direct-limit argument already proved in AB.9.

For injectivity, take two \(K_x\)-cocycles whose images have the same
class in \(D_x\). Lift the conjugating \(D_x\)-element to \(K_x\) by the
surjection in (GB.33), and change one cocycle by that cochain. They now
have the same image pointwise. Write them \(c_\sigma\) and
\(c'_\sigma\). The elements \(b_\sigma=c'_\sigma c_\sigma^{-1}\) belong
to \(P_F^0\) and form a cocycle for the twisted action \[
 \sigma * p=c_\sigma\,\sigma(p)\,c_\sigma^{-1}.
 \tag{GB.39}
\] The full \(K_x\) normalizes \(P_F^0\). Its generic conjugations
therefore extend uniquely to the connected model, by the integral-point
criterion GB.8. The cocycle identity gives a semilinear descent datum on
that model. All its coefficients lie in a finite unramified field, as
above. The affine descent proof AB.9 supplies a smooth connected
descended twisted model, and its integral Lang proof kills \(b_\sigma\).
Thus \(c_\sigma\) and \(c'_\sigma\) are cohomologous in \(K_x\), proving
(GB.37). This argument includes possible nontrivial torus and diagram
components; it assumes none of them vanish.

If \(gK_x\) is fixed, then \(c_\sigma=g^{-1}\sigma(g)\in K_x\) is a
cocycle. Changing the representative multiplies it by a
\(K_x\)-coboundary. If its \(D_x\)-class is zero, lift the
\(D_x\)-cochain to \(b\in K_x\); the adjusted cocycle
\(b^{-1}c_\sigma\sigma(b)\) is \(P_F^0\)-valued and, by (GB.36), equals
\(a^{-1}\sigma(a)\) for some \(a\in P_F^0\). Then \(gba^{-1}\) is
Galois-fixed and represents the same coset. Conversely a fixed
representative gives the zero cocycle. This proves the exact component
criterion. Finally any \(K_x\)-cocycle trivial in \(G(K)\) is
\(g^{-1}\sigma(g)\) for a \(g\), and therefore gives a fixed coset. Two
such cosets are in the same \(G(k)\)-orbit exactly when their
\(K_x\)-cocycles are cohomologous, by multiplying their representatives
and checking Galois invariance. This proves (GB.38). \(\square\)

The connected-coset equality \[
 (G(K)/P_F^0)^\Gamma=G(k)/P_F^{0,\Gamma}
 \tag{GB.40}
\] has no component obstruction, by (GB.36) and the same explicit
computation. Equations (GB.37)--(GB.38), rather than (GB.40) with its
denominator replaced, are the correct full-stabilizer statement.

\subsubsection{GB.12. Torus lifting and transporters for the constructed
models}\label{gb.12.-torus-lifting-and-transporters-for-the-constructed-models}

The following proof gives the specific deformation mechanism used here.
It does not invoke Artin approximation to turn a formal homomorphism
into an actual one.

\textbf{Lemma (finite free linearization).} Every smooth affine model in
GB.8 or GB.10 over the strictly henselian \(R\) has a faithful closed
representation on a finite free \(R\)-module. The same holds for its
unramified descended models.

\textbf{Proof.} Its integral-point subgroup is bounded and Zariski dense
in its generic group. Density follows directly from smoothness near the
identity: an étale coordinate chart gives a valuation-open polydisc of
integral points. Such a polydisc is Zariski dense. To verify the latter
assertion, pass to the completion, where the étale inverse is a
convergent local power series. A nonzero regular function has a nonzero
lowest homogeneous term and cannot vanish on an entire polydisc over an
infinite valued field; choose values at which that term is nonzero and
then scale them sufficiently far toward zero. The original field is
dense in its completion, and henselianity supplies the inverses there
already in the original field, so the same conclusion holds before
completion. This also applies to the torus models.

Let \(A\) be the model's coordinate algebra. Take a finite-dimensional
generic right-regular comodule \(V\subset K[G]\) containing a finite
generating set of \(A\); existence is the actual finite-comodule
construction in the earlier \href{../AG-GS/AG-GS-01.html}{Group schemes,
actions and Hopf algebras}, Lemma 5.1. Put \(M=V\cap A\). The
integral-point criterion GB.8 says exactly that \(M\) consists of the
functions in \(V\) integral on the integral-point subgroup. Choose
\(\dim V\) integral group points whose evaluation functionals are
linearly independent; density supplies them by induction. Their inverse
evaluation matrix bounds every coefficient of an element of \(M\). On
the other hand clearing denominators of a basis of \(V\) gives a full
lattice contained in \(M\). Projection/induction over the DVR therefore
makes \(M\) a finite free full lattice.

Right translation by every integral group point preserves \(M\).
Relative to its free basis, the generic representation matrices
consequently have integral values on all integral group points. GB.8
extends their coefficient functions to \(A\), including their
inverse-determinant coordinate. The Hopf identities hold generically and
hence over \(R\). This is an actual representation of the model on
\(M\). It is a closed immersion: evaluating its comultiplication at the
identity expresses every chosen generator of \(A\) as an \(R\)-linear
combination of its matrix coefficients, so the matrix coordinate algebra
surjects onto \(A\).

For an unramified descended model, first base change to the strictly
henselian \(R\), so that the algebraically closed residue hypothesis of
GB.8 applies. Choose the initial generators and generic comodule there
together with their finitely many Galois translates. The lattice
\(M=V\cap A\) is then Galois-stable. The explicit finite-free-module
descent in AB.9 descends it and the representation, and faithful
flatness descends the coordinate surjection. This proves the descended
assertion too. \(\square\)

\textbf{Lemma (lifting a fiber torus).} Let \(\mathcal H\) be one of
these smooth affine models, over \(R\) or over the original henselian
valuation ring. An embedded torus in its special fiber lifts to an
embedded torus in \(\mathcal H\), with the same residue torus.
Transporters and centralizers between embedded tori are smooth. A
prescribed subtorus may be kept fixed when lifting in its smooth
centralizer.

\textbf{Proof of the finite parameter space.} First split the source
torus by a finite unramified cover. In the finite free faithful
representation just constructed, its special-fiber action has a finite
set \(S\) of weights, with positive multiplicities \(r_m\).
Homomorphisms with these weights are represented by finitely many
endomorphism variables \(P_m\) satisfying \[
 P_mP_n=\delta_{m,n}P_m,\qquad \sum_{m\in S}P_m=1,
 \tag{GB.41}
\] and the open-and-closed conditions \(\operatorname {rank}P_m=r_m\).
Their homomorphism is \(t\mapsto\sum_m t^mP_m\). To require its image to
lie in \(\mathcal H\), substitute this matrix into the finitely many
equations of the closed faithful representation; each gives a finite
Laurent polynomial, so vanishing of its coefficients gives finitely many
additional equations. Thus this is a finite-presentation parameter
scheme near the original embedding. The weights generate the character
lattice, since the original torus embedding is faithful. Near its given
point, each positive-rank projector has a unit matrix entry, and matrix
coefficients of \(P_m\rho(t)\) give \(t^m\); coefficients of
\(\rho(t)^{-1}P_m\) give \(t^{-m}\). Hence pullback on the torus
coordinate algebra is surjective, and every homomorphism in this
neighborhood is an embedding. The source torus over the unsplit original
ring is obtained by descending its finite-field character lattice along
the finite unramified cover, using AB.9. That same descent, with the
Galois permutation of \(S\), descends the finite parameter scheme and
its equations.

\textbf{Proof of smoothness.} For a square-zero extension, smoothness of
the affine target first lifts the source's scheme map. Its
multiplication defect is a regular two-cocycle in the Lie algebra
tensored with the square-zero ideal, with action through the original
homomorphism. For a split torus integration means taking the constant
Laurent coefficient. It is invariant under translation in the integrated
variable, over any base ring. If \(c(t,u)\) is a regular two-cocycle and
\(b(t)=\int c(t,v)\,dv\), integration of its cocycle equation gives \[
 c(t,u)=t\cdot b(u)-b(tu)+b(t).
 \tag{GB.42}
\] Thus correcting the map by this one-cochain removes the defect. For a
one-cocycle \(a(t)\), integration of \(a(tu)=a(t)+t\cdot a(u)\) gives
\(a(t)=A-t\cdot A\), where \(A=\int a(u)\,du\). Hence two lifts are
infinitesimally conjugate. The same assertions descend from a split
cover by exactness of the weight-zero projection. A finite free torus
representation acquires no additional weights over a square-zero
extension: its weight projectors reducing to zero are zero by Nakayama.
Thus the corrected map remains in the finite parameter space (GB.41).
The infinitesimal lifting criterion proves that space smooth at the
original embedding.

Henselian lifting in a smooth finite-presentation scheme now gives an
actual homomorphism over the original ring, in the embedding
neighborhood. This follows from an étale coordinate chart and the
simple-root Hensel argument of AB.9.3.

For a transporter point, lift its ambient group element smoothly. The
transported embedding and the required embedding have the same
reduction; the one-cocycle calculation corrects that lift by
infinitesimal conjugation. The transporter equations are
finite-presentation closed equations: expand equality of the two
conjugated torus homomorphisms in their finitely many matrix Laurent
coefficients. If only equality of subgroups is required, a fixed
character-lattice isomorphism gives one such chart, and these charts are
disjoint open pieces indexed by the integral lattice automorphisms. The
same argument gives smoothness of the centralizer, where the two
embeddings are identical. An automorphism of the source torus has no
infinitesimal change, as is immediate from its character lattice. These
arguments establish precisely the transporter and centralizer
assertions. Working in the centralizer of the prescribed subtorus keeps
it fixed. \(\square\)

The primary free source for this deformation problem is Conrad's author
manuscript \emph{Reductive group schemes}, Proposition 2.1.3 and Remark
2.1.4 (PDF pp.~58--60), together with Bruhat--Tits II 5.1.10. The actual
preceding programme proof \href{../AG-RG/AG-RG-01.html}{Tori, maximal
tori and their conjugacy} Lemma 3.1, supplies the weight-cohomology
mechanism. The finite projector parameter space and henselian lifting
above replace the formal-completion/Artin-approximation passage of that
earlier presentation for the models used here.

\subsubsection{GB.13. Galois action, connected residue tori, and stable
apartments}\label{gb.13.-galois-action-connected-residue-tori-and-stable-apartments}

Here \(G\) is defined over the original \(k\), and the explicit
hypothesis \(G_K\) quasi-split is retained. The statement below supplies
stable apartments; it deliberately does not call every fixed affine
subspace a maximal \(k\)-apartment.

\textbf{Lemma (the semilinear action).} The valuation-preserving action
of \(\Gamma\) on \(G(K)\) induces an isometric action on the enlarged
building of GB.9. On every apartment an open subgroup of \(\Gamma\) acts
trivially. Facets and connected models are carried to their
corresponding conjugates.

\textbf{Proof.} A rational Borel subgroup of a quasi-split group is
conjugate to the chosen one by \(G(K)\): apply the absolute Bruhat
charts to its point in the Borel variety. The rational point's absolute
cell label is Galois-fixed, its normalizer representative is rational by
GB.1, and uniqueness of the ordered cell coordinates makes the remaining
unipotent coordinates rational. The same descent argument was used for
actual group elements in GB.5. Within a fixed rational Borel, two
rational maximal tori are conjugate by its unipotent radical. Indeed
projection to its torus quotient identifies each maximal torus with that
quotient, since its multiplicative-type kernel cannot lie in a unipotent
group. Filter the unipotent radical by absolute root height, grouping
Galois-stable relative-root coordinates. Successive quotients are vector
groups; the degree-one torus calculation of GB.12 kills the difference
of the two sections in each quotient, and induction conjugates them. The
assertion about the kernel follows from that same weight projection: a
group of multiplicative type has no nontrivial morphism into a vector
group. For a diagonalizable source this is the degree-one integration
calculation with trivial coefficient action, using its group-algebra
character basis; it works for finite non-smooth sources such as
\(\mu_p\) too. Descent and the filtration give none into the unipotent
radical.

Thus an image of the pinned Borel--torus pair under \(\sigma\) can be
compared with the original pair by an actual \(G(K)\)-conjugation. The
pinning comparison in GB.2 gives precisely the resulting affine change
of root coordinates: the valuation change on each root is its pairing
with the comparison translation. On the central factor use the canonical
character valuation \(\nu_Z(g)(\chi)=-\omega(\chi(g))\), with the same
sign convention as the apartment action; compensate the comparison
conjugation by its central translation. These comparisons carry every
defining root level and normalizer stabilizer to its transformed one.
They therefore preserve exactly the quotient equivalence relation
defining the building.

Independence of the comparison follows by composing two comparisons:
their conjugating elements differ by a torus element, whose
root-coordinate valuation change and entire affine translation are
exactly the torus covariance formula GB.2. The compensating translation
cancels both its root and central contributions. Composition gives the
comparison for \(\sigma\tau\) for the same reason. Hence the quotient
maps form an actual semilinear action. They are Euclidean isometries on
apartments, after averaging the root inner product over the finite
based-root-datum action; they preserve chain lengths, so AB.6's chain
metric makes them global isometries. The coordinates of the chosen pair,
pinning, conjugation and splitting field are finite field data. For an
open Galois subgroup fixing those data the comparison is the identity
and its character action is trivial. Thus it fixes that apartment
pointwise. GB.10's intrinsic integral-point characterization gives the
asserted action on connected facet models. \(\square\)

\textbf{Lemma (the torus rank of a connected facet fiber).} If \(S\) is
the maximal \(K\)-split torus for an apartment through \(F\), its split
integral torus \(S_R\) embeds in \(\mathcal G_F\), and its special image
is a geometrically maximal torus of \((\mathcal G_F)_f\).

\textbf{Proof.} First embed \(S_R\) in \(\mathcal T^0\). Its integral
points are in \(T(K)^b\), so GB.8 extends its generic inclusion to
\(\mathcal T^b\); connectedness of its special fiber factors the map
through \(\mathcal T^0\). This extension is an embedding, including on
the special fiber. For every absolute character \(\chi\) of \(T_L\),
\(\chi\) and \(\chi^{-1}\) have integral values on \(T(K)^b\). Expand
them in the finite free basis of \(B/R\) and apply GB.8; they become
actual characters on \(\mathcal T^b_B\). Their special pullbacks to
\(S_f\) give all its characters, because \(S_L\subset T_L\) is a
subtorus and restriction on character lattices is surjective. This
proves the special embedding. Over \(R\), grade the target coordinate
algebra by translation by \(S_R\). The coefficient ideal of each
character in the pullback to \(R[X^*(S)]\) has full special fiber and
hence is the unit ideal in this local ring. Thus pullback is surjective
in every weight, proving the actual closed immersion. This is also the
explicit graded closed-immersion criterion in
\href{../AG-RG/AG-RG-01.html}{Tori, maximal tori and their conjugacy}.

Any special-fiber torus of \(\mathcal T^0\) lifts by GB.12 to an
embedded torus over \(R\). Strict henselianity splits it, so its generic
rank is at most \(\operatorname {rank}S\) by maximality of \(S\) in
\(T\). Together with the preceding embedding this proves \[
 \operatorname {rank}(\mathcal T^0_f)=\operatorname {rank}S.
 \tag{GB.43}
\]

The Gaussian maps in GB.10 and their product map are mutually inverse on
an open neighborhood of the identity, also on the special fiber: both
compositions are the identity generically and hence on their flat smooth
domains. Consequently the special tangent space decomposes into the
torus-model tangent space and the positive and negative root-model
tangent spaces. Conjugation by \(S_R\) on the latter spaces has the
nonzero characters \(a\) or \(2a\), in the explicit coordinates of GB.7.
These are nonzero group characters even in characteristic two. The
weight-zero tangent space is therefore exactly that of
\(\mathcal T^0_f\).

The centralizer of this split torus in the smooth facet fiber is smooth,
by the transporter proof GB.12. Its identity component consequently has
dimension \(\dim T\). The special torus-model map is injective: any
kernel element maps into the Gaussian neighborhood of the identity,
whose inverse Gaussian torus map recovers it. Its connected commutative
image has dimension \(\dim T\), by that same local inverse. Its map to
the smooth centralizer identity component has invertible differential
everywhere by translation, and trivial kernel; it is an étale
monomorphism and hence an open immersion. Its image is an open subgroup
of a connected group, so it is the entire identity component. A maximal
torus containing the image of \(S_f\) lies in this identity component,
whose torus rank is (GB.43); it cannot properly enlarge \(S_f\). Finally
all maximal tori of a smooth connected affine group over an
algebraically closed field are conjugate, by the preceding programme
\href{../AG-RG/AG-RG-01.html}{Tori, maximal tori and their conjugacy}
Theorem 2.2. This proves the assertion. \(\square\)

\textbf{Theorem (stable apartments through stable facets).} Every
\(\Gamma\)-stable facet \(F\) is contained in an apartment associated to
a maximal \(K\)-split torus defined over \(k\). In particular every
Galois-fixed building point lies in such a stable apartment.

\textbf{Proof.} The descended connected facet model exists by GB.10. Its
special fiber over the finite residue field has a geometrically maximal
torus \(Q\) defined over that field: this is the complete finite-field
Lang proof in actual \href{../AG-RG/AG-RG-01.html}{Tori, maximal tori
and their conjugacy} Theorem 4.2, also supplied in AB.9. Lift \(Q\) to a
torus \(\mathcal Q\) inside the descended smooth facet model using
GB.12. After base change to \(R\), it is split. Its generic torus \(S'\)
is therefore \(K\)-split, is defined over \(k\), and has rank
\(\operatorname {rank}S\) by the preceding lemma. Hence it is maximal
\(K\)-split.

Over the algebraically closed residue field, the special tori \(S_f\)
and \(Q_f\) are conjugate by a special facet-model point. Identify their
split character lattices by that conjugation and extend this
identification to their split \(R\)-tori. Their transporter is smooth by
GB.12; henselian lifting gives a conjugating element
\(g\in\mathcal G_F(R)=P_F^0\). On generic fibers \(gSg^{-1}=S'\). Since
\(P_F^0\) fixes every point of \(F\), the apartment \(gA\) contains
\(F\). The torus \(S'\) is defined over \(k\), so the semilinear action
preserves its apartment. A fixed point's unique open facet is stable,
giving the final assertion. \(\square\)

On a stable apartment the Galois action factors through a finite affine
isometry group, so its fixed affine subspace is nonempty: average the
orbit of any point. Its direction is
\((X_*(S')\otimes\mathbf R)^\Gamma=X_*(S'_k{}_{\mathrm{split}})\otimes\mathbf R\).
The enlarged central fixed factor is likewise the full invariant
character/cocharacter subspace, with its Euclidean metric; it is not
discarded.

These fixed affine subspaces need not all have the maximal \(k\)-split
rank. For example, in a split rank-one group an unramified nonsplit
maximal torus gives a stable apartment with zero-dimensional reduced
fixed subspace. A proof that the fixed set is the \(G(k)\)-union of
apartments for maximal \(k\)-split tori, with the precise descended
facet arrangement and locally finite metric, must therefore still
establish that stronger apartment theorem. Stable apartments alone do
not prove it.

\subsubsection{GB.14. The arithmetic dimension-one input and finite
unramified
definition}\label{gb.14.-the-arithmetic-dimension-one-input-and-finite-unramified-definition}

\textbf{Proposition.} Every finite algebraic extension
\(L/k^{\mathrm{ur}}\) has zero Brauer group. In particular the precise
field condition ``all finite-extension Brauer groups vanish'' is proved
here, in both mixed and equal characteristic.

\textbf{Proof.} First let \(L/K\) be finite separable. In GB.9 we
obtained \(L=MK\), with \(M/E\) a finite totally ramified extension of a
finite unramified local field \(E/k\). This field \(MK\) is
\(M^{\mathrm{ur}}\): the compositum of \(M\) with each unramified
degree-\(d\) extension of \(E\) is unramified of degree \(d\) over
\(M\), as follows either from its Eisenstein basis in GB.9 or from the
identical residue extensions and ramification indices. Conversely the
unramified extension of any given degree is unique by the
finite-field/Hensel construction AB.9.1. Taking the union proves the
equality.

A central simple algebra over \(M^{\mathrm{ur}}\) descends to some
finite unramified local field \(M_d/M\): choose its finite basis and
finitely many multiplication and unit constants there. Associativity
holds at that level. A proper ideal or an extra central element there
would persist after faithful field extension, so the descended algebra
is central simple. The preceding programme
\href{../NT-CFT/brauer-groups-of-local-and-global-fields.html}{Brauer
groups of local and global fields}, Theorem 24.2, proves
\(\operatorname {Br}(M_d)=\mathbf Q/\mathbf Z\), with restriction
multiplying its invariant by the extension degree. Its class is
consequently killed by a further finite unramified extension whose
degree is divisible by the invariant's denominator. That field is still
in \(M^{\mathrm{ur}}\), so the original algebra splits. This proves the
separable case, including every characteristic and every ramification of
\(M/E\). The complete earlier proof of the local theorem, rather than a
citation to a local-class-field textbook, is the provider.

In characteristic \(p>0\), each complete equal-characteristic local
field \(M\) is \(\mathbf F_{q_M}((\Pi))\). The coefficient field is
obtained by lifting the simple roots of \(X^{q_M}-X\); successive
subtraction of the unique coefficient times a power of \(\Pi\), followed
by completeness, gives the unique Laurent series expansion. Its
unramified extensions change only this coefficient field. Grouping
powers of \(\Pi\) by their residues modulo \(p\), and taking the unique
\(p\)-roots of finite-field coefficients, proves \[
 M^{\mathrm{ur}}=\bigoplus_{i=0}^{p-1}
             \Pi^i(M^{\mathrm{ur}})^p.
 \tag{GB.44}
\] Independence follows from the distinct valuation residues modulo
\(p\). Thus this field has \(p\)-degree exactly \(p\). Its finite purely
inseparable extensions are exactly \((M^{\mathrm{ur}})^{1/p^r}\): if the
exponent of such an extension is \(r\), an element of exponent \(r\)
generates degree \(p^r\), which is already the degree of the entire
\(p^r\)-root field by (GB.44). Frobenius identifies that extension
abstractly with \(M^{\mathrm{ur}}\), so it too has zero Brauer group.
Every finite algebraic extension of \(K\) is purely inseparable over a
finite separable one: take sufficiently high \(p\)-powers of a finite
generating set to obtain separable generators. Apply the preceding
separable case and then this Frobenius argument. In characteristic zero
there is no further case. \(\square\)

\textbf{Proposition (finite definition of quasi-splitting).} If
\(G_{k^{\mathrm{ur}}}\) is quasi-split, then \(G\) is quasi-split over
some finite unramified extension of \(k\).

\textbf{Proof.} Its Borel subgroup has a finitely presented coordinate
algebra and finitely many defining equations inside \(G_K\). Descend
those equations, its group identities and its identity/inverse maps to a
finite unramified field \(E/k\). Faithful extension from \(E\) to \(K\)
detects smoothness, geometric connectedness, solvability and the
geometric maximal dimension of that solvable subgroup. Therefore the
descended subgroup is a Borel of \(G_E\). A chosen Borel torus or
pinning likewise descends after enlarging \(E\), since it consists of
finitely many equations and maps. The converse implication is immediate
by field extension. This proves the finite-definition assertion; it does
not assume that \(K\) is perfect. \(\square\)

Over a field whose finite algebraic extensions have zero Brauer group,
GB.16 proves the required torus-cocycle vanishing, GB.17 gives the
rational regular-class representative and GB.18 proves inner-form
descent to a rational Borel. Together with the arithmetic proposition
and finite-definition proposition above, these prove quasi-splitting of
every connected reductive group over a finite unramified extension.
Steinberg, \emph{Regular elements of semi-simple algebraic groups}, and
Borel--Springer, \emph{Rationality properties of linear algebraic groups
II}, Section 8.6, give the historical context, including imperfect
fields.

\subsubsection{GB.15. Fixed points and the descended CAT(0)
metric}\label{gb.15.-fixed-points-and-the-descended-cat0-metric}

For the next metric argument assume \(G_K\) quasi-split; GB.18 proves
this hypothesis for every connected reductive group. The semilinear
action of GB.13 has finite orbit on every point: an apartment is fixed
pointwise by an open subgroup, and a group translate uses finitely many
field coefficients. Since the building of GB.9 is complete CAT(0), the
entire Galois group has a fixed point. Here is the needed existence
proof without assuming properness.

For a finite orbit \(O\), put \(r(z)=\max_{o\in O}d(z,o)\) and
\(r_0=\inf_zr(z)\). Choose \(z_n\) with \(r(z_n)^2\leq r_0^2+1/n\).
Strong convexity AB.7 at the midpoint gives \[
 r((z_n+z_m)/2)^2
 \leq\frac{r(z_n)^2+r(z_m)^2}{2}
                 -\frac{d(z_n,z_m)^2}{4}.
 \tag{GB.45}
\] Its left side is at least \(r_0^2\), so \(d(z_n,z_m)^2\leq2/n+2/m\).
Completeness gives a limit attaining the infimum. The same inequality
proves uniqueness, and invariance of \(O\) makes this circumcenter
Galois-fixed.

The fixed set \(X=\mathcal B_K^\Gamma\) is closed and convex: an
isometry fixing both endpoints fixes their unique geodesic. Thus its
induced metric is complete CAT(0), with the same geodesic distances and
comparison inequalities. By GB.13 every point of \(X\) lies in a stable
apartment, and the induced metric on that apartment's fixed affine
subspace is exactly its Euclidean metric. The central product is \[
 X=(\mathcal B_K^0)^\Gamma\times V_Z^\Gamma.
 \tag{GB.46}
\] The central factor is the entire invariant cocharacter/character
space, hence exactly the split central directions defined over \(k\),
with its retained Euclidean metric. The semilinear central action is
linear by the character-valuation construction in GB.13, so no central
translation is hidden in this equality.

The maximal rational apartment and the rational stabilizers are
established in GB.19--GB.23. In particular GB.22 proves \[
 X=G(k)\,A_k,\qquad A_k=A^\Gamma,
 \tag{GB.47}
\] where the apartment is attached to a maximal \(K\)-split torus whose
maximal \(k\)-split subtorus is maximal in \(G\). GB.19 proves the
unique reduced fixed point and bounded rational subgroup for its
root-zero centralizer, retaining the full split centre. GB.20 gives the
weighted rational root filtration and its exact half-apartment bounds;
GB.21 identifies facet cofaces with parabolics of the connected residue
reductive quotient; GB.22 proves opposite completion and descended wall
exchange. The full-stabilizer component test is GB.11, and GB.23 proves
local compactness, properness, cocompactness and equality of the
descended quotient and inherited metrics. Thus the metric conclusions
above combine with these exact rational apartment and face-group
statements.

\subsubsection{GB.16. Vanishing for every torus over the strict
unramified
field}\label{gb.16.-vanishing-for-every-torus-over-the-strict-unramified-field}

Let \(K=k^{\mathrm{ur}}\), or any finite extension of it. GB.14 proves
\(\operatorname{Br}(E)=0\) for every finite extension \(E/K\), including
the imperfect equal-characteristic case. We prove \[
                         H^1(K,T)=1
                 \quad\hbox{for every \(K\)-torus \(T\)}.       \tag{GB16.1}
\] All cohomology below is discrete. The complete bar resolution and its
Tate extension are proved in the actual programme lesson
\href{../NT-CFT/brauer-groups-of-local-and-global-fields.html}{Brauer
groups of local and global fields}, §3: induced modules are
Tate-acyclic, short exact sequences give long exact sequences,
restriction followed by corestriction is multiplication by the subgroup
index, and all Tate groups are killed by the group order. We supply the
extra tensor argument; vanishing for \(\mathbf G_m\) does not by itself
settle a twisted torus.

\textbf{Finite-group algebra.} Let \(D\) be finite and \(A\) a
\(D\)-module. If \(H^1(H,A)=H^2(H,A)=0\) for every subgroup
\(H\subset D\), then there is a resolution \[
 0\longrightarrow Q\longrightarrow P\longrightarrow A
                      \longrightarrow0                       \tag{GB16.2}
\] with \(P,Q\) projective over \(\mathbf Z[D]\). Consequently
\(A\otimes_{\mathbf Z}\Lambda\), with the diagonal action, is
Tate-acyclic for every subgroup if \(\Lambda\) is a finite free abelian
group with a \(D\)-action. Here are full algebraic details, including
torsion in \(A\).

For a finite \(p\)-group \(D\), the augmentation ideal
\(J\subset\mathbf F_p[D]\) is nilpotent. Induct on \(|D|\): the class
equation gives a nontrivial centre, hence a central element \(z\) of
order \(p\). The central ideal \((z-1)\) has \(p\)-th power zero, and
the quotient is \(\mathbf F_p[D/\langle z\rangle]\). Induction gives
\(J^m\subset(z-1)\), hence \(J^{mp}=0\). Thus a module with zero
coinvariants is zero, even if infinitely generated.

If a module \(V\) killed by \(p\) has \(H_1(D,V)=0\), lift a basis of
\(V/JV\). The resulting map from a free \(\mathbf F_p[D]\)-module \(F\)
is surjective, since its cokernel has zero coinvariants. For its kernel
\(W\), the homology exact sequence gives \[
 0\longrightarrow W/JW\longrightarrow F/JF
                    \longrightarrow V/JV\longrightarrow0 .
\] The last map is an isomorphism, so \(W=0\) and \(V\) is free. If any
single Tate group of \(V\) vanishes, use the usual two exact sequences
with induced modules to shift that degree to \(-2\), keeping all modules
killed by \(p\). The shifted module has \(H_1=0\) and is free. All its
Tate groups vanish; shifting back proves the same for \(V\), and the
argument with \(H_1\) makes \(V\) itself free.

If \(B\) has no \(p\)-torsion and two consecutive Tate groups vanish,
the exact sequence \[
 0\longrightarrow B\xrightarrow{p}B\longrightarrow B/pB
                       \longrightarrow0
\] makes one Tate group of \(B/pB\) zero. Thus \(B/pB\) is free over
\(\mathbf F_p[D]\). Multiplication by \(p\) is now an isomorphism on
every Tate group of \(B\) and of every subgroup. These groups are killed
by powers of \(p\), so all are zero.

A free abelian, cohomologically trivial \(D\)-module \(B\) is
\(\mathbf Z[D]\)-projective. Map
\(\operatorname{Ind}(B)=\mathbf Z[D]\otimes_{\mathbf Z}B\) onto \(B\) by
\(g\otimes b\mapsto gb\), and call the kernel \(B_*\). It is free
abelian and cohomologically trivial. A basis of \(B\) gives an
underlying abelian-group section, so \[
 0\to\operatorname{Hom}_{\mathbf Z}(B,B_*)\to
 \operatorname{Hom}_{\mathbf Z}(B,\operatorname{Ind}(B))\to
 \operatorname{Hom}_{\mathbf Z}(B,B)\to0                       \tag{GB16.3}
\] is exact. Restrict to a Sylow \(p\)-subgroup \(D_p\). We just proved
that \(B/pB\) is \(\mathbf F_p[D_p]\)-free. Since \(B\) is free abelian
and \(B_*\) has no \(p\)-torsion, \[
 \operatorname{Hom}_{\mathbf Z}(B,B_*)/
       p\operatorname{Hom}_{\mathbf Z}(B,B_*)
 =\operatorname{Hom}_{\mathbf F_p}(B/pB,B_*/pB_*).
\] The right module is coinduced: on every free group-algebra summand
evaluate at its group basis and untwist the diagonal action by applying
the inverse group element to its value. Products over the basis index
still have the form \(\operatorname{Maps}(D_p,C)\), because \(D_p\) is
finite. Thus its Tate groups vanish. The integral Hom module has no
\(p\)-torsion, and the preceding \(p\)-multiplication argument makes it
Tate-acyclic. Restriction/corestriction for each Sylow group proves its
cohomological triviality for \(D\) and every subgroup. Taking invariants
in (GB16.3) lifts the identity of \(B\) to an equivariant section of
\(\operatorname{Ind}(B)\to B\). Hence \(B\) is projective.

Subgroups of the possibly infinite free abelian modules here are free:
well-order an ambient basis, intersect the subgroup with successive
initial spans, and at a successor stage project to the new coordinate.
Its image is \(d\mathbf Z\). If nonzero, lift \(d\) and split off its
free cyclic span; the preceding intersection is the kernel. At limit
stages take the unions, since vectors have finite support. The chosen
lifts form a basis. No finite generator list was presupposed.

Finally choose a free \(\mathbf Z[D]\)-module \(P\) surjecting onto
\(A\), with kernel \(Q\). For every \(H\subset D\), the Tate sequence
gives \(H^2(H,Q)=H^1(H,A)=0\) and \(H^3(H,Q)=H^2(H,A)=0\). The module
\(Q\) is free abelian. The no-\(p\)-torsion argument, followed by Sylow
restriction/corestriction, makes it cohomologically trivial. The
projectivity criterion proves (GB16.2). Tensor it with \(\Lambda\);
exactness remains because \(\Lambda\) is free abelian. Projectives
remain projective after this tensor: for a free summand use \[
 \mathbf Z[D]\otimes_{\mathbf Z}\Lambda
 \simeq\mathbf Z[D]\otimes_{\mathbf Z}\Lambda_{\rm underlying},
 \qquad g\otimes v\longmapsto g\otimes g^{-1}v ,
\] with the action on the first factor on the right. Projectives
restrict to projectives and are summands of induced modules. Their Tate
groups vanish, and the exact sequence proves the tensor assertion. This
includes all roots-of-unity torsion in \(A\).

\textbf{Apply the lemma.} Choose a finite Galois splitting field \(L/K\)
of \(T\), with group \(D\). Hilbert90 gives \(H^1(H,L^\times)=0\) for
every \(H\subset D\); its actual programme averaging proof is
\emph{Hilbert's Theorem90 and Kummer theory}, Theorem3.1. The
crossed-product identification of \(H^2(H,L^\times)\) with the relative
Brauer kernel is the actual provider bound in GB.14 and
\href{../NT-CFT/brauer-groups-of-local-and-global-fields.html}{Brauer
groups of local and global fields}, §1. Since
\(\operatorname{Br}(L^H)=0\), degree two also vanishes. With the
cocharacter lattice \(\Lambda=X_*(T)\), the lemma gives \[
       T(L)=\Lambda\otimes_{\mathbf Z}L^\times,\qquad
                         H^1(D,T(L))=0 .
\] Over \(L\), the torus is a product of multiplicative groups, so
Hilbert90 kills its first cohomology. The degree-one
inflation/restriction sequence, proved by the double-bar complex in
\href{../NT-CFT/brauer-groups-of-local-and-global-fields.html}{Brauer
groups of local and global fields}, §1, identifies \(H^1(K,T)\) with
\(H^1(D,T(L))\). One can apply that same calculation over a larger
finite Galois field containing all values of any given continuous
cocycle. This proves (GB16.1).

The integral lattice argument is the Tate-cohomology method discussed in
R. Sharifi, \emph{Group and Galois Cohomology}, Section 1.11, Lemmas
1.11.4--1.11.9 and Theorem 1.11.11. This result holds in mixed
characteristic, in residue characteristic two, and for the imperfect
equal-characteristic strict unramified field. It also proves norm
surjectivity for every finite cyclic \(L/K\): cyclic Tate periodicity
identifies \(K^\times/N_{L/K}L^\times=H^2(D,L^\times)\), whose relative
Brauer group has just vanished.

\subsubsection{GB.17. A rational representative of a strongly regular
class}\label{gb.17.-a-rational-representative-of-a-strongly-regular-class}

We need the following specific part of Steinberg's conjugacy argument.
Over \(K=k^{\mathrm{ur}}\), if \(H\) is quasi-split, semisimple and
simply connected, then every geometrically strongly regular semisimple
conjugacy class invariant under
\(\operatorname{Gal}(K^{\mathrm{sep}}/K)\) has a representative in
\(H(K)\). Invariance here is an assertion about the geometric class, not
a prior assertion that its conjugating element is rational. The proof is
given below. It uses the freely available original Steinberg paper,
§§6--7 and10--11, for the character/section mechanism; the unitary case
is proved directly by a trace form and GB.16. No perfectness of \(K\),
smoothness of a centre, or division by a residue characteristic is
assumed.

\textbf{Integral representations used in the argument.} Let
\(H_{\mathbf Z}\) be the pinned simply connected split group of a given
root system. The actual programme construction is
\href{../AG-RG/AG-RG-05.html}{Pinnings and the classification of split
reductive groups}, Theorems5.1 and10.1, and its §6 supplies the rational
characteristic-zero modules \(V_i=L(\lambda_i)\) of fundamental highest
weights. We explain the needed integral extension of those particular
modules.

Choose a highest vector \(v_i\) and expand \(u^-v_i\), for the universal
ordered negative-root product, as a polynomial in its root coordinates.
Let \(M_i\) be the \(\mathbf Z\)-span of its finitely many coefficients
in \(V_i\). It is a finite free lattice: it is finitely generated and
torsion-free, and it spans \(V_i\) over \(\mathbf Q\), since the highest
vector generates the characteristic-zero module under negative root
operators. Its weight decomposition is integral, because each
coefficient has the weight \(\lambda_i\) minus the corresponding sum of
positive roots. The weight-\(\lambda_i\) part is exactly
\(\mathbf Z v_i\).

The ordered integral multiplication in \(U^-\),
\href{../AG-RG/AG-RG-05.html}{Pinnings and the classification of split
reductive groups} §7, proves stability under every negative root group.
For a positive root group, factor \(x_\alpha(t)u^-\) formally at
\((t,u^-)=(0,1)\) in the integral big cell \(U^-TU^+\). The integral
cell and the formal inverse at its identity are provided by the integral
group of \href{../AG-RG/AG-RG-05.html}{Pinnings and the classification
of split reductive groups}; the root and torus coordinates of that
factorization are power series with integer coefficients. For the
inverse, solve its coordinates degree by degree: its linear coefficient
matrix at the identity is the integral identity, so each new coefficient
is obtained by subtracting an already integral polynomial. Acting on the
highest vector kills \(U^+\). The torus factor acts by the integral
character \(\lambda_i\); expansion of a positive or negative integral
power of a series with constant term one has integral binomial
coefficients. The coefficient of every monomial in \(t\) and the
negative-root parameters consequently lies in \(M_i\). On the
characteristic-zero module the result is already polynomial; comparing
coefficients proves that every coefficient of \(x_\alpha(t)\) preserves
\(M_i\). The torus preserves its weight lattice, and the simple
normalizers, being products of positive and negative root elements with
parameters \(1,-1\), preserve \(M_i\) as well.

These actions extend to \(H_{\mathbf Z}\). On its open cell their
ordered product gives a morphism to \(\operatorname{GL}(M_i)\). Its
determinant and the determinant of the inverse are integral units, since
each root factor is unipotent and the torus factor is diagonal with
integral characters. The homomorphism identity on the open subset where
\(x,y,xy\) all lie in the cell holds over \(\mathbf Q\), hence
integrally: that subset is flat over \(\mathbf Z\), and its regular
functions inject into their rational scalar extension.
\href{../AG-RG/AG-RG-05.html}{Pinnings and the classification of split
reductive groups} Lemma4.1, whose complete translation proof is, extends
the cell map to the entire group. Base change gives a module in every
characteristic with highest weight \(\lambda_i\), a one-dimensional
highest line, and all weights \(\lambda_i-\sum n_\alpha\alpha\),
\(n_\alpha\geq0\). Nothing requires its reduction to be irreducible.
Write \(\chi_i\) for its trace.

The rank-one action on the highest line in the \(i\)-th module is the
two-dimensional standard \(\operatorname{SL}_2\)-action: the pairing
with \(\alpha_i^\vee\) is one; for \(j\ne i\) that pairing is zero. Its
usual matrices over \(\mathbf Q\) preserve the lattice just constructed,
with the two primitive highest-string vectors and coefficients \(1\).
Thus these statements hold after reduction in every characteristic.
Transporting a pinning transports this construction and the traces: a
pinned diagram automorphism permutes \(\lambda_i,M_i,\chi_i\). One can
also see trace uniqueness directly from their characteristic-zero weight
multiplicities and the integral lattice decomposition.

\textbf{Characters separate semisimple classes.} Work first over an
algebraically closed field. On a maximal torus, every \(\chi_i\) is a
Weyl-invariant Laurent polynomial with highest dominant weight
\(\lambda_i\), whose coefficient is one; all other dominant weights are
strictly lower in the positive-root order. Every dominant weight is
\(\lambda=\sum n_i\lambda_i\), \(n_i\geq0\). The product
\(\prod\chi_i^{n_i}\) has highest weight \(\lambda\) with coefficient
one and otherwise lower dominant weights. The sums of the distinct Weyl
translates of a dominant weight form a basis of the Weyl-invariant
Laurent polynomials. Subtracting this highest term and inducting
expresses every such orbit sum as a polynomial in the \(\chi_i\). The
induction terminates: for weights between zero and a fixed dominant
highest weight, pair with \(\rho^\vee\), where
\(\alpha_i(\rho^\vee)=1\); every positive-root subtraction strictly
lowers a bounded discrete nonnegative value on dominant weights. There
are only finitely many lattice points with bounded value in the dominant
cone.

The invariant Laurent polynomials separate finite Weyl orbits in all
characteristics. Given two disjoint finite orbits, interpolate a Laurent
polynomial equal to zero on one and one on the other; the maximal ideals
of distinct torus points are comaximal. Its product over all Weyl
translates is invariant and still has values zero and one. This uses no
division by the Weyl-group order. Finally, semisimple elements are
conjugate into tori by the actual proof
\href{../AG-RG/AG-RG-02.html}{Regular elements and centralizers}
Theorem3.1. If two elements of a torus are conjugate, their torus
centralizers can first be conjugated within the centralizer of the
common element using \href{../AG-RG/AG-RG-01.html}{Tori, maximal tori
and their conjugacy} Theorem2.2; the resulting conjugator normalizes the
torus. Thus semisimple classes are precisely Weyl orbits and are
separated by the \(\chi_i\).

Trace on any module has the same value on an element and its semisimple
Jordan factor: the commuting unipotent factor acts unipotently on each
eigenspace, so its trace is the eigenspace dimension. A fibre of the
tuple \(\chi=(\chi_i)\) whose semisimple class is strongly regular
therefore consists only of that class. Indeed its semisimple factor is
conjugate to the given torus element, and the commuting unipotent factor
lies in its torus centralizer, hence is one. Jordan factors belong to
the group, as in the actual calculation in
\href{../AG-RG/AG-RG-02.html}{Regular elements and centralizers}.

\textbf{A section when simple-root orbits are orthogonal.} Use the
pinned representatives \[
 n_i=x_{\alpha_i}(1)x_{-\alpha_i}(-1)x_{\alpha_i}(1),
 \qquad N(c)=\prod_i x_{\alpha_i}(c_i)n_i .                  \tag{GB17.1}
\] Choose an order putting each Galois orbit of simple roots together.
Suppose roots in each orbit are mutually orthogonal. Its factors
commute, so the morphism (GB17.1) is Galois-equivariant for the
permutation action on its coordinates. These coordinates form the vector
space descended from the permutation root set.

The morphism \(\chi N\) is a polynomial isomorphism. Here is the trace
computation, valid on the integral modules just constructed. On a weight
\(\mu\), the factor \(x_{\alpha_j}(c_j)n_j\) first sends the weight to
\(s_j\mu\) and then adds nonnegative multiples of \(\alpha_j\). A
contribution to a diagonal matrix block of the complete product must
leave each simple-root coordinate unchanged: each \(\alpha_j\) occurs in
just that one factor and the simple roots are independent. Thus each
factor in that contribution returns to the same weight \(\mu\). Such a
contribution is zero unless every
\(\langle\mu,\alpha_j^\vee\rangle\geq0\), so only dominant weights
contribute. For a dominant \(\mu\), its diagonal block depends only on
those \(c_j\) for which \(\langle\mu,\alpha_j^\vee\rangle>0\), because
the positive root action adds exactly this many root steps. On the
highest line of \(M_i\), all factors except the \(i\)-th fix the line;
the standard rank-one matrix gives the diagonal coefficient \(c_i\).

For any other dominant weight \(\mu\) in \(M_i\), write
\(\mu=\sum q_j\lambda_j\). If \(q_j>0\), then \[
 \langle\lambda_j,\rho^\vee\rangle
 \leq\langle\mu,\rho^\vee\rangle
 <\langle\lambda_i,\rho^\vee\rangle .
\] Consequently \[
 \chi_i(N(c))=c_i+F_i\bigl(c_j:
    \langle\lambda_j,\rho^\vee\rangle
       <\langle\lambda_i,\rho^\vee\rangle\bigr),             \tag{GB17.2}
\] with integral coefficients. Solve recursively in increasing
\(\rho^\vee\)-value. This yields an integral polynomial inverse. Since
\(\chi N\) is equivariant, its inverse is equivariant; it descends over
the original field. An invariant tuple of character values thus produces
a rational \(N(c)\). For a strongly regular semisimple class the
preceding fibre argument proves that this is the desired representative.
No descent through an inseparable singleton has been assumed: the
inverse is an actual polynomial morphism.

\textbf{The remaining diagram orbits.} An orbit of adjacent simple roots
in an irreducible finite Dynkin diagram occurs only in type \(A_{2m}\).
Here is a proof requiring no unchecked classification list. Normalize
simple vectors to squared length two. Their Gram matrix has diagonal two
and an edge entry \(-\sqrt{d}\), where \(d\in\{1,2,3\}\) is the product
of the two Cartan integers. Positive definiteness first makes the graph
a tree: a cycle with coefficient one at all its vertices would have
squared norm at most zero. A vertex of degree four is impossible, by
assigning coefficient two to that vertex and one to four neighbours. Two
degree-three vertices are impossible too: choose the path between them,
put coefficient two on its vertices and one on two external neighbours
at either end. The resulting norm is at most zero. Thus there is at most
one branch vertex. There is also at most one multiple edge on a path.
For the segment between two nearest such edges, replace their weights by
\(\sqrt2\) and all intermediate weights by one. Its Gram matrix has a
positive null vector: use coefficient one at the two ends and \(\sqrt2\)
at every intervening vertex. Increasing either edge weight makes its
quadratic value nonpositive, contradicting positive definiteness.

A tree automorphism preserves its centre, obtained by repeatedly
removing leaves. Two adjacent vertices in the same automorphism orbit
must be the two vertices of the central edge: away from it adjacent
vertices have different distances from the centre. Some automorphism
therefore swaps that edge and its two sides. There can be no unique
branch vertex, because this swap fixes no vertex. The diagram is
consequently a path. There can be no multiple edge off the centre, since
it would have a distinct mirrored companion; the central edge is not
multiple either, since its swapped endpoints have equal root lengths.
Thus the path has only simple edges and an even number \(2m\) of
vertices, which is exactly \(A_{2m}\). The pinned matrix realization of
that path as \(\operatorname{SL}_{2m+1}\) follows from its displayed
standard root datum and \href{../AG-RG/AG-RG-05.html}{Pinnings and the
classification of split reductive groups} Theorem5.1, rather than an
algebraic-group classification citation. For reducible data, first
descend a transitive set of simple components to its stabilizer field
\(E\); the original group is its restriction of scalars. Its rational
points and invariant geometric classes are identified with those over
\(E\) by projecting to one component. GB.14 and GB.16 apply to that
finite extension too. Thus it remains to handle the quasi-split simply
connected unitary group of dimension \(2m+1\) over a separable quadratic
\(F/E\). Its identification with the pinned outer \(A_{2m}\)-form is
\href{../AG-RG/AG-RG-06.html}{Automorphisms, forms and parabolic
subgroups} Proposition5.1 together with the explicit transpose-inverse
pinned automorphism and the pinned isomorphism theorem, already bound in
GB.1.

In fact the following argument handles every unitary dimension. Let the
strongly regular class have characteristic polynomial \(P(t)\in F[t]\),
of degree \(n\). It is separable, its constant term is nonzero, the
product of its roots is one, and its multiset of roots is invariant
under \(z\mapsto\bar z^{-1}\). These are exactly the conditions forced
by the special-unitary equation on its geometric class and its
invariance. The coefficient descent assertion is literal: after the
group is identified over \(F\) with \(\operatorname{SL}_n\), the
subgroup fixing \(F\) fixes every characteristic coefficient, and the
remaining involution sends the polynomial to its conjugate reciprocal.
Hence \(P\) is defined over \(F\), with that reciprocal identity, even
when \(E\) is imperfect.

Put \(A=F[t]/P(t)\). It is a finite separable \(F\)-algebra, and the
reciprocal identity defines a semilinear involution \(*\) with
\(t^*=t^{-1}\). The form \[
                    h(x,y)=\operatorname{Tr}_{A/F}(x^*y)    \tag{GB17.3}
\] is hermitian and nondegenerate: after a separable splitting field its
trace pairing is the sum of coordinate products, and \(*\) is an
invertible semilinear map. Multiplication by \(t\) preserves \(h\),
since \(t^*t=1\), and its determinant is one.

Every nondegenerate hermitian form for a separable quadratic \(F/E\)
admits an orthogonal basis with diagonal entries in \(E^\times\), also
in characteristic two. There is a vector \(v\) with \(h(v,v)\ne0\).
Otherwise \(h(x,x)=h(y,y)=0\) and the equality for \(x+cy\), as \(c\)
ranges over \(F\), would give \(\bar c\,h(y,x)+c\,h(x,y)=0\) for every
\(c\). Taking \(c=1\), then a \(c\) not fixed by the involution, forces
\(h(x,y)=0\), contrary to nondegeneracy. Split off this line, whose
orthogonal complement is nondegenerate, and induct. Norm surjectivity
\(F^\times\to E^\times\), proved in GB.16 from the vanished relative
Brauer group, rescales every diagonal entry to one. Hence all
nondegenerate hermitian forms of dimension \(n\) are isometric. In
particular (GB17.3) is isometric to the standard quasi-split form: the
latter is itself a nondegenerate form of that dimension. This statement
does not silently divide by two; the orthogonalization argument works in
characteristic two.

In such an isometry, multiplication by \(t\) becomes a matrix in the
standard \(\operatorname{SU}_n(E)\) with characteristic polynomial
\(P\). Over an algebraic closure its distinct eigenvalues identify its
\(\operatorname{GL}_n\)-class. The same class is a single
\(\operatorname{SL}_n\)-class: adjust the determinant of a conjugator by
a diagonal centralizer element in an eigenbasis. Thus it is a
representative of the specified geometric special-unitary class. This
proves the missing type \(A_{2m}\) case without invoking the
perfect-field semisimple-part construction of a cross-section.

Products and restrictions of scalars now prove the stated
rational-representative result for every quasi-split semisimple simply
connected \(H/K\). The construction used exact integral coordinates and
separable trace forms. In particular it includes wildly ramified
quadratic extensions in characteristic two and groups with nonsmooth
finite centre.

\subsubsection{\texorpdfstring{GB.18. Quasi-splitting over
\(k^{\mathrm{ur}}\), including nonsmooth
centres}{GB.18. Quasi-splitting over k\^{}\{\textbackslash mathrm\{ur\}\}, including nonsmooth centres}}\label{gb.18.-quasi-splitting-over-kmathrmur-including-nonsmooth-centres}

For every connected reductive \(G/k\), \[
                 G_{k^{\mathrm{ur}}}\text{ is quasi-split}.    \tag{GB18.1}
\] Its Borel is defined over some finite unramified extension. We prove
the assertion without citing the dimension-one theorem as a proof.

Put \(K=k^{\mathrm{ur}}\). First let \(H/K\) be quasi-split, semisimple
and adjoint, and let \(c_\sigma\in H(K^{\mathrm{sep}})\) be a continuous
cocycle. Its simply connected covering \(\widetilde H\) is quasi-split
too: the inverse image of its Borel is a Borel, as is checked in the
central-quotient root coordinates of \href{../AG-RG/AG-RG-04.html}{Root
data, Weyl chambers and the Bruhat decomposition} Theorem7.1. The
adjoint group acts algebraically on \(\widetilde H\) by conjugation. The
central kernel acts trivially, so conjugation descends faithfully flat
through the central quotient; this is also the group-scheme automorphism
decomposition of \href{../AG-RG/AG-RG-06.html}{Automorphisms, forms and
parabolic subgroups} Theorem3.1. Thus \(c\) twists \(\widetilde H\) to a
reductive \(K\)-group \({}_c\widetilde H\), even if that kernel is
nonsmooth. We do not lift the individual cocycle values through it.

A smooth connected affine group over \(K\) has Zariski-dense rational
points, by the local henselian argument proved in GB.12: take an étale
coordinate chart at the identity, spread its finitely many equations
over a finite local-field subextension, and Hensel-lift a sufficiently
small valuation polydisc. In the completion, the étale inverse expands a
nonzero regular function into a nonzero convergent power series: the
local completion is faithfully flat and its étale coordinates are formal
parameters. Choose parameter values at which its first nonzero
homogeneous polynomial is nonzero, possible by the one-variable root
bound and induction over the infinite field. Scaling those values
sufficiently toward zero makes that term dominate the later terms.
Henselianity places the inverse points in the original field, dense in
its completion. Thus no nonzero regular function vanishes throughout the
polydisc. This proves density in the chart and hence the group. It
applies to the twisted group and does not assume that \(K\) is perfect.

The geometrically strongly regular semisimple locus of \(\widetilde H\)
is nonempty and open. In a torus avoid \(\alpha(t)=1\) for every root.
The \href{../AG-RG/AG-RG-02.html}{Regular elements and centralizers} §3
root-weight proof makes the centralizer smooth with identity component
that torus. Every other component normalizes it because the identity
component is characteristic. Avoid also the finitely many proper loci
\(w(t)=t\), \(w\ne1\), in the faithful Weyl action. No other component
remains. The differential of conjugation is surjective on this locus, by
\href{../AG-RG/AG-RG-02.html}{Regular elements and centralizers}, so its
image is open. It is invariant under automorphisms and therefore
descends to the twist. Density supplies a strongly regular
\(y\in{}_c\widetilde H(K)\). On identifying the two groups over
\(K^{\mathrm{sep}}\), its twisted rationality is \[
                     c_\sigma\sigma(y)c_\sigma^{-1}=y,       \tag{GB18.2}
\] where the formula means the algebraic adjoint action. The geometric
conjugacy class in the original quasi-split \(\widetilde H\) is
invariant. GB.17 gives a representative \(y_0\in\widetilde H(K)\) of
that class.

There is a conjugator \(a\in H(K^{\mathrm{sep}})\) with
\(y=ay_0a^{-1}\). Here geometric conjugacy alone would not justify the
field of definition. Its transporter in the adjoint group is a torsor
under \(T=C_H(y_0)\). This is a \(K\)-torus: over an algebraic closure
the root factors \(\alpha(y_0)-1\) are units and absence of a Weyl
stabilizer excludes additional components, while the diagonalizable
fixed-point argument makes the centralizer scheme smooth. The geometric
transporter is consequently nonempty and smooth over
\(K^{\mathrm{sep}}\). A nonempty smooth finite-type scheme over a
separably closed field has a point: take an étale chart to affine space,
choose a rational point of its nonempty open image, and a point in its
étale fibre has a finite separable residue field, hence the same field.
This proves the choice of \(a\) without a purely inseparable adjunction.

Substitution in (GB18.2) shows that \(a^{-1}c_\sigma\sigma(a)\) is a
cocycle in \(T(K^{\mathrm{sep}})\). GB.16 kills this class, so \(c\) is
trivial. We have proved \(H^1(K,H)=1\) for every quasi-split semisimple
adjoint \(H\), with nonsmooth covering centres treated explicitly.

For arbitrary connected reductive \(G/K\), choose a split pinned form
\(G_0\) of its root datum. The
\href{../AG-RG/AG-RG-06.html}{Automorphisms, forms and parabolic
subgroups} Theorem3.1 proves, as group schemes, \[
 \operatorname{Aut}(G_0)=
 G_0^{\mathrm{ad}}\rtimes\operatorname{Aut}(\text{based datum}).
\] The form is split over a finite separable extension: a geometrically
maximal torus spreads through its smooth transporter, the torus
character lattice has a finite separable splitting field, and the split
root lines can then be pinned using the
\href{../AG-RG/AG-RG-03.html}{Roots and reductive groups of rank one}
root construction. Its descent cocycle has an outer datum part
\(e_\sigma\) and an inner part \(c_\sigma\). The datum part has finite
image even when the full integral-lattice automorphism group is
infinite, since the descent is defined over that finite splitting field.
Twisting by \(e\) gives a quasi-split \(G_{\rm qs}\): its pinned Borel
is preserved and descends, as proved in
\href{../AG-RG/AG-RG-06.html}{Automorphisms, forms and parabolic
subgroups} Proposition5.1. The remaining \(c\) lies in the quasi-split
semisimple adjoint group \(G_{\rm qs}^{\rm ad}\). Its first cohomology
has just been proved trivial. Correcting the descent isomorphism by the
coboundary identifies \(G\) with \(G_{\rm qs}\) and transports its
rational Borel. This proves (GB18.1).

The Borel and its inclusion use finitely many coefficients in \(K\), so
are already defined over a finite unramified extension \(E/k\).
Smoothness, connectedness, solvability and the geometric Borel dimension
descend faithfully flat from \(K\) to \(E\). Equivalently its point in
the projective Borel scheme uses finitely many such coordinates. This
proves the finite-unramified assertion with all the original arithmetic
hypotheses. In particular imperfect equal characteristic, arbitrary
residue characteristic, the full reductive centre and nonsmooth finite
centres have been retained.

\subsubsection{\texorpdfstring{GB.19. The centralizer of a maximal
\(k\)-split
torus}{GB.19. The centralizer of a maximal k-split torus}}\label{gb.19.-the-centralizer-of-a-maximal-k-split-torus}

This section closes the root-zero and boundedness part of the descent
problem. Let \(S\subset G\) be a maximal \(k\)-split torus and put
\(M=Z_G(S)\). Then \(M\) is connected reductive by the actual
torus-centralizer proof \href{../AG-RG/AG-RG-02.html}{Regular elements
and centralizers} Theorems2.1--2.2. The torus \(S\) is central in \(M\);
there is no larger \(k\)-split torus in \(M\), because a larger one in
\(M\) would also be one in \(G\). More explicitly any noncentral
\(k\)-split torus \(S_1\subset M\) commutes with \(S\), and their
product is a split torus of dimension greater than \(\dim S\). Thus all
\(k\)-split directions of \(M\) are exactly those of its central \(S\).

By GB.18, \(M_K\), \(K=k^{\mathrm{ur}}\), is quasi-split. GB.0--GB.6 and
AB.6--AB.7 give its complete enlarged building \(\mathcal B(M,K)\), its
root coordinates, and its product
\(\mathcal B(M,K)^0\times V_{Z(M),K}\). GB.13 supplies its natural
Galois isometric action and GB.15 supplies a fixed point. We prove \[
          \mathcal B(M,K)^\Gamma
                    =\{z_M\}\times V_{Z(M),K}^\Gamma,
       \qquad \dim V_{Z(M),K}^\Gamma=\dim S .                \tag{GB19.1}
\] The point \(z_M\) is unique in the reduced building. The second
equality is the torus-lattice calculation AB.10, applied to the full
centre; anisotropic central directions have no invariant cocharacters.

Let \(x\) be a fixed point, and \(F\) its open facet. The connected
model \(\mathcal M_F\) constructed in GB.10 descends to
\(\mathcal O_k\), by GB.13 and AB.9. The central split torus
\(S_{\mathcal O_k}\) embeds in it: in a \(K\)-apartment every central
torus character and its inverse have value zero on the bounded central
torus, so the character-extension and faithful weight argument of GB.13
embeds its integral model in the bounded torus model and then the facet
model. It is central there, since it is central generically and the
models and source are flat.

Every geometric maximal torus of the special fibre contains this central
\(S_{\mathbf F_q}\): the product with that central torus is a torus, so
maximality forces inclusion. Choose a maximal torus \(Q\) defined over
\(\mathbf F_q\). Its existence is the actual finite-field Lang proof in
\href{../AG-RG/AG-RG-01.html}{Tori, maximal tori and their conjugacy}
§4, already bound in GB.13. In using that proof one may start with a
torus containing the central \(S\), and all conjugations retain \(S\).
Lift \(Q\) to an unramified torus \(\mathcal Q\subset\mathcal M_F\) by
GB.12, prescribing the already embedded subtorus \(S_{\mathcal O_k}\).
The prescribed-subtorus part follows from the same finite-projector and
degree-one transporter calculation there; centrality ensures that its
centralizer is the whole group. Hence this is an actual lift containing
\(S\), not merely a formal lift with an unspecified embedding of that
subtorus.

The maximal split part of an unramified torus has the same rank
generically and on its finite residue field: both are the invariant
cocharacters of its finite étale character lattice. Consequently the
generic maximal \(k\)-split part of \(\mathcal Q_k\) contains \(S\).
Maximality of \(S\) in \(M\) makes it exactly \(S\). GB.13's
special-fibre torus-conjugacy and smooth transporter argument identifies
\(\mathcal Q_K\) with a maximal \(K\)-split torus of \(M_K\), and gives
a \(\Gamma\)-stable apartment \(A_x\) containing the whole facet \(F\).
In this paragraph \(\mathcal Q_K\) means the maximal \(K\)-split part of
the generic unramified torus; its generic rank equals the geometric
torus rank of the facet model, as proved in GB.13. That apartment has \[
                       A_x^\Gamma=x+V_S.                   \tag{GB19.2}
\] Indeed its invariant cocharacters are exactly the just identified
split part \(S\), and its affine Galois action has the fixed point
\(x\). Central directions in the enlarged building are global, so this
is a translate of the same central vector space at every \(x\).

To deduce uniqueness of the reduced fixed point, take two fixed points
\(x,y\). They have a common \(K\)-apartment by GB.5 and AB.4, and their
segment is fixed by convexity. That segment meets finitely many facets:
in that apartment the affine root hyperplanes form the locally finite
arrangement proved in GB.4, and only finitely many meet a bounded
segment. Every facet whose open part contains a fixed segment point is
\(\Gamma\)-stable, by uniqueness of the open facet of a point. For each
such \(F'\), repeat the preceding construction. Its stable apartment
contains the whole \(F'\), and (GB19.2) shows that \(F'^\Gamma\) has
only central directions. Partition the segment into its finitely many
facet pieces; on each piece its reduced-building projection is constant.
Continuity at their endpoints makes this projection constant on the
entire segment. Thus \(x,y\) have the same reduced projection.
Conversely one fixed point plus all invariant central directions is
fixed, by the product action in GB.13--GB.15. This proves (GB19.1). This
centralizer argument uses only the stable-apartment and centralizer
statements already established; the general residue/coface and
maximal-apartment results are proved separately in GB.21--GB.22.

The rational group \(M(k)\) preserves its unique reduced fixed point. On
the remaining \(S\)-flat it acts by translations. Its central character
valuation is \[
 \langle\chi,\nu_M(m)\rangle=-\omega(\chi(m)),
                       \qquad \chi\in X_k^*(M),             \tag{GB19.3}
\] with the same sign as BF.2, GB.2 and AB.10. These characters span the
dual of the split central factor: the root datum of \(M/M^{\rm der}\)
identifies them with a finite-index sublattice of the centre's rational
character space. An isometry commuting with all central
\(S(k)\)-translations has identity linear part on this flat, so (GB19.3)
gives its full action. The image is discrete because each displayed
valuation is integral; it contains a full lattice from the cocharacters
\(S(k)\), by the nonsingular character/cocharacter pairing. Thus it is
itself a full lattice.

The kernel of (GB19.3) is compact. Its elements fix an enlarged point,
hence belong to the full norm stabilizer constructed in GB.10. In the
faithful integral representation over a finite splitting extension, the
displayed weighted norm bounds every matrix coordinate and inverse
coordinate. To transfer those bounds to any faithful
\(k\)-representation, express its matrix entries and inverse entries as
regular functions in the first closed matrix immersion. Each is a
polynomial in finitely many matrix and inverse coordinates with fixed
coefficients in that splitting field, so has a uniform valuation bound.
Its \(k\)-values consequently lie in closed bounded subsets of the
original locally compact field. The same holds for its inverse matrices.
The subgroup is closed, and its image in the affine closed immersion
over \(k\) is a closed subset of a compact product of valuation balls.
Hence it is compact. This uses compactness of the original \(k\), not of
\(K\). In particular an anisotropic connected reductive \(k\)-group has
compact rational points and a unique fixed point in the enlarged
\(K\)-building. With a split centre present, that entire centre has
instead been retained in (GB19.1).

For completeness the centralizer building embeds in that of \(G_K\) with
the same apartment metric. Choose a maximal \(K\)-split torus of \(M_K\)
containing \(S\). It is maximal in \(G_K\): its centralizer in \(G_K\)
is contained in \(M_K\), since it contains \(S\), and is the torus
centralizer already computed for quasi-split \(M_K\). Its apartment is
the same full cocharacter space. Restrict the integral root and weight
representation of \(G_K\) to \(M_K\). The \(M\)-root coordinates are
exactly the roots vanishing on \(S\), as follows from the actual
centralizer weight calculation and big-cell root product. The full norm
stabilizer at a point therefore satisfies \[
                K_x(G,K)\cap M(K)=K_x(M,K).                 \tag{GB19.4}
\] One may use the restricted faithful representation in GB.10: its
integral root coefficients and weight coordinates are still those
already used, and it detects a normalizer's movement of every apartment
coordinate. Thus the quotient map \(M(K)\times A\to G(K)\times A\) is
well defined and injective after the point-stabilizer relations. Two
points of the \(M\)-building have a common \(M\)-apartment, on which
this map is the identity Euclidean map; hence it is isometric. Its image
is the union of apartments whose maximal \(K\)-split tori contain \(S\),
since such tori lie in \(M\) and are conjugate there by the quasi-split
root/Borel argument of GB.13. The image is \(\Gamma\)-invariant and its
fixed set is the single full \(S\)-flat (GB19.1). This gives an actual
maximal-split flat in the general fixed-point space, together with its
exact root-zero compact kernel.

\subsubsection{GB.20. Rational filtered root parameters and exact
half-apartment
bounds}\label{gb.20.-rational-filtered-root-parameters-and-exact-half-apartment-bounds}

Let \(S\subset G\) be maximal \(k\)-split, and choose the
\(\Gamma\)-stable \(K\)-apartment \(A\) through the centralizer flat of
GB.19. Its maximal \(K\)-split torus is defined over \(k\) and contains
\(S\); choose \(x_0\in A^\Gamma\). For a primitive character
\(b\in X^*(S)\) on a nonzero root-weight ray, let \(\Psi_b\) be the set
of nondivisible \(K\)-roots \(a\) with \[
                              a|_S=m_a b,\qquad m_a>0 .
\] The integers \(m_a\) are the actual restrictions of the characters.
Internal \(2a\) coordinates in a nonreduced \(K\)-root group have weight
\(2m_ab\). Define \(U[b]_K\) as the product of these \(K\)-root groups,
in any order refining positive \(S\)-weight. It is a smooth connected
unipotent subgroup and descends to \(k\). To verify descent without an
assumed nonsplit root classification, choose a cocharacter of \(S\)
positive on \(b\); the ambient positive group has ordered root
coordinates. Conjugation by \(S\) assigns each coordinate its character.
A semilinear automorphism fixing \(S\) sends an output coordinate to a
polynomial of the same \(S\)-weight. On the ray subgroup all its inputs
have weights \(m b\), \(m>0\), so every output of a weight off that ray
is zero. Hence the subgroup is \(\Gamma\)-stable. Positive closure of
the ray and the integral root collection formula give its group law.
This construction allows every actual positive integer multiple; it does
not assume an unproved claim that all original nonsplit rank-one groups
are \(\operatorname{SL}_2\) or \(\operatorname{SU}_3\).

Normalize the \(K\)-root valuations at \(x_0\): \[
 \phi_{a,x_0}(u)=\phi_a(u)+a(x_0),\qquad
 U[b]_r(K)=
      \prod_{a\in\Psi_b}\{u\in U_a(K):
                         \phi_{a,x_0}(u)\geq m_a r\}.        \tag{GB20.1}
\] The same definition includes the \(2a\) coordinates at twice that
bound. The coordinate product is a subgroup by the weighted commutator
calculation GB.2. Root order changes preserve it, because every new
collected coordinate has the sum of the input weights and at least the
sum of their valuation bounds. The \(x_0\)-normalization makes it
\(\Gamma\)-stable: GB.13's affine covariance, evaluated at the fixed
\(x_0\), removes the translation part. Define \(U[b]_{r+}\) by strict
inequalities on the intrinsic root valuations, equivalently
\(U[b]_{r+\epsilon}\) for a sufficiently small positive \(\epsilon\).
There is such a common \(\epsilon\), since only finitely many coordinate
weights occur and all their valuation sets are discrete.

The exact integral models of these groups are products of the
finite-free additive and unitary parameter spaces of GB.7, with the
scaled bounds (GB20.1). Their group law and inverse are the integral
collected polynomials, so they are smooth affine models with connected
fibres. They are canonically \(\Gamma\)-stable. Indeed the generic
semilinear maps preserve their entire \(R\)-point groups. The
integral-point criterion of GB.8 extends these maps and their inverses,
since \(R=\mathcal O_K\) is henselian with algebraically closed residue
field; the group and action identities extend by flatness. Infinite
continuous descent is effective by AB.9, after putting the finitely many
equations and character weights over a finite unramified subextension.

Here is the cohomology needed for rational lifting: \[
                   H^1_{\rm cts}(\Gamma,U[b]_r(R))=1,
       \qquad H^1_{\rm cts}(\Gamma,U[b]_{r+}(R))=1.          \tag{GB20.2}
\] Here \(U[b]_r(R)\) denotes the point group of its integral model,
namely the bounded group in (GB20.1); no separate integral-point
assumption on a \(K\)-coordinate is intended. Filter this group by
increasing positive \(S\)-character weight. In coordinates the subgroup
with all weights less than \(j b\) zero is closed and normal, and its
quotient by the next such subgroup is a vector group on a finite free
\(R\)-module. The commutator of weights \(i b,j b\) has weight
\((i+j)b\), so these quotients are central at the relevant stage. A
semilinear \(S\)-equivariant automorphism on a quotient of the one
weight \(j b\) is linear: every monomial of degree different from one
has a different algebraic character, even when the integer \(j\) is
divisible by the residue characteristic. Differentiating the character
is neither necessary nor sufficient for this assertion.

The positive cohomology of every such semilinear \(R\)-module is zero.
AB.9 descends it to a finite free \(\mathcal O_k\)-module. A given
continuous cochain and the finite descent equations lie over a finite
unramified \(E/k\). The ring \(\mathcal O_E\) has a normal integral
basis: choose a normal basis in its separable residue extension,
Hensel-lift its generator, and the determinant of its conjugates is a
unit by residue independence. Thus
\(\mathcal O_E\otimes_{\mathcal O_k}N\), under
\(\operatorname{Gal}(E/k)\), is an induced module for every coefficient
module \(N\). The bar insertion contraction proves vanishing in all
positive degrees, without dividing by the extension degree. Increasing
\(E\) to contain all values of a continuous cocycle proves the same for
\(\Gamma\). This proof works for finite-length modules too: the same
normal basis identifies their scalar extensions with induced modules,
regardless of \(\mathcal O_k\)-torsion.

Now kill a group cocycle in successive central vector quotients. After
making its first vector image a coboundary, correct by a lift of that
vector; the remaining cocycle lies in the next group. Each vector
correction lifts, because the actual coordinates give surjective maps on
\(R\)-points. The finitely many weights make the induction terminate.
This proves (GB20.2). The same proof applies to the strict model and, if
required, to an inner twist of these actions by a cocycle: inner
conjugation changes a parameter only by terms of larger weight, so its
action on each vector quotient is unchanged.

The boundary map on rational points is consequently surjective: \[
 U[b]_r(k)\longrightarrow
     \bigl(U[b]_r(K)/U[b]_{r+}(K)\bigr)^\Gamma .             \tag{GB20.3}
\] Indeed lift a fixed boundary class to \(u\in U[b]_r(K)\). Its defect
\(u^{-1}\sigma(u)\) is a cocycle in the strict group. Write it
\(a^{-1}\sigma(a)\) using (GB20.2). Then \(ua^{-1}\) is fixed, remains
in the same bounded group, and has the chosen boundary class. The same
argument applies to quotients by a positive-weight tail, such as the
\(2b\)-tail when that is the only higher root multiple. Thus there is no
unproved averaging of nonlinear unitary coordinates and no failure at
residue characteristic two.

The boundary group itself has a finite central filtration with
successive finite-length \(R\)-modules. To see this, filter by its
coordinate character weights as above. At a fixed weight the change from
non-strict to strict bounds is inclusion of fractional lattices; its
quotient is a finite-length \(R\)-module. Changes of unitary coordinates
may add quadratic terms, but those have twice the character weight and
disappear in the preceding vector quotient. Its semilinear action
descends to the corresponding finite residue modules by AB.9 and the
normal-basis calculation. In particular a nonzero vector boundary piece
has nonzero rational points after descent to \(\mathbf F_q\), including
\(q=2\). Apply (GB20.3) to lift them.

Finally these numerical bounds really give the half-apartment fixed by
each parameter. For \(u=\prod u_a\in U[b](k)\), put \[
                 \phi_b(u)=
                   \min_{u_a\ne1}
                           \frac{\phi_{a,x_0}(u_a)}{m_a}.
                                                               \tag{GB20.4}
\] This minimum is independent of the ordered coordinates: reordering
adds terms with summed weights and no smaller weighted valuation.
Applying the inverse reordering gives equality of the two minima. In the
unitary coordinate group use the exact BF.7 formula
\(\phi_a(u,v)=\omega(v)/2\) and \(\omega(u)\geq\omega(v)/2\), with no
lost \(2a\) parameter. The filtered model and its bounds, rather than
the chosen trace-one coordinate change, define (GB20.4).

The full point-stabilizer calculation GB.10 gives \[
 u\text{ fixes }x\in A^\Gamma
     \quad\Longleftrightarrow\quad
                 b(x-x_0)+\phi_b(u)\geq0 .                  \tag{GB20.5}
\] For this equivalence, the root Gauss coordinates of \(u\) have torus
component one and denominator one. Membership in the full weighted norm
stabilizer is therefore exactly the root-coordinate inequalities
\(\phi_{a,x_0}(u_a)+a(x-x_0)\geq0\), with the doubled coordinate bounds
in the unitary group. Substituting \(a|_{A^\Gamma}=m_a b\) proves
(GB20.5). The sign is the same as the split formula; replacing it by its
negative would interchange the fixed and moving half-apartments. These
levels, rational boundary lifting and all same-sign weighted commutator
bounds are now proved for the original arbitrary connected reductive
group.

The ray construction does not identify an original nonsplit rank-one
group with a classical matrix group. GB.21 proves the exact residue wall
data, and GB.22 proves the opposite rational completion and wall
exchange at every attained wall using (GB20.3). All actual restricted
character weights and higher tails remain in the ordered ray groups;
that proof route needs no additional standalone classification of the
original nonsplit forms.

\subsubsection{GB.21. The residue group and its coface
parabolics}\label{gb.21.-the-residue-group-and-its-coface-parabolics}

Let \(F\) be a \(K\)-facet, let \(\mathcal H_F\) be the connected model
of GB.10, and write \(H_F\) for its special fibre over \(\bar f\), the
algebraically closed residue field. Its maximal torus is the special
torus \(S_K\) of GB.13. Its centralizer is the connected bounded torus
special fibre \(T_F\), with reductive torus quotient \(S_K\). We prove
that reduction identifies the cofaces of \(F\), in reverse order, with
the parabolics of \[
                              \overline H_F=H_F/R_u(H_F).
                                                               \tag{GB21.1}
\] It identifies cofaces in a given apartment with parabolics containing
that apartment's special maximal torus. This is the concrete input for
maximal-apartment descent. The primary free comparison is
BTII4.6.8--4.6.14 and4.6.33--4.6.35; the root and point-group arguments
needed here follow.

Fix \(x\in F\). A root coordinate is called strict if its weighted value
at \(x\) is positive. These strict subgroups do not depend on
\(x\in F\). Non-strict boundary coordinates occur precisely at the
affine root walls containing \(F\); the opposite coordinate has a
boundary too, by the exact normalizer completion and valuation
covariance GB.0--GB.2. For a reduced field root, its boundary quotient
is one additive \(\bar f\)-line: consecutive fractional ideals in the
totally ramified root field have the same algebraically closed residue
field by GB.9. We give the corresponding calculation for every ramified
unitary wall, rather than assume its residue quotient.

Normalize the valuation on its intermediate field \(E\) to integers. A
quadratic \(F_1/E\) has an Eisenstein uniformizer \(\theta\), with \[
 \theta^2-t\theta+n=0,\qquad \omega_E(n)=1,\quad
 \omega_E(t)\geq1,\quad
 \mathcal O_{F_1}=\mathcal O_E\oplus\mathcal O_E\theta .
                                                               \tag{GB21.2}
\] Here \(\omega_{F_1}(\theta)=1/2\); in characteristic two separability
says \(t\ne0\). This is the finite-extension theorem AB.0/GB.9. For
\(s=2r\in\frac12\mathbf Z\), the lattice \(I_s\) of the \(v\)-coordinate
is \[
 \begin{array}{ll}
 s=j:& I_s=\pi^j\mathcal O_E\oplus\pi^j\mathcal O_E\theta,\\
 s=j+\frac12:&
 I_s=\pi^{j+1}\mathcal O_E\oplus\pi^j\mathcal O_E\theta .
 \end{array}
\] Its trace ideal is generated by the displayed basis traces
\(2\pi^j,t\pi^j\), or \(2\pi^{j+1},t\pi^j\), respectively. Put
\(h_2=\omega_E(2)\), with \(h_2=\infty\) in characteristic two, and
\(h_t=\omega_E(t)\), with \(h_t=\infty\) if \(t=0\). Choose a basis
element \(c\) whose trace generates this ideal; if both traces have
equal valuation, choose the element with the larger coordinate
valuation. Then \(\lambda=c/\operatorname{Tr}(c)\) has trace one. The
formula of GB.7 is \[
 u\in\{u:\omega_E(Nu)\geq\tau\},\qquad
 v=-\lambda Nu+\delta z,\qquad
 \tau=\omega_E(\operatorname{Tr}I_s),
                                                               \tag{GB21.3}
\] where \(\delta\mathcal O_E=I_s\cap\ker\operatorname{Tr}\). In
characteristic two this trace-zero line is \(E\). Otherwise use
\(\theta-t/2\); its valuation is \(1/2\) if \(h_t\geq h_2+1\), and the
integer \(h_t-h_2\) if \(h_t\leq h_2\). These alternatives follow from
unequal valuations of its two terms, and include \(t=0\).

Exactly one scalar parameter survives modulo the strict group. At
\(s=j\), it is the leading \(u\)-coordinate when \(h_t\geq h_2+1\), and
the leading trace-zero coordinate when \(h_t\leq h_2\). At \(s=j+1/2\),
the choices are reversed. To check this list, the minimum valuation of
\(-\lambda Nu\) is \(\omega(c)\), because \(\omega(Nu)=\tau\) is
attained and \(u\) ranges over the principal ideal of valuation
\(\tau/2\). The minimum trace-zero valuation has the integer or
half-integer part just computed. In each case one minimum equals \(s\)
and the other is strictly larger. The surviving principal-ideal quotient
has dimension one over \(\bar f\). In characteristic two \(h_2=\infty\):
at integral \(s\) the trace-zero parameter survives, and at
half-integral \(s\) the \(u\)-parameter survives. Thus no missing
division by two occurs at these walls.

The surviving coordinate is additive. If it is \(u\), this follows from
the first coordinate of the exact group law; if it is trace-zero \(v\),
the product term \(-\bar u u'\) has strictly larger valuation and
disappears. The strict model maps onto the kernel of that one coordinate
in the special root model; (GB21.3) and the explicit polynomial group
law prove the assertion on schemes, not just their field points. The
torus character of the surviving line is respectively \(a\) or \(2a\).
The opposite line has its negative character, since the boundary
normalizer conjugates the two models and fixes the wall. Rescaling
\(\omega_E\) back to the original \(k\)-normalization divides all these
valuations by the same ramification index and preserves the list. This
proves the boundary-line assertion for every quadratic ramification and
every residue characteristic.

Let \(R_F\) be the reduced connected closed subgroup generated by all
special strict root subgroups and \(R_u(T_F)\). It is normal in \(H_F\).
Same-sign normalization is the integral commutator collection: one
strict input makes the resulting weighted value positive. For opposite
coordinates, one strict input makes the rank-one Gauss denominator a
unit congruent to one in its effective boundary coordinate; BF.7/GB.3's
Gauss identities put its root corrections in the strict groups and its
torus correction in \(R_u(T_F)\). These identities use the original
unitary formula, so the list above includes both effective \(a\) and
\(2a\) walls. The bounded torus normalizes these subgroups. The wall
normalizers preserve strict values and normalize \(R_u(T_F)\). Finally
reduction of the exact group \(P_F^0\) is surjective, by henselian
smooth lifting, and \(P_F^0\) is generated by the root groups and
\(T^0\), by GB.10. Thus the normalization checks on these generators
prove normality in the entire special group.

This normal subgroup is unipotent. Use the faithful lattice-chain
representation from GB.10 and its associated graded representation.
Every strict root matrix raises the weighted filtration and is trivial
on its associated graded. The special maximal torus acts faithfully
there: the weights of a closed faithful torus immersion generate its
character lattice, as in GB.13. The reduced identity kernel of this
graded representation has no torus. Indeed every torus of \(H_F\) is
conjugate into its maximal torus, and normality of the kernel would put
a conjugate of any kernel torus in that faithful torus. The no-torus
argument \href{../AG-RG/AG-RG-02.html}{Regular elements and
centralizers} §1 makes its identity kernel unipotent. The image of
\(R_u(T_F)\) is central in the graded image, since its commutators with
boundary root lines are strict: a unipotent group has no nontrivial
homomorphism to the scalar character group of a boundary line. Thus the
connected subgroup generated by the strict kernel and this
torus-centralizer radical is an extension of unipotent groups, hence
unipotent (triangularize its normal kernel and then its quotient). This
proves \(R_F\subset R_u(H_F)\).

Pass to the smooth affine quotient \(Q=H_F/R_F\). Here quotient
existence has a short local proof. Choose a finite-dimensional
translation comodule in \(\bar f[H_F]\) containing generators of the
ideal of \(R_F\), using the actual finite-comodule argument AG-GS
\href{../AG-GS/AG-GS-01.html}{Group schemes, actions and Hopf algebras},
Lemma5.1. The top exterior power of its intersection with that ideal
gives a line whose scheme stabilizer is \(R_F\): stability of the
subspace is exactly stability of its generated ideal, and translation
preserves \(R_F\) exactly for its own points. The normal unipotent
\(R_F\) acts trivially on that line, since it has no characters, and
normality makes it fix the span of all its translates. The resulting
representation has kernel exactly \(R_F\). Its reduced image is a closed
affine group: a constructible subgroup contains an open of its closure,
and translations make it the whole closure. The homomorphism onto it is
flat by generic flatness and translations; its fibres are translates of
the smooth \(R_F\), so it is a smooth surjection and an \(R_F\)-torsor.
This represents the quotient and makes \(Q\) smooth connected affine.

The special Gauss cell shows that \(R_F\) has precisely the
strict-coordinate and \(R_u(T_F)\) directions. Its ordered product has
that dimension. Conversely every word in these generators, whenever in
the Gauss cell, has zero effective boundary coordinates, by the strict
commutator and opposite-Gauss calculations above; the same equations
hold on their reduced closure. Hence its intersection with that cell has
no additional directions. Its scheme intersection with each boundary
additive line is trivial: that line's effective coordinate is one of the
displayed linear \(u\)- or trace-zero coordinates, while this coordinate
vanishes on \(R_F\). It follows that \(Q\) has a torus \(S_K\),
one-dimensional boundary root lines, and no other moving tangent
directions. Its Gauss cell is the product of those boundary lines and
that torus.

There is no further unipotent radical in \(Q\). If its smooth normal
radical had a nonzero boundary weight, centralize the codimension-one
kernel of that weight and take its positive limit group. The actual
smooth cocharacter construction \href{../AG-RG/AG-RG-03.html}{Roots and
reductive groups of rank one} Lemma3.1, isolates a smooth connected
one-dimensional subgroup with that tangent line; in \(Q\) the only
positive weight on this ray is its one boundary line. The equal
dimensions and closed inclusion force that whole additive line into the
radical. Normality and the boundary normalizer put the opposite line
into it too. But the exact GB.0 normalizer formula expresses the wall
representative as a product of those two lines modulo strict
coordinates. It would then be in the radical while acting by a
nontrivial reflection on the surviving torus. This is impossible: a
normal unipotent radical has trivial scheme intersection with a torus,
by diagonalizing multiplicative-type representations and triangularizing
unipotent ones, so that torus embeds in the radical quotient, where a
radical element acts trivially. This also excludes an infinitesimal
boundary intersection: a nontrivial torus-stable finite subgroup of an
additive line has nonzero tangent, giving the same positive-limit
contradiction. Weight zero in a putative radical lies in the centralizer
of the torus, which in \(Q\) is the torus itself, by the Gauss
coordinates. It therefore vanishes too. A smooth connected
positive-dimensional group has nonzero Lie algebra, so the radical is
trivial. Thus \(Q\) is reductive and \(R_F=R_u(H_F)\). The quotient
(GB21.1) is therefore reductive. Its roots are precisely the surviving
boundary characters, one of \(a,2a\) for each wall direction, with the
negative character on the opposite line. Its maximal torus is the image
of \(S_K\), and its Weyl group is the finite wall group \(W_F\): every
root reflection has the reduced boundary representative, and these
generate its Weyl group by the \href{../AG-RG/AG-RG-04.html}{Root data,
Weyl chambers and the Bruhat decomposition} root/Weyl proof. Their
actions are exactly the linear actions of the walls through \(F\). This
identifies the local spherical root data without replacing ramified
unitary roots by a reduced-field formula.

Let \(F'\) be a coface of \(F\) in the base apartment and \(d\) a small
rational displacement from \(F\) into \(F'\). A boundary root at \(F\)
keeps its full level if its character is nonnegative on \(d\), and
becomes strict if negative. Previously strict roots stay strict: choose
the displacement smaller than every positive gap in the finitely many
discrete root levels. The maps \(\mathcal H_{F'}\to\mathcal H_F\) exist
by the integral-point extension criterion GB.8, because
\(P_{F'}^0\subset P_F^0\). Their special image contains \(R_F\) and the
torus, and its reductive image is generated by the boundary lines
nonnegative on \(d\). It is exactly the cocharacter parabolic
\(P_{\overline H_F}(d)\), by the actual classification proof
\href{../AG-RG/AG-RG-06.html}{Automorphisms, forms and parabolic
subgroups} Theorem1.1. All parabolics containing this torus occur by the
same finite root-chamber classification. Different cofaces give
different sign/zero patterns.

The exact inverse-image assertion is \[
 P_{F'}^0=
 \{g\in P_F^0:
          \operatorname{red}(g)\in
        \operatorname{preimage}(P_{\overline H_F}(d))(\bar f)\}.
                                                               \tag{GB21.4}
\] The reduction kernel is contained in every such \(P_{F'}^0\): in the
Gauss coordinates of an element reducing to one, each root parameter is
strict, and the torus parameter is in \(T^0\). The desired parabolic's
points are generated by its torus, its permitted boundary root groups
and the radical. Their parameters lift in the displayed root models, and
their torus parameters lift by smooth henselian lifting. They
consequently lift to \(P_{F'}^0\). This proves both inclusions in
(GB21.4).

Finally all cofaces can be moved into the base apartment by \(P_F^0\).
Choose a chamber of their common \(K\)-apartment having \(F\) in its
closure. The finite-wall Bruhat decomposition of GB.5/AB.4 makes \(P_F\)
transitive on the chambers containing \(F\). Apartments containing a
chamber are its Iwahori translates, by the explicit apartment argument
AB.6. The factors \(T_b\) normalize the base apartment and can be
absorbed, leaving \(P_F^0\); GB.10 gives \(P_F=T_bP_F^0\). Thus the
coface orbits are \(P_F^0/P_{F'}^0\). By (GB21.4) these are exactly the
residue flag orbits. The latter exhaust the parabolics by
\href{../AG-RG/AG-RG-06.html}{Automorphisms, forms and parabolic
subgroups} Theorem1.1, and equal parabolics have equal stabilizers. This
proves the claimed bijection, its order, and the apartment assertion.

When \(F\) is \(\Gamma\)-stable all maps and radicals are canonical and
descend. A stable coface therefore corresponds exactly to an
\(\mathbf F_q\)-parabolic of the descended reductive quotient. The
parabolic scheme and its faithful-flat descent are proved in
\href{../AG-RG/AG-RG-06.html}{Automorphisms, forms and parabolic
subgroups} Theorem2.1. Its finite set of \(\mathbf F_q\)-points gives
finitely many stable cofaces. Smooth henselian lifting gives
\(P_F^0(k)\twoheadrightarrow H_F(\mathbf F_q)\). Lang for the connected
special group, and for its connected parabolic preimages, makes this
connected rational point group transitive on the cofaces of any rational
parabolic type. This is a statement about the connected model. No
assertion that \(H^1\) of the full disconnected bounded stabilizer
vanishes has been used.

\subsubsection{GB.22. Maximal rational apartments and descended wall
exchange}\label{gb.22.-maximal-rational-apartments-and-descended-wall-exchange}

Let \(X=\mathcal B(G,K)^\Gamma\), with its closed convex metric from
GB.15. We prove that its apartments are the maximal-split flats of GB.19
and that one such flat \(A_k\) satisfies \[
                              X=G(k)A_k .                   \tag{GB22.1}
\] We use the coface theorem just proved, not an unproved equality of
fixed cosets for a disconnected stabilizer.

We first record the finite-field torus fact needed in this argument. In
a connected reductive group over \(\mathbf F_q\), every maximal split
torus has torus centralizer and all maximal split tori are rationally
conjugate. Such a group is quasi-split: apply the pinned outer/inner
decomposition \href{../AG-RG/AG-RG-06.html}{Automorphisms, forms and
parabolic subgroups}, and kill its adjoint inner cocycle by the actual
finite-field Lang proof \href{../AG-RG/AG-RG-01.html}{Tori, maximal tori
and their conjugacy} §4. If the centralizer of a maximal split torus had
a nontrivial semisimple derived factor, that factor would be quasi-split
by the same argument and would have a nontrivial split cocharacter: sum
the coroots in a diagram orbit as in GB.1. Multiplying by the original
central split torus would enlarge it. Its centralizer is therefore a
torus. A generic cocharacter in the split torus now determines a
rational Borel, since no absolute root is zero on it. Rational Borels
are conjugate: their smooth connected stabilizer has trivial \(H^1\) by
Lang, so the rational points of the Borel variety are one rational
orbit. Within a rational Borel, the root-height vector filtration and
additive Hilbert90 conjugate the maximal tori by its unipotent radical,
as in the \href{../AG-RG/AG-RG-06.html}{Automorphisms, forms and
parabolic subgroups} Proposition5.1 proof. This proves the split-torus
conjugacy assertion. For a smooth connected nonreductive group over
\(\mathbf F_q\), lift through its split unipotent radical. Tori have
trivial intersection with it, their transporters are smooth by GB.12,
and the vector-filtration/Lang calculation makes the same conjugacy
rational. These arguments also show that a maximal split torus belongs
to every rational parabolic after a rational conjugation.

At a fixed point \(x\), with stable open facet \(F\), choose a maximal
\(\mathbf F_q\)-split torus \(Q_0\) of \(H_F\), then a rational
geometric maximal torus \(Q\) containing it. GB.12 lifts them, with the
inclusion prescribed, to actual unramified tori in the descended
\(\mathcal H_F\). GB.13 gives a stable \(K\)-apartment containing \(F\),
whose generic maximal \(k\)-split part is \(S_0\), the lift of \(Q_0\).
We claim that \(S_0\) is maximal split in \(G\).

The centralizer-model comparison used in this claim is exact. For any
integral split subtorus \(\mathcal S_0\subset\mathcal H_F\), the smooth
centralizer \(C_{\mathcal H_F}(\mathcal S_0)\) has connected fibres, by
the \href{../AG-RG/AG-RG-02.html}{Regular elements and centralizers}
Lemma1.1 and its weight proof. Its generic fibre is \(M_0=Z_G(S_0)\).
Its Gauss coordinates are the bounded torus and precisely the root
coordinates of zero \(S_0\)-weight. Thus its integral points are
\(P_F^0\cap M_0(K)\), which is the connected facet point group of
\(M_0\): the torus model is the same \(T^0\), and root coordinates of
nonzero \(S_0\)-weight vanish in the centralizer. GB.10 and the integral
uniqueness criterion GB.8 identify it with the actual \(M_0\)-facet
model. Its special reductive quotient is the centralizer of the image of
\(Q_0\) in \(\overline H_F\). Indeed the same zero-weight Gauss
coordinates remain after killing the strict radical of GB.21. This is a
torus by the finite-field fact above. This proves the comparison on
models and point groups; it does not assume that centralizers commute
with arbitrary nonsmooth reduction.

The \(M_0\)-building embeds in the \(G_K\)-building by the root-zero
restriction argument at the end of GB.19; that embedding argument
applies to any split subtorus, irrespective of maximality. The point
\(x\) belongs to it. Since the residue reductive group of its
\(x\)-facet is a torus, GB.21 gives no proper stable coface there. A
sufficiently small ball about \(x\) in the \(M_0\)-building meets only
cofaces of this facet, by the explicit local-star radius of AB.7. Its
fixed part is therefore locally contained in the facet's fixed affine
space. That space has precisely the directions of \(S_0\): choose a
special maximal torus containing \(Q_0\) and apply GB.13; every
additional invariant cocharacter would enlarge the maximal residue split
torus. Since \(S_0\) is central in \(M_0\), these directions are a
global central metric factor. After removing it, the fixed set has an
isolated point \(x\). The fixed set is convex and connected; a geodesic
from an isolated point cannot reach a second point. Hence its entire
fixed set is this one \(S_0\)-flat.

If \(S_0\) were contained in a larger split torus, enlarge to one
maximal split torus \(S_1\) in \(M_0\). GB.19 applied to
\(Z_{M_0}(S_1)\), and its isometric embedding, gives an \(S_1\)-flat in
the same fixed set of the \(M_0\)-building. Its dimension is larger, a
contradiction. This proves the claim without importing rational
conjugacy of maximal split tori as a premise. In particular every fixed
point belongs to a stable apartment whose fixed affine space is a
maximal-split flat.

The fixed cells are the nonempty intersections \(F^\Gamma\) of stable
\(K\)-facets with the fixed set. GB.13 puts each in a stable apartment,
so it is a convex polyhedral cell. Its cofaces are exactly the rational
parabolics of GB.21. They are finite in number. Maximal cofaces
correspond to rational Borels, because the finite residue quotient is
quasi-split, as just proved. Their fixed dimensions agree locally: all
those Borels contain conjugate maximal residue split tori, and the lift
has the same generic split rank. These locally equal dimensions agree
throughout \(X\). To check this last assertion without a purity
assumption, take a segment between any two points in a common
\(K\)-apartment. Its finitely many facet pieces have overlapping local
stars at their endpoints; the maximal-coface dimensions therefore agree
successively along it. Write the resulting constant dimension as \(r\),
including the invariant centre. Every maximal \(k\)-split flat has
dimension \(r\): its local cells have that dimension, and every
residue-maximal lift is a maximal split torus, while a split torus of
largest dimension contributes its own flat by GB.19. Thus no hidden
lower-rank apartment has been substituted.

The maximal cells are gallery connected, after dividing out the full
central Euclidean factor. A segment meets finitely many cells. In each
local star the rational parabolics form the finite spherical building of
the residue group, by GB.21 and the actual finite Bruhat theorem; its
maximal chambers are gallery connected (write any finite Weyl label as a
word in simple reflections). Replace the finitely many star crossings of
the segment by these galleries. This gives a gallery between any two
maximal cells, also when the segment originally crossed a face of larger
codimension. In rank one it is the succession of vertex stars; in rank
zero the reduced fixed space is a point.

The group \(P_F^0(k)\) is transitive on the maximal cofaces at every
face \(F^\Gamma\), by the final paragraph of GB.21. Moving across a
gallery thus makes \(G(k)\) transitive on maximal fixed cells. Choose
one of them \(C_k\), in a maximal-split flat \(A_k\). Every point lies
in the closure of a maximal cell, so it lies in some \(G(k)\)-translate
of \(\overline C_k\), proving (GB22.1).

We also need apartment transitivity with a fixed cell. Let two
maximal-split flats contain a fixed point \(x\). Their integral split
tori embed in \(\mathcal H_F\): their units fix \(F^\Gamma\), and the
weighted root-character and integral extension argument of GB.13, over
\(R\), gives their actual split integral embeddings; on \(F\), take the
minimal \(K\)-facet containing \(x\), whose norm inequalities those
units also satisfy. Their special split tori are maximal in \(H_F\), by
the claim and its residue-rank argument. Conjugate them by
\(H_F(\mathbf F_q)\) using the finite-field fact. The transporter of the
two prescribed split torus embeddings is smooth by GB.12, so that
conjugator lifts to \(P_F^0(k)\). Its generic conjugation carries the
two split tori, hence carries their unique centralizer flats by GB.19.
It fixes \(x\). When both flats contain a maximal cell, use its model
and the same transporter; the lifted element fixes that cell pointwise.
Thus the action is strongly transitive on a flat together with a maximal
cell in it.

The walls in \(A_k\) are exactly the nonzero restrictions of the
\(K\)-walls that meet it. They are locally finite, by the finite
discrete root-value sets of GB.4. The cell containing a given open fixed
point is the intersection with its unique \(K\)-facet, so this
restricted arrangement is the actual fixed-cell arrangement. Its
noncentral normal directions span: a cocharacter annihilating all
restricted roots is central in \(G\), by the absolute big-cell
centralizer computation. Translation by \(S(k)\) supplies a full
lattice. Its maximal cells are bounded modulo the \(k\)-split centre of
\(G\): each root direction has walls at an unbounded arithmetic sequence
in both signs, so their finite independent directions bound a cell.

Every wall reflection lifts to \(N_G(S)(k)\). At a relative panel the
residue group has relative semisimple rank one. Its relative normalizer,
computed from the finite-field orbit classification GB.1, contains its
reflection. Lift that residue normalizer point preserving the split
torus. First correct its lift in the special group's unipotent radical
so that it normalizes the chosen split torus (the split-torus
transporter and Lang argument above); then use the smooth normalizer of
GB.12 and henselian lifting. The resulting
\(n\in P_F^0(k)\cap N_G(S)(k)\) fixes the panel pointwise and acts on
\(A_k\) by its nontrivial rank-one linear reflection. An affine isometry
with those properties is precisely the wall reflection, with no
translation along the panel. It preserves the full arrangement because
it acts on \(X\). AB.2's proved gallery-disc argument now gives the
affine Coxeter system for this arrangement and the alcove \(C_k\). This
proof does not identify its translation lattice with a coroot lattice
that may be inappropriate for the original group.

Here is the group exchange, including the actual opposite-parameter
input. Let \(I_k\) be the full pointwise stabilizer of \(C_k\), let
\(\alpha=b+r\geq0\) be one of its walls, and let \(s\) be the reflection
lift. Use \(L=Z_G((\ker b)^0)\), with its split central tangential torus
removed; its rational split semisimple rank is one. Its positive and
negative groups are the ray groups \(U[b],U[-b]\) of GB.20. The exact
centralizer-model comparison above applies to its panel model. Its
special reductive quotient has finite-field relative rank one.
GB.0--GB.1, with the zero valuation and the finite root fields, prove
its two Bruhat cells, including the unitary finite-field cell.

Put \(A=U[b]_r(k)\) and \(A^+=U[b]_{r+}(k)\). The boundary lifting GB.20
makes \(A\) surject onto the positive residue root group. The coface
inverse-image formula (GB21.4) consequently gives \[
               I_k=J A,\qquad
               J=I_k\cap s^{-1}I_ks .
\] The reduction kernel and the residue torus are in \(J\), and the
rank-one Borel is its torus times its positive root group. For the full
\(I_k\), first use transitivity of the connected point group on flats
containing \(C_k\): an element fixing \(C_k\) can be adjusted by its
connected facet group to normalize \(A_k\), where it belongs to the
pointwise kernel \(M(k)^0\). Thus \(I_k\) is its connected coface group
times \(M(k)^0\). This compact kernel is normal under \(s\) and lies in
\(J\). Hence the same factorization holds for the full \(I_k\),
including its component classes; they have not been discarded.

For a nontrivial residue class \(u\in A\setminus A^+\), the residue
element \(su\) is in the opposite Gauss cell. This cell has an actual
integral inverse. Apply \href{../AG-RG/AG-RG-03.html}{Roots and
reductive groups of rank one} Lemma3.1 to the panel model of \(L\) and
the split cocharacter normal to the panel: its positive, zero and
negative weight models are precisely the filtered positive ray, the
\(M=Z_G(S)\) model and the filtered negative ray. Their multiplication
is an open immersion. Its special image is exactly the inverse image of
the residue opposite Gauss cell. Here is the radical check in this
assertion. The smooth connected radical over the perfect residue field
has a torus-stable central filtration with additive-line quotients:
triangularize its semidirect product with the acting torus and take the
reduced identity components of the successive upper-entry kernels, as in
the full solvable-structure proof \href{../AG-RG/AG-RG-01.html}{Tori,
maximal tori and their conjugacy} §2. Each quotient is an additive line
with one algebraic torus character. On it, multiplication of its
negative-, zero- and positive-weight limit subgroups is an isomorphism,
since exactly one of those three groups is that line. Induct up the
central filtration: a factorization in the quotient lifts, and its
remaining central coordinate is put in its uniquely prescribed sign
factor; uniqueness follows by projecting and then using the central
line. This proves that the corresponding three factors exhaust the
radical, with polynomial inverses. Their formation is the actual
cocharacter limit construction and commutes with reduction by
\href{../AG-RG/AG-RG-03.html}{Roots and reductive groups of rank one}
Lemma3.1. Thus killing the radical sends the integral cell to the
residue Gauss cell and the sign-factor argument lifts its entire inverse
image. The strict-coordinate calculations of GB.21 identify these
factors with the displayed positive, zero and negative filtered models,
including the unitary higher-coordinate tails. Thus \(su\), whose
reduction is in that cell, lies in the integral open cell, and its
inverse coordinates give \[
                 su=v\,m\,w,\qquad
 v\in U[b]_r(k),\quad m\in M(k)^0,\quad
                            w\in U[-b]_{-r}(k).             \tag{GB22.2}
\] Here \(M(k)^0=\ker\nu_M\): the middle factor fixes the panel point,
hence its centralizer-flat translation is zero by GB.19. The coordinates
are rational because the open-cell inverse is defined over
\(\mathcal O_k\); an external Bruhat assertion has not replaced this
factorization.

Equation (GB22.2) is also the normalizer completion. Rearranging gives
\[
 u=(s^{-1}vs)(s^{-1}m)w ;
\] the outside factors are in the negative ray, while \(s^{-1}m\)
normalizes \(S\) and acts by the wall reflection. Thus every nonzero
boundary parameter has its exact opposite rational completion. Apply
this same panel-model argument at each attainable wall to obtain the
assertion at every level; it does not require different wall levels to
be conjugate under \(S(k)\). All levels and half-apartment signs are
those of (GB20.4)--(GB20.5). Higher central parameters in a nonreduced
ray remain in these groups throughout.

For \(n\) in the apartment normalizer, \(J\) can be moved through \(s\).
A strict \(u\in A^+\) conjugates through \(s\) into \(I_k\). For a
boundary \(u\), if \(n^{-1}\alpha\) is positive on \(C_k\), (GB20.5)
gives \(n^{-1}un\in I_k\). If it is negative, use (GB22.2):
\(v,m\in I_k\), and \(n^{-1}wn\in I_k\) by that same exact
half-apartment inequality. These two cases prove \[
                sI_kn\subset I_ksnI_k\ \cup\ I_knI_k .      \tag{GB22.3}
\] The group nontriviality required by AB.4 is also exact. The positive
boundary group at this panel is nonzero and has a nonzero
\(\mathbf F_q\)-point, by GB.20 and its finite-field rank-one
identification. Lift it to \(u\in U[b]_r(k)\setminus U[b]_{r+}(k)\). It
belongs to \(I_k\), by (GB21.4), and its exact inequality (GB20.5) fixes
the base side of the wall and fails on the interior of \(sC_k\). Thus
\(u\notin sI_ks^{-1}\), and these two subgroups differ. The intersection
of \(I_k\) with the apartment normalizer is its pointwise kernel,
because \(C_k\) has an open subset of the apartment.

Finally galleries and panel transitivity generate \(G(k)\) by \(I_k\)
and the normalizer. Each adjacent chamber is an \(I_k\)-translate of
\(sC_k\) at the base panel. Move this argument along a gallery. A
residual element stabilizing \(C_k\) can be adjusted by \(I_k\) to
normalize \(A_k\), using the just proved transitivity on flats
containing that cell. Thus no extra rational-group generation premise is
needed. Apply AB.4 to the affine Coxeter subgroup, as follows. Let
\(W_{\rm aff}\) be the reflection subgroup proved above, let
\(N_{\rm aff}\subset N_G(S)(k)\) be its inverse image, and put
\(G_{\rm aff}=\langle I_k,N_{\rm aff}\rangle\). Its quotient
\(N_{\rm aff}/(I_k\cap N_{\rm aff})\) is precisely \(W_{\rm aff}\): the
intersection is the compact pointwise kernel \(M(k)^0\), proved directly
in GB.23 below. The reflection group is normal in the full affine
normalizer, since conjugation sends a wall reflection to the reflection
in its image wall. Every full normalizer element can be followed by an
element of \(N_{\rm aff}\) to preserve the base alcove; the remainder
normalizes \(I_k\). Consequently \(G_{\rm aff}\) is normal in the
generated full group and the full group is \(G_{\rm aff}\) times these
residual normalizer elements. Gallery induction itself makes
\(G_{\rm aff}\) transitive on maximal cells. It is also transitive on
apartments containing the base cell: the connected transporter used
above lies in its subgroup \(I_k\). Thus \(X=G_{\rm aff}A_k\), and the
Coxeter arrangement, generation within \(G_{\rm aff}\) and (GB22.3) are
exactly AB.4's hypotheses. Its complete double-coset proof supplies
common maximal rational apartments for any two cells and the exact
finite-face parahorics. This use does not misidentify the full extended
affine normalizer, which can contain central translations and residual
alcove symmetries, with a Coxeter group. Thus the descended construction
has the common maximal apartments and wall exchange required for the
enlarged building.

\subsubsection{\texorpdfstring{GB.23. Full stabilizers, topology, metric
and properness over
\(k\)}{GB.23. Full stabilizers, topology, metric and properness over k}}\label{gb.23.-full-stabilizers-topology-metric-and-properness-over-k}

The group \(G(k)\) is locally compact and second countable: a faithful
affine \(k\)-matrix immersion realizes it as a closed subgroup of
\(\operatorname{GL}_n(k)\), and the original nonarchimedean local field
has compact valuation balls and a countable valuation/residue basis. The
action on \(X\) is continuous. Indeed each point stabilizer is open: the
faithful weighted norm description GB.10 is a finite set of closed
matrix inequalities with discrete coefficient valuations, and a
sufficiently small matrix neighbourhood of the identity satisfies them.
Joint continuity then follows from the isometry inequality
\(d(gx,gy)=d(x,y)\), using a fixed point's open stabilizer.

Every full point stabilizer \(K_x(k)=K_x(G,K)\cap G(k)\) is compact. The
valuation-bound transfer to a faithful \(k\)-matrix immersion is exactly
the polynomial coordinate argument of GB.19: the weighted norm bounds
matrices and inverses in a finite splitting field, hence bounds each
regular \(k\)-matrix coordinate. The stabilizer is closed, so is a
closed subset of a compact product of original \(k\)-valuation balls. No
compactness of \(K=k^{\mathrm{ur}}\) is being asserted.

We first prove the pointwise and setwise stabilizers of the entire
maximal-split flat; neither assertion requires the double-coset
argument. If \(g\in G(k)\) fixes every point of \(A_k\), then
\(sgs^{-1}\) fixes \(x_0\) for every \(s\in S(k)\). Those matrices lie
in the single uniformly bounded full norm stabilizer just proved.
Decompose the faithful representation over a finite splitting field into
its algebraic \(S\)-weight spaces. Its matrix block from weight \(\eta\)
to weight \(\xi\) is multiplied under this conjugation by
\((\xi-\eta)(s)\). If a block of nonzero difference weight were nonzero,
choose a \(k\)-cocharacter \(\lambda\) with nonzero pairing and put
\(s=\lambda(\pi)^j\), letting \(j\) tend to either sign of infinity.
That block would violate the uniform valuation bound. Thus all
nonzero-difference blocks vanish. Faithfulness makes \(g\) centralize
\(S\), so \(g\in M(k)\). Its translation on the centralizer flat is
zero, hence \(g\in M(k)^0=\ker\nu_M\). Conversely GB.19 shows that this
kernel fixes the whole flat.

The subgroup \(M(k)^0\) is Zariski dense in \(M\). It contains an
ordinary open neighbourhood of the identity in \(M(k)\), since the
finitely many character valuations in (GB19.3) vanish near one. A
\(k\)-analytic open neighbourhood of a smooth rational point is Zariski
dense in the smooth connected group: take an étale coordinate chart
there; the finite polynomial Jacobian equations, after scaling to a
sufficiently small valuation ball, have unit Jacobian and their Newton
corrections converge, giving a branch over a full \(k\)-ball. A nonzero
regular function vanishing on that branch would have a nonzero field
norm to the coordinate function field (the chart is generically finite
and separable) vanishing on that ball wherever its denominators are
defined. A nonzero polynomial cannot vanish on a full product of
infinite \(k\)-balls, by induction on the number of variables and the
one-variable root bound. This is a contradiction. One may restrict the
chart and ball to remove all those denominators; the coordinate norm
remains nonzero. This proves the density also in imperfect equal
characteristic.

An element carrying \(A_k\) to itself conjugates its pointwise kernel to
itself. By this density it therefore normalizes the algebraic group
\(M\). The torus \(S\) is the unique maximal \(k\)-split subtorus of
\(Z(M)^0_{\rm red}\): it is central and split, and any additional
central split direction would enlarge it in \(G\). Equivalently this
unique torus is defined by the invariant cocharacter sublattice of the
centre, so every \(k\)-automorphism of \(M\) preserves it. Hence the
element normalizes \(S\). Conversely \(N_G(S)(k)\) preserves the
centralizer building and its unique fixed flat by GB.19. We have proved
that the setwise stabilizer of \(A_k\) is exactly \(N=N_G(S)(k)\), and
its pointwise kernel is exactly the compact \(M(k)^0\).

Let this \(N\) act on \(A_k\). Its linear image is finite: it permutes
the finite nonzero root restrictions on \(S\), acts trivially on the
rational characters of \(G\), and these together span the apartment's
dual. For each linear label its translation vectors form a coset of the
discrete lattice \(\nu_M(M(k))\), by (GB19.3). Thus the full affine
image is discrete and its point stabilizers are finite; \(N/M(k)^0\)
acts properly on the apartment. Wall reflections are the ones proved in
GB.22, so these facts supply the finite-stabilizer hypothesis of AB.2
used there. Wall-pair translations span the semisimple directions: every
wall direction occurs in an unbounded arithmetic family; parallel
reflection products give nonzero translations in its coroot direction.
The full centre has the lattice from \(S(k)\) in GB.19. Thus the affine
Coxeter alcove has compact closure modulo that entire central factor,
and the extended normalizer has only a finite residual alcove action in
the semisimple directions.

For \(x\in A_k\), the full stabilizer formula is \[
                 K_x(k)=P_F^0(k)\,N_x,\qquad
                 N_x=\{n\in N:nx=x\},                       \tag{GB23.1}
\] where \(F\) is the stable minimal \(K\)-facet of \(x\). If \(g\)
fixes \(x\), its translated maximal-split flat \(gA_k\) also contains
\(x\). The connected-model torus transporter of GB.22 supplies
\(p\in P_F^0(k)\) fixing \(x\) and carrying that flat back to \(A_k\).
Then \(pg\in N_x\). Conversely both factors fix \(x\). The
setwise-normalizer assertion here is precisely the uniform-conjugation
and Zariski-density proof above; thus the adjustment gives an actual
element of \(N_G(S)(k)\).

Connected point groups are exactly those of the actual descended models
GB.10--GB.13. Their full bounded models can have nontrivial finite
component groups. Formula (GB23.1) retains \(N_x\), including its
compact \(M(k)^0\) kernel. If a fixed full-stabilizer coset is used, the
exact component-cocycle test is GB.11: its class must vanish in the
finite component group, after which the connected Lang calculation
supplies the correction. This argument used connected coface and torus
transporters, so it never assumed vanishing of \(H^1(\Gamma,K_x(G,K))\).
In particular it has not equated the full bounded torus with its
connected integral model \(T^0\).

There is an exact quotient description \[
              X\simeq
             (G(k)\times A_k)/\!\sim,\qquad
 (g,x)\sim(h,y)\ \Longleftrightarrow\
 \exists n\in N:\ y=nx,\quad g^{-1}hn\in K_x(k),             \tag{GB23.2}
\] with the normalizer convention chosen consistently with AB.1.
Surjectivity is (GB22.1). For injectivity, if the two represented points
agree, the connected point transporter in GB.22 carries their two
maximal flats to each other while fixing that point; after it is removed
the remaining comparison is a normalizer element, giving the displayed
witness. The reverse implication follows from the point stabilizer
identity. Thus the quotient uses the actual full stabilizers.

All hypotheses of the topological theorem AB.1 now hold over the
original locally compact field. Compactness and closedness of \(K_x\)
were just proved. The local upper-containment \(K_y\subset K_x\) near
\(x\), and closed incidence, follow directly from the finite weighted
norm inequalities: round their finitely many real weight differences to
the discrete coefficient valuation lattices. At a boundary a
neighbouring inequality may become stronger, while it cannot become
weaker than the inequality at that boundary; away from a boundary it is
constant. Matrices and inverses give the same conclusion for norm
equality. Normalizer covariance and properness were proved above.

A chamber's pointwise stabilizer \(I_k\) is compact open. A face has
only finitely many adjacent rational chambers, by the finite set of
residue parabolics GB.21. For a point of a face in \(\overline C_k\),
its compact full point stabilizer contains the open \(I_k\) and
therefore has finite index over \(I_k\): its disjoint open cosets have a
finite subcover by compactness. The same applies to the pointwise face
stabilizer. A setwise chamber stabilizer has only a finite residual
action on its vertices modulo the central factor; its possible central
translations are retained in the normalizer, while a point stabilizer
has zero central translation. Root wall exchange (GB22.3) gives
precisely the finite face groups required by AB.4. Finally the compact
set consisting of the closed semisimple alcove times a central lattice
parallelepiped covers all orbits, by (GB22.1) and the retained lattice
translations. AB.1 therefore proves that (GB23.2) has a locally compact
Hausdorff second-countable topology and a proper continuous
\(G(k)\)-action with compact orbit space.

The quotient metric and the metric inherited from \(\mathcal B(G,K)\)
agree. AB.4 supplies a common maximal rational apartment for any two
cells. Its straight segment is a geodesic in the \(K\)-building, with
the same Euclidean metric as \(A_k\). Conversely every quotient chamber
chain has at least that endpoint distance, since each chamber metric is
the restricted \(K\)-metric and the triangle inequality holds there.
Thus both distances equal the apartment distance. AB.6's retractions and
AB.7's local-star comparison show that this metric gives exactly the
quotient topology; finiteness of the residue panel sets makes bounded
closed balls compact. Completeness and the CAT(0) inequality already
hold on the closed convex fixed set by GB.15, and are thereby also
established for the descended quotient. The central factor is exactly \[
                             V_{Z(G),K}^\Gamma ,
\] with its full Euclidean metric and full lattice of rational central
translations. No reduced-building properness assertion has silently
removed it.

Lastly this enlarged space is a numerable universal proper
\(G(k)\)-space. Compact subgroups have bounded orbits in its proper
metric, and the circumcentre argument GB.15 makes their fixed sets
nonempty; they are convex and contractible. Point stabilizers are
compact. Properness, second countability and the compact quotient give
the invariant finite partition of unity on slice neighbourhoods
constructed in AB.8. Its partition-of-unity and
squared-distance-minimization proof then gives the universal property
for numerable proper spaces. Theorem 3.2 also gives the mapping property
directly for every locally compact proper source. This proves the
arbitrary connected reductive local-field building and properness scope,
in every residue characteristic, with the exact finite-unramified
arithmetic, connected models, component correction and metric
comparisons bound above.

The argument used generating root rays and their exact integral
filtration, opposite completion and wall exchange. A separate
classification of the original nonsplit \(k\)-rank-one groups as
classical matrix groups is not needed or claimed: their residue rank-one
group is quasi-split over the finite field and has the proved orbit
classification GB.1. Internal \(2a\) unitary parameters and every actual
restricted character weight have remained in (GB20.1)--(GB22.2).
Consequently no unproved standalone list of nonsplit classical forms is
a prerequisite of this construction.

\subsection{What this lesson does not
prove}\label{what-this-lesson-does-not-prove}

The equivariant Kasparov-cycle, product and homotopy comparison results
used in Proposition 4.4 are proved in \href{KT-KK-17.html}{Equivariant
KK-theory and the Green--Julg theorem}, Theorems 2.3, 4.3 and
Proposition 6.1. The topological and Haar prerequisites for Theorems 1.3
and 2.1 are proved in
\href{supporting/noncompact-foundations/noncompact-foundations.html}{Noncompact
foundations for equivariant induction}, NCF.1 and NCF.3--NCF.5. The
exact earlier algebraic-group and local-field cohomology proofs used by
the building construction are named in BF.0, AB.9 and GB.14--GB.18. They
are prerequisites, not replacements by external literature citations.

The numerable recognition theorem, the almost-connected
homogeneous-space theorem, the Rips-complex theorem and the
nonequivariant geometric/analytic comparison are stated and unused in
Section 5. Their precise statements and freely accessible sources are
given there. The assembly homomorphism, its conjectural isomorphism and
its consequences belong to the following lesson.

\subsection{References}\label{references}

\begin{itemize}
\tightlist
\item
  Wolfgang Lück,
  \href{https://arxiv.org/abs/math/0312378v2}{\emph{Survey on
  Classifying Spaces for Families of Subgroups}}, free author preprint,
  revised 2004. Definitions 1.8 and 2.1, Theorem 2.5, and Theorems 4.3,
  4.5--4.7, 4.13 and 4.16 discuss the mapping category and the models
  named here.
\item
  Maria Paula Gómez Aparicio, Pierre Julg and Alain Valette,
  \href{https://arxiv.org/abs/1905.10081}{\emph{The Baum--Connes
  conjecture: an extended survey}}, 2019, Section 4.2, Definition 4.2,
  and Section 4.3, Remark 4.3. This gives the universal-space and
  equivariant K-homology formulation due to Baum, Connes and Higson,
  including cutoff normalization.
\item
  Paul Baum, Nigel Higson and Thomas Schick,
  \href{https://arxiv.org/abs/math/0701484v4}{\emph{On the Equivalence
  of Geometric and Analytic K-Homology}}, free author preprint, revised
  2009. Definitions 5.1--5.7 and Theorem 6.2 give the finite CW-pair
  comparison stated in Section 5.
\item
  François Bruhat and Jacques Tits,
  \href{https://www.numdam.org/item/PMIHES_1972__41__5_0/}{\emph{Groupes
  réductifs sur un corps local I. Données radicielles valuées}}, 1972.
  Sections 2.3--2.5 and 3.2 provide the apartment, retraction and metric
  constructions; Definition 6.2.1 and Examples 6.2.3 give the
  root-valuation conventions. The proofs used here are in BF, AB and GB.
\item
  François Bruhat and Jacques Tits,
  \href{https://www.numdam.org/item/PMIHES_1984__60__5_0/}{\emph{Groupes
  réductifs sur un corps local II. Schémas en groupes. Existence d'une
  donnée radicielle valuée}}, 1984. Sections 4.1--4.2 and 5.1 treat
  nonreduced roots and unramified descent.
\item
  Michel Demazure,
  \href{https://webusers.imj-prg.fr/~patrick.polo/SGA3/Exp23-13oct24.pdf}{\emph{SGA
  3, Exposé XXIII}}, freely accessible revised text, 2024. Sections
  1.1--1.2 and Corollaries 5.2--5.9 concern integral pinned root
  coordinates.
\item
  Brian Conrad,
  \href{https://math.stanford.edu/~conrad/papers/luminysga3smf.pdf}{\emph{Reductive
  group schemes}}, free author text. Lemma 6.2.8 treats integral root
  constants. His
  \href{https://math.stanford.edu/~conrad/249CS13Page/handouts/langunirat.pdf}{\emph{Lang's
  theorem and unirationality}}, Section 2, gives the finite-field
  context of the local proof in AB.9.
\item
  James S. Milne,
  \href{https://www.jmilne.org/math/CourseNotes/ANT.pdf}{\emph{Algebraic
  Number Theory}}, free author notes, Proposition 2.29 and Theorem 7.38.
  The finite-extension valuation and DVR proof needed here is given in
  AB.0.
\item
  Todd Trimble,
  \href{https://math.ucr.edu/home/baez/trimble/BN-pairs_and_Boolean_Hecke_algebras.html}{\emph{BN-pairs
  and Boolean Hecke algebras}}, 2005. AB.4 proves the exchange
  consequences needed here.
\item
  Bas Edixhoven and Matthieu Romagny,
  \href{https://arxiv.org/abs/1204.1799v3}{\emph{Group schemes out of
  birational group laws, Néron models}}, revised 2013. The integral
  model arguments are supplied in GB.7--GB.13.
\item
  Armand Borel and Tonny A. Springer,
  \href{https://www.jstage.jst.go.jp/article/tmj1949/20/4/20_4_443/_article}{\emph{Rationality
  properties of linear algebraic groups II}}, 1968, Section 8.6. The
  arithmetic-dimension-one argument, including imperfect fields, is
  proved in GB.14 and GB.16--GB.18.
\item
  Robert Steinberg,
  \href{https://www.numdam.org/item/PMIHES_1965__25__49_0/}{\emph{Regular
  elements of semi-simple algebraic groups}}, 1965. The regular-class
  and cocycle arguments needed for quasi-splitting are proved in
  GB.17--GB.18.
\item
  Romyar Sharifi,
  \href{https://www.math.ucla.edu/~sharifi/groupcoh.pdf}{\emph{Group and
  Galois Cohomology}}, free author notes, Section 1.11, Lemmas
  1.11.4--1.11.9 and Theorem 1.11.11. GB.16 proves the integral lattice
  criterion and torus vanishing used here.
\end{itemize}

\end{document}
